% Options for packages loaded elsewhere
\PassOptionsToPackage{unicode}{hyperref}
\PassOptionsToPackage{hyphens}{url}
\PassOptionsToPackage{dvipsnames,svgnames*,x11names*}{xcolor}
%
\documentclass[
  11pt,
  a4paper,
]{report}
\usepackage{lmodern}
\usepackage{amssymb,amsmath}
\usepackage{ifxetex,ifluatex}
\ifnum 0\ifxetex 1\fi\ifluatex 1\fi=0 % if pdftex
  \usepackage[T1]{fontenc}
  \usepackage[utf8]{inputenc}
  \usepackage{textcomp} % provide euro and other symbols
\else % if luatex or xetex
  \usepackage{unicode-math}
  \defaultfontfeatures{Scale=MatchLowercase}
  \defaultfontfeatures[\rmfamily]{Ligatures=TeX,Scale=1}
\fi
% Use upquote if available, for straight quotes in verbatim environments
\IfFileExists{upquote.sty}{\usepackage{upquote}}{}
\IfFileExists{microtype.sty}{% use microtype if available
  \usepackage[]{microtype}
  \UseMicrotypeSet[protrusion]{basicmath} % disable protrusion for tt fonts
}{}
\makeatletter
\@ifundefined{KOMAClassName}{% if non-KOMA class
  \IfFileExists{parskip.sty}{%
    \usepackage{parskip}
  }{% else
    \setlength{\parindent}{0pt}
    \setlength{\parskip}{6pt plus 2pt minus 1pt}}
}{% if KOMA class
  \KOMAoptions{parskip=half}}
\makeatother
\usepackage{xcolor}
\IfFileExists{xurl.sty}{\usepackage{xurl}}{} % add URL line breaks if available
\IfFileExists{bookmark.sty}{\usepackage{bookmark}}{\usepackage{hyperref}}
\hypersetup{
  colorlinks=true,
  linkcolor=blue,
  filecolor=Maroon,
  citecolor=Blue,
  urlcolor=blue,
  pdfcreator={LaTeX via pandoc}}
\urlstyle{same} % disable monospaced font for URLs
\usepackage[margin=1in]{geometry}
\usepackage{longtable,booktabs}
% Correct order of tables after \paragraph or \subparagraph
\usepackage{etoolbox}
\makeatletter
\patchcmd\longtable{\par}{\if@noskipsec\mbox{}\fi\par}{}{}
\makeatother
% Allow footnotes in longtable head/foot
\IfFileExists{footnotehyper.sty}{\usepackage{footnotehyper}}{\usepackage{footnote}}
\makesavenoteenv{longtable}
\setlength{\emergencystretch}{3em} % prevent overfull lines
\providecommand{\tightlist}{%
  \setlength{\itemsep}{0pt}\setlength{\parskip}{0pt}}
\setcounter{secnumdepth}{3}

% Headers and footers
\usepackage{fancyhdr}
\pagestyle{fancy}
\fancyhf{}
\fancyhead[L]{\small\leftmark}
\fancyhead[R]{\small CMC in Drug Development}
\fancyfoot[C]{\thepage}
\renewcommand{\headrulewidth}{0.4pt}
\renewcommand{\footrulewidth}{0.2pt}

% Better section formatting
\usepackage{titlesec}
\titleformat{\chapter}[display]{\normalfont\huge\bfseries}{\chaptertitlename\ \thechapter}{20pt}{\Huge}
\titlespacing*{\chapter}{0pt}{-20pt}{30pt}

\author{}
\date{}

\begin{document}

% ============ TITLE PAGE ============
\begin{titlepage}
\centering
\vspace*{2cm}

{\Huge\bfseries CMC in Drug Development\par}
\vspace{0.5cm}
{\LARGE\bfseries The Definitive Guide\par}
\vspace{1cm}
{\Large Chemistry, Manufacturing, and Controls\par}
{\Large From Discovery to Post-Market Lifecycle\par}

\vspace{2cm}

\rule{\textwidth}{0.4pt}
\vspace{0.5cm}

{\large\itshape A single, self-contained expert-level reference covering every dimension of CMC in pharmaceutical drug development. Maps CMC activities to clinical trial phases, details regulatory expectations across FDA and EMA, and provides the practical depth needed for strategic client engagements.\par}

\vspace{0.5cm}
\rule{\textwidth}{0.4pt}

\vspace{2cm}

{\large\bfseries Coverage\par}
\vspace{0.5cm}
{\normalsize
34 Sections $\bullet$ 80+ Glossary Terms $\bullet$ 2,000+ Lines\\[0.3cm]
ICH Q1--Q14 $\bullet$ FDA 21 CFR $\bullet$ EMA Guidelines\\[0.3cm]
Small Molecules $\bullet$ Biologics $\bullet$ Cell \& Gene Therapy $\bullet$ mRNA\\[0.3cm]
ADCs $\bullet$ Bispecifics $\bullet$ Oligonucleotides $\bullet$ Biosimilars $\bullet$ Generics
\par}

\vfill

{\large Compiled: March 2026\par}
\vspace{0.5cm}
{\small Sources: ICH Guidelines, FDA Regulations \& Guidances, EMA Guidelines, USP/Ph. Eur. Standards\par}

\end{titlepage}

% ============ TABLE OF CONTENTS ============
\setcounter{tocdepth}{2}
{
\hypersetup{linkcolor=blue}
\tableofcontents
}
\hypertarget{cmc-in-drug-development-the-definitive-guide}{%
\chapter{CMC in Drug Development: The Definitive
Guide}\label{cmc-in-drug-development-the-definitive-guide}}

\hypertarget{chemistry-manufacturing-and-controls-from-discovery-to-post-market-lifecycle}{%
\section{Chemistry, Manufacturing, and Controls --- From Discovery to
Post-Market
Lifecycle}\label{chemistry-manufacturing-and-controls-from-discovery-to-post-market-lifecycle}}

\begin{center}\rule{0.5\linewidth}{0.5pt}\end{center}

\begin{quote}
\textbf{Purpose:} This document is a single, self-contained expert-level
reference covering every dimension of CMC in pharmaceutical drug
development. It maps CMC activities to clinical trial phases, details
regulatory expectations across FDA and EMA, and provides the practical
depth needed for strategic client engagements.
\end{quote}

\begin{center}\rule{0.5\linewidth}{0.5pt}\end{center}

\hypertarget{table-of-contents}{%
\section{Table of Contents}\label{table-of-contents}}

\begin{enumerate}
\def\labelenumi{\arabic{enumi}.}
\tightlist
\item
  \protect\hyperlink{1-what-is-cmc-and-why-it-matters}{What is CMC and
  Why It Matters}
\item
  \protect\hyperlink{2-the-four-pillars-of-cmc}{The Four Pillars of CMC}
\item
  \protect\hyperlink{3-regulatory-framework--ich-guidelines}{Regulatory
  Framework --- ICH Guidelines}
\item
  \protect\hyperlink{4-the-common-technical-document-ctd-module-3}{The
  Common Technical Document (CTD) Module 3}
\item
  \protect\hyperlink{5-drug-substance-api--deep-dive}{Drug Substance
  (API) --- Deep Dive}
\item
  \protect\hyperlink{6-drug-product--deep-dive}{Drug Product --- Deep
  Dive}
\item
  \protect\hyperlink{7-analytical-methods--deep-dive}{Analytical Methods
  --- Deep Dive}
\item
  \protect\hyperlink{8-container-closure-systems}{Container Closure
  Systems}
\item
  \protect\hyperlink{9-stability-studies}{Stability Studies}
\item
  \protect\hyperlink{10-process-validation}{Process Validation}
\item
  \protect\hyperlink{11-clinical-trial-phases-as-cmc-triggers}{Clinical
  Trial Phases as CMC Triggers}
\item
  \protect\hyperlink{12-cmc-strategy-and-planning}{CMC Strategy and
  Planning}
\item
  \protect\hyperlink{13-regulatory-interactions-for-cmc}{Regulatory
  Interactions for CMC}
\item
  \protect\hyperlink{14-post-approval-cmc-changes-and-lifecycle-management}{Post-Approval
  CMC Changes and Lifecycle Management}
\item
  \protect\hyperlink{15-biologics-specific-cmc}{Biologics-Specific CMC}
\item
  \protect\hyperlink{16-small-molecule-vs-biologics-cmc--key-differences}{Small
  Molecule vs.~Biologics CMC --- Key Differences}
\item
  \protect\hyperlink{17-advanced-therapies--cell-gene-and-mrna-cmc}{Advanced
  Therapies --- Cell, Gene, and mRNA CMC}
\item
  \protect\hyperlink{18-emerging-trends-in-cmc}{Emerging Trends in CMC}
\item
  \protect\hyperlink{19-common-cmc-pitfalls-and-how-to-avoid-them}{Common
  CMC Pitfalls and How to Avoid Them}
\item
  \protect\hyperlink{20-fda-vs-ema--cmc-comparison}{FDA vs.~EMA --- CMC
  Comparison}
\item
  \protect\hyperlink{21-technology-transfer}{Technology Transfer}
\item
  \protect\hyperlink{22-cmocdmo-management-and-quality-agreements}{CMO/CDMO
  Management and Quality Agreements}
\item
  \protect\hyperlink{23-supply-chain-and-vendor-qualification}{Supply
  Chain and Vendor Qualification}
\item
  \protect\hyperlink{24-biosimilar-cmc-pathway}{Biosimilar CMC Pathway}
\item
  \protect\hyperlink{25-combination-products-and-fixed-dose-combinations}{Combination
  Products and Fixed-Dose Combinations}
\item
  \protect\hyperlink{26-nitrosamine-impurity-assessment}{Nitrosamine
  Impurity Assessment}
\item
  \protect\hyperlink{27-emerging-modalities--adcs-bispecifics-oligonucleotides}{Emerging
  Modalities --- ADCs, Bispecifics, Oligonucleotides}
\item
  \protect\hyperlink{28-drug-master-files-dmfs}{Drug Master Files
  (DMFs)}
\item
  \protect\hyperlink{29-data-integrity-and-electronic-records}{Data
  Integrity and Electronic Records}
\item
  \protect\hyperlink{30-facility-design-environmental-controls-and-utilities}{Facility
  Design, Environmental Controls, and Utilities}
\item
  \protect\hyperlink{31-cmc-due-diligence-for-licensing-and-ma}{CMC Due
  Diligence for Licensing and M\&A}
\item
  \protect\hyperlink{32-generic-drug-anda-cmc-requirements}{Generic Drug
  (ANDA) CMC Requirements}
\item
  \protect\hyperlink{33-master-reference-table}{Master Reference Table
  --- ICH Guidelines, FDA Guidances, and Regulations}
\item
  \protect\hyperlink{34-glossary-of-key-cmc-terms}{Glossary of Key CMC
  Terms}
\end{enumerate}

\begin{center}\rule{0.5\linewidth}{0.5pt}\end{center}

\hypertarget{what-is-cmc-and-why-it-matters}{%
\section{1. What is CMC and Why It
Matters}\label{what-is-cmc-and-why-it-matters}}

\textbf{CMC = Chemistry, Manufacturing, and Controls.} It encompasses
all technical information about a drug's composition, how it is made,
and how its quality is assured. CMC is the quality backbone of every
regulatory submission --- without robust CMC data, no drug can receive
regulatory approval regardless of how strong the clinical efficacy data
may be.

\hypertarget{the-core-proposition}{%
\subsection{The Core Proposition}\label{the-core-proposition}}

CMC is about \textbf{proving your drug is reliably made using a
reproducible, scalable process and is safe for patients.} Every tablet,
vial, or syringe that reaches a patient must be: - \textbf{Identified}
--- it is what the label says it is - \textbf{Pure} --- free of harmful
impurities above qualified thresholds - \textbf{Potent} --- it delivers
the intended therapeutic effect at the stated strength - \textbf{Stable}
--- it maintains its quality throughout its shelf life -
\textbf{Consistently manufactured} --- every batch is equivalent to the
batches that proved safe and effective in clinical trials

\hypertarget{why-cmc-matters-strategically}{%
\subsection{Why CMC Matters
Strategically}\label{why-cmc-matters-strategically}}

\begin{itemize}
\tightlist
\item
  \textbf{CMC deficiencies are a leading cause of FDA Complete Response
  Letters (CRLs)}, delaying approvals by months or years.
\item
  \textbf{CMC timelines often become the critical path} --- process
  validation, stability data generation, and manufacturing site
  readiness cannot be accelerated beyond physical and regulatory
  constraints.
\item
  \textbf{CMC investment is substantial} --- 15-25\% of total
  development costs for small molecules, 30-50\% for biologics, and
  40-60\% for cell/gene therapies.
\item
  \textbf{Post-approval CMC changes} (site transfers, process
  improvements, new formulations) require ongoing regulatory management
  for the life of the product.
\end{itemize}

\begin{center}\rule{0.5\linewidth}{0.5pt}\end{center}

\hypertarget{the-four-pillars-of-cmc}{%
\section{2. The Four Pillars of CMC}\label{the-four-pillars-of-cmc}}

\hypertarget{pillar-1-drug-substance-active-pharmaceutical-ingredient-api}{%
\subsection{Pillar 1: Drug Substance (Active Pharmaceutical Ingredient /
API)}\label{pillar-1-drug-substance-active-pharmaceutical-ingredient-api}}

Everything about the active ingredient before it is formulated: -
\textbf{Structure and characterization}: Molecular formula,
stereochemistry, solid-state properties - \textbf{Manufacturing
process}: Synthetic route (small molecules) or cell culture/fermentation
(biologics) - \textbf{Impurity profile}: Process-related impurities,
degradation products, elemental impurities, residual solvents, mutagenic
impurities - \textbf{Specifications}: Tests and acceptance criteria for
identity, assay, purity, physical properties - \textbf{Stability}:
Retest period supported by ICH-compliant stability data

\hypertarget{pillar-2-drug-product-finished-dosage-form}{%
\subsection{Pillar 2: Drug Product (Finished Dosage
Form)}\label{pillar-2-drug-product-finished-dosage-form}}

The formulated product as administered to patients: -
\textbf{Formulation}: Quantitative composition, excipient selection and
justification - \textbf{Manufacturing process}: Unit operations, batch
formula, in-process controls - \textbf{Container closure system}:
Primary packaging, extractables/leachables - \textbf{Specifications}:
Release and stability testing (identity, assay, content uniformity,
dissolution, degradation products, microbial limits) -
\textbf{Stability}: Shelf life supported by ICH-compliant stability data

\hypertarget{pillar-3-analytical-methods}{%
\subsection{Pillar 3: Analytical
Methods}\label{pillar-3-analytical-methods}}

Tools used to measure and ensure quality: - \textbf{Identity tests}: IR,
UV, HPLC retention time, mass spectrometry - \textbf{Assay/Potency}:
HPLC, UV spectrophotometry, bioassays (biologics) - \textbf{Impurity
testing}: HPLC, GC-HS (residual solvents), ICP-MS (elemental impurities)
- \textbf{Physical tests}: Dissolution, content uniformity, particle
size, hardness - \textbf{Microbiological tests}: Sterility, endotoxins,
microbial enumeration - \textbf{Validation}: Per ICH Q2(R2) ---
specificity, linearity, accuracy, precision, range, LOD/LOQ, robustness

\hypertarget{pillar-4-manufacturing-controls}{%
\subsection{Pillar 4: Manufacturing
Controls}\label{pillar-4-manufacturing-controls}}

Systems ensuring reproducible, consistent quality: - \textbf{cGMP
compliance}: 21 CFR 210/211 (US), EU GMP Annexes, ICH Q7 (APIs) -
\textbf{Process validation}: FDA 2011 lifecycle approach (Stages 1-3) -
\textbf{In-process controls}: Tests and measurements during
manufacturing - \textbf{Environmental controls}: Cleanroom
classifications, environmental monitoring - \textbf{Equipment
qualification}: IQ/OQ/PQ - \textbf{Cleaning validation}: Product and
cleaning agent residue removal - \textbf{Supply chain controls}: Vendor
qualification, incoming material testing

\begin{center}\rule{0.5\linewidth}{0.5pt}\end{center}

\hypertarget{regulatory-framework-ich-guidelines}{%
\section{3. Regulatory Framework --- ICH
Guidelines}\label{regulatory-framework-ich-guidelines}}

The International Council for Harmonisation (ICH) develops guidelines
accepted by FDA, EMA, PMDA (Japan), Health Canada, Swissmedic, and
others. The Quality (``Q'') guidelines are the backbone of CMC
regulatory science.

\hypertarget{ich-q1-stability-testing}{%
\subsection{ICH Q1: Stability Testing}\label{ich-q1-stability-testing}}

\begin{longtable}[]{@{}lll@{}}
\toprule
\begin{minipage}[b]{0.32\columnwidth}\raggedright
Guideline\strut
\end{minipage} & \begin{minipage}[b]{0.21\columnwidth}\raggedright
Title\strut
\end{minipage} & \begin{minipage}[b]{0.38\columnwidth}\raggedright
Key Content\strut
\end{minipage}\tabularnewline
\midrule
\endhead
\begin{minipage}[t]{0.32\columnwidth}\raggedright
\textbf{Q1A(R2)}\strut
\end{minipage} & \begin{minipage}[t]{0.21\columnwidth}\raggedright
Stability Testing of New Drug Substances and Products\strut
\end{minipage} & \begin{minipage}[t]{0.38\columnwidth}\raggedright
Core protocol --- storage conditions (25°C/60\% RH long-term, 40°C/75\%
RH accelerated), testing frequency, minimum 12 months long-term + 6
months accelerated at filing\strut
\end{minipage}\tabularnewline
\begin{minipage}[t]{0.32\columnwidth}\raggedright
\textbf{Q1B}\strut
\end{minipage} & \begin{minipage}[t]{0.21\columnwidth}\raggedright
Photostability Testing\strut
\end{minipage} & \begin{minipage}[t]{0.38\columnwidth}\raggedright
Minimum 1.2 million lux-hours visible + 200 Wh/m² UV exposure\strut
\end{minipage}\tabularnewline
\begin{minipage}[t]{0.32\columnwidth}\raggedright
\textbf{Q1C}\strut
\end{minipage} & \begin{minipage}[t]{0.21\columnwidth}\raggedright
Stability Testing for New Dosage Forms\strut
\end{minipage} & \begin{minipage}[t]{0.38\columnwidth}\raggedright
Reduced data acceptable for new dosage forms of approved
substances\strut
\end{minipage}\tabularnewline
\begin{minipage}[t]{0.32\columnwidth}\raggedright
\textbf{Q1D}\strut
\end{minipage} & \begin{minipage}[t]{0.21\columnwidth}\raggedright
Bracketing and Matrixing\strut
\end{minipage} & \begin{minipage}[t]{0.38\columnwidth}\raggedright
Statistical designs to reduce stability testing burden\strut
\end{minipage}\tabularnewline
\begin{minipage}[t]{0.32\columnwidth}\raggedright
\textbf{Q1E}\strut
\end{minipage} & \begin{minipage}[t]{0.21\columnwidth}\raggedright
Evaluation of Stability Data\strut
\end{minipage} & \begin{minipage}[t]{0.38\columnwidth}\raggedright
Statistical analysis, shelf-life extrapolation rules (up to 2× long-term
data, max 12 months beyond)\strut
\end{minipage}\tabularnewline
\bottomrule
\end{longtable}

\textbf{Stability Storage Conditions:}

\begin{longtable}[]{@{}llll@{}}
\toprule
\begin{minipage}[b]{0.17\columnwidth}\raggedright
Study Type\strut
\end{minipage} & \begin{minipage}[b]{0.17\columnwidth}\raggedright
Conditions\strut
\end{minipage} & \begin{minipage}[b]{0.26\columnwidth}\raggedright
Minimum Duration\strut
\end{minipage} & \begin{minipage}[b]{0.29\columnwidth}\raggedright
Testing Frequency\strut
\end{minipage}\tabularnewline
\midrule
\endhead
\begin{minipage}[t]{0.17\columnwidth}\raggedright
Long-term\strut
\end{minipage} & \begin{minipage}[t]{0.17\columnwidth}\raggedright
25°C ± 2°C / 60\% RH ± 5\%\strut
\end{minipage} & \begin{minipage}[t]{0.26\columnwidth}\raggedright
12 months at filing\strut
\end{minipage} & \begin{minipage}[t]{0.29\columnwidth}\raggedright
0, 3, 6, 9, 12, 18, 24, 36 months\strut
\end{minipage}\tabularnewline
\begin{minipage}[t]{0.17\columnwidth}\raggedright
Intermediate\strut
\end{minipage} & \begin{minipage}[t]{0.17\columnwidth}\raggedright
30°C ± 2°C / 65\% RH ± 5\%\strut
\end{minipage} & \begin{minipage}[t]{0.26\columnwidth}\raggedright
12 months\strut
\end{minipage} & \begin{minipage}[t]{0.29\columnwidth}\raggedright
0, 6, 9, 12 months\strut
\end{minipage}\tabularnewline
\begin{minipage}[t]{0.17\columnwidth}\raggedright
Accelerated\strut
\end{minipage} & \begin{minipage}[t]{0.17\columnwidth}\raggedright
40°C ± 2°C / 75\% RH ± 5\%\strut
\end{minipage} & \begin{minipage}[t]{0.26\columnwidth}\raggedright
6 months\strut
\end{minipage} & \begin{minipage}[t]{0.29\columnwidth}\raggedright
0, 3, 6 months\strut
\end{minipage}\tabularnewline
\bottomrule
\end{longtable}

\textbf{``Significant Change'' Triggers (Drug Product):} ≥5\% assay
change, degradant exceeding limit, pH outside range, dissolution failure
for 12 units, appearance/physical property failure.

\hypertarget{ich-q2r2-validation-of-analytical-procedures}{%
\subsection{ICH Q2(R2): Validation of Analytical
Procedures}\label{ich-q2r2-validation-of-analytical-procedures}}

All analytical methods used for release and stability must be validated:

\begin{longtable}[]{@{}ll@{}}
\toprule
\begin{minipage}[b]{0.43\columnwidth}\raggedright
Parameter\strut
\end{minipage} & \begin{minipage}[b]{0.51\columnwidth}\raggedright
Requirement\strut
\end{minipage}\tabularnewline
\midrule
\endhead
\begin{minipage}[t]{0.43\columnwidth}\raggedright
\textbf{Specificity}\strut
\end{minipage} & \begin{minipage}[t]{0.51\columnwidth}\raggedright
Resolve analyte from all degradants, impurities, and matrix\strut
\end{minipage}\tabularnewline
\begin{minipage}[t]{0.43\columnwidth}\raggedright
\textbf{Linearity}\strut
\end{minipage} & \begin{minipage}[t]{0.51\columnwidth}\raggedright
Minimum 5 concentration levels; demonstrate proportional response\strut
\end{minipage}\tabularnewline
\begin{minipage}[t]{0.43\columnwidth}\raggedright
\textbf{Range}\strut
\end{minipage} & \begin{minipage}[t]{0.51\columnwidth}\raggedright
80-120\% (assay), 70-130\% (content uniformity), ±20\% beyond spec
(dissolution)\strut
\end{minipage}\tabularnewline
\begin{minipage}[t]{0.43\columnwidth}\raggedright
\textbf{Accuracy}\strut
\end{minipage} & \begin{minipage}[t]{0.51\columnwidth}\raggedright
Minimum 9 determinations across 3 concentrations (3 replicates
each)\strut
\end{minipage}\tabularnewline
\begin{minipage}[t]{0.43\columnwidth}\raggedright
\textbf{Precision}\strut
\end{minipage} & \begin{minipage}[t]{0.51\columnwidth}\raggedright
Repeatability (min 9 determinations), intermediate precision (different
days/analysts/equipment), reproducibility (between labs)\strut
\end{minipage}\tabularnewline
\begin{minipage}[t]{0.43\columnwidth}\raggedright
\textbf{Detection Limit}\strut
\end{minipage} & \begin{minipage}[t]{0.51\columnwidth}\raggedright
S/N ≥ 3:1 or standard deviation of response/slope\strut
\end{minipage}\tabularnewline
\begin{minipage}[t]{0.43\columnwidth}\raggedright
\textbf{Quantitation Limit}\strut
\end{minipage} & \begin{minipage}[t]{0.51\columnwidth}\raggedright
S/N ≥ 10:1 or standard deviation/slope\strut
\end{minipage}\tabularnewline
\begin{minipage}[t]{0.43\columnwidth}\raggedright
\textbf{Robustness}\strut
\end{minipage} & \begin{minipage}[t]{0.51\columnwidth}\raggedright
Deliberate parameter variation (pH ±0.2, temperature ±5°C, flow rate
±10\%)\strut
\end{minipage}\tabularnewline
\bottomrule
\end{longtable}

The R2 revision added Analytical Target Profile (ATP), enhanced
statistical treatment, and lifecycle management concepts.

\hypertarget{ich-q3-impurities}{%
\subsection{ICH Q3: Impurities}\label{ich-q3-impurities}}

\textbf{Q3A(R2) --- Drug Substance Impurity Thresholds:}

\begin{longtable}[]{@{}llll@{}}
\toprule
Max Daily Dose & Reporting & Identification &
Qualification\tabularnewline
\midrule
\endhead
≤2 g/day & ≥0.05\% & ≥0.10\% or 1.0 mg/day & ≥0.15\% or 1.0
mg/day\tabularnewline
\textgreater2 g/day & ≥0.03\% & ≥0.05\% & ≥0.05\%\tabularnewline
\bottomrule
\end{longtable}

\textbf{Q3B(R2) --- Drug Product Degradation Thresholds:}

\begin{longtable}[]{@{}llll@{}}
\toprule
\begin{minipage}[b]{0.24\columnwidth}\raggedright
Max Daily Dose\strut
\end{minipage} & \begin{minipage}[b]{0.17\columnwidth}\raggedright
Reporting\strut
\end{minipage} & \begin{minipage}[b]{0.24\columnwidth}\raggedright
Identification\strut
\end{minipage} & \begin{minipage}[b]{0.23\columnwidth}\raggedright
Qualification\strut
\end{minipage}\tabularnewline
\midrule
\endhead
\begin{minipage}[t]{0.24\columnwidth}\raggedright
≤1 g/day\strut
\end{minipage} & \begin{minipage}[t]{0.17\columnwidth}\raggedright
≥0.1\%\strut
\end{minipage} & \begin{minipage}[t]{0.24\columnwidth}\raggedright
0.2\% (low dose) to 0.1\% (high dose)\strut
\end{minipage} & \begin{minipage}[t]{0.23\columnwidth}\raggedright
1.0\% (low dose) to 0.15\% (high dose)\strut
\end{minipage}\tabularnewline
\begin{minipage}[t]{0.24\columnwidth}\raggedright
\textgreater1 g/day\strut
\end{minipage} & \begin{minipage}[t]{0.17\columnwidth}\raggedright
≥0.05\%\strut
\end{minipage} & \begin{minipage}[t]{0.24\columnwidth}\raggedright
≥0.05\%\strut
\end{minipage} & \begin{minipage}[t]{0.23\columnwidth}\raggedright
≥0.05\%\strut
\end{minipage}\tabularnewline
\bottomrule
\end{longtable}

\textbf{Q3C(R8) --- Residual Solvents:} - \textbf{Class 1} (avoid):
Benzene (2 ppm), carbon tetrachloride (4 ppm), 1,2-dichloroethane (5
ppm) - \textbf{Class 2} (limit): Methylene chloride (600 ppm),
chloroform (60 ppm), acetonitrile (410 ppm), DMF (880 ppm), toluene (890
ppm) - \textbf{Class 3} (low toxicity): Acetone, ethanol, ethyl acetate,
isopropanol --- limit 5000 ppm

\textbf{Q3D(R2) --- Elemental Impurities:} - \textbf{Class 1} (highest
toxicity): As, Cd, Hg, Pb --- controlled in all drug products -
\textbf{Class 2A}: Co, Ni, V --- evaluated across all routes -
\textbf{Class 2B}: Ag, Au, Ir, Os, Pd, Pt, Rh, Ru, Se, Tl --- evaluated
if intentionally added - \textbf{Class 3}: Ba, Cr, Cu, Li, Mo, Sb, Sn, W
--- significant only for inhalation

\textbf{Q3E --- Extractables and Leachables:} Risk-based approach using
Safety Concern Threshold (SCT) and Qualification Threshold (QT).

\textbf{ICH M7(R2) --- Mutagenic Impurities:} - Threshold of
Toxicological Concern (TTC): 1.5 µg/day for lifetime exposure -
Less-than-lifetime (LTL) adjustments allow higher limits for shorter
clinical durations - In silico assessment (Derek Nexus, Sarah Nexus) for
all synthetic intermediates and plausible by-products

\hypertarget{ich-q5-quality-of-biotechnological-products}{%
\subsection{ICH Q5: Quality of Biotechnological
Products}\label{ich-q5-quality-of-biotechnological-products}}

\begin{longtable}[]{@{}lll@{}}
\toprule
\begin{minipage}[b]{0.32\columnwidth}\raggedright
Guideline\strut
\end{minipage} & \begin{minipage}[b]{0.21\columnwidth}\raggedright
Title\strut
\end{minipage} & \begin{minipage}[b]{0.38\columnwidth}\raggedright
Key Content\strut
\end{minipage}\tabularnewline
\midrule
\endhead
\begin{minipage}[t]{0.32\columnwidth}\raggedright
\textbf{Q5A(R2)}\strut
\end{minipage} & \begin{minipage}[t]{0.21\columnwidth}\raggedright
Viral Safety Evaluation\strut
\end{minipage} & \begin{minipage}[t]{0.38\columnwidth}\raggedright
Three-pronged approach: cell line testing, viral clearance validation
(≥4 log₁₀ reduction), product testing\strut
\end{minipage}\tabularnewline
\begin{minipage}[t]{0.32\columnwidth}\raggedright
\textbf{Q5B}\strut
\end{minipage} & \begin{minipage}[t]{0.21\columnwidth}\raggedright
Analysis of Expression Construct\strut
\end{minipage} & \begin{minipage}[t]{0.38\columnwidth}\raggedright
Characterize vector, host cell, genetic stability through production
culture\strut
\end{minipage}\tabularnewline
\begin{minipage}[t]{0.32\columnwidth}\raggedright
\textbf{Q5C}\strut
\end{minipage} & \begin{minipage}[t]{0.21\columnwidth}\raggedright
Stability of Biotech Products\strut
\end{minipage} & \begin{minipage}[t]{0.38\columnwidth}\raggedright
Potency is primary stability indicator; accelerated data are supportive
only\strut
\end{minipage}\tabularnewline
\begin{minipage}[t]{0.32\columnwidth}\raggedright
\textbf{Q5D}\strut
\end{minipage} & \begin{minipage}[t]{0.21\columnwidth}\raggedright
Cell Substrate Characterisation\strut
\end{minipage} & \begin{minipage}[t]{0.38\columnwidth}\raggedright
MCB/WCB testing: identity, purity, stability, adventitious agents\strut
\end{minipage}\tabularnewline
\begin{minipage}[t]{0.32\columnwidth}\raggedright
\textbf{Q5E}\strut
\end{minipage} & \begin{minipage}[t]{0.21\columnwidth}\raggedright
Comparability\strut
\end{minipage} & \begin{minipage}[t]{0.38\columnwidth}\raggedright
Tiered approach: analytical → functional → nonclinical → clinical (if
needed)\strut
\end{minipage}\tabularnewline
\bottomrule
\end{longtable}

\hypertarget{ich-q6-specifications}{%
\subsection{ICH Q6: Specifications}\label{ich-q6-specifications}}

\textbf{Q6A (Chemical Substances):} - Universal drug substance tests:
appearance, identity, assay, impurities - Universal drug product tests:
appearance, identity, assay, degradation products - Specific tests as
appropriate: dissolution, content uniformity, water content, residual
solvents, microbial limits, pH, sterility/endotoxins (parenterals) -
Decision trees for polymorphism, dissolution, microbial testing

\textbf{Q6B (Biotechnological Products):} - Identity: immunochemical,
peptide mapping, electrophoretic - Purity: SEC, IEX, CE-SDS, rp-HPLC -
Potency: biological activity assay (mechanism-based) - Process-related
impurities: HCP, host cell DNA, Protein A, media components -
Product-related impurities: aggregates, fragments, charge variants,
deamidated/oxidized variants

\hypertarget{ich-q7-gmp-for-apis}{%
\subsection{ICH Q7: GMP for APIs}\label{ich-q7-gmp-for-apis}}

The definitive GMP guidance for API manufacturing: - Applies from the
``regulatory starting material'' onward - Covers: quality management,
personnel, facilities, equipment, documentation, materials management,
production controls, packaging, storage, lab controls, validation,
change control, complaints/recalls - Section 18: APIs by cell
culture/fermentation (biotech) - Section 19: APIs for clinical trials
(reduced requirements)

\hypertarget{ich-q8r2-pharmaceutical-development-quality-by-design}{%
\subsection{ICH Q8(R2): Pharmaceutical Development (Quality by
Design)}\label{ich-q8r2-pharmaceutical-development-quality-by-design}}

The foundation of modern CMC development: - \textbf{Quality Target
Product Profile (QTPP)}: Prospective summary of desired quality
characteristics - \textbf{Critical Quality Attributes (CQAs)}:
Properties within appropriate limits for desired quality -
\textbf{Design Space}: Multidimensional combination of input variables
demonstrated to assure quality --- working within design space is NOT a
regulatory change - \textbf{Control Strategy}: Planned controls ensuring
process performance and product quality - \textbf{Process Analytical
Technology (PAT)}: Real-time monitoring and control -
\textbf{Traditional vs.~Enhanced approach}: Enhanced (QbD) enables
greater regulatory flexibility

\hypertarget{ich-q9r1-quality-risk-management}{%
\subsection{ICH Q9(R1): Quality Risk
Management}\label{ich-q9r1-quality-risk-management}}

\begin{itemize}
\tightlist
\item
  \textbf{Process}: Risk Assessment → Risk Control → Risk Review → Risk
  Communication
\item
  \textbf{Tools}: FMEA/FMECA, Fault Tree Analysis, HACCP, HAZOP, Risk
  Ranking and Filtering
\item
  \textbf{R1 revision} (2023): Emphasis on avoiding subjectivity and
  bias, promoting formality
\end{itemize}

\hypertarget{ich-q10-pharmaceutical-quality-system}{%
\subsection{ICH Q10: Pharmaceutical Quality
System}\label{ich-q10-pharmaceutical-quality-system}}

\begin{itemize}
\tightlist
\item
  \textbf{Four PQS elements}: (1) Process Performance and Product
  Quality Monitoring, (2) CAPA System, (3) Change Management, (4)
  Management Review
\item
  \textbf{Enablers}: Knowledge Management, Quality Risk Management
\item
  The ``ICH Q-Trio'' (Q8+Q9+Q10) represents the modern pharmaceutical
  quality paradigm
\end{itemize}

\hypertarget{ich-q11-drug-substance-development-and-manufacture}{%
\subsection{ICH Q11: Drug Substance Development and
Manufacture}\label{ich-q11-drug-substance-development-and-manufacture}}

\begin{itemize}
\tightlist
\item
  Extends Q8 QbD principles to API development
\item
  Starting material justification criteria: well-defined, sufficient
  steps to final API, full GMP from starting material onward
\item
  Manufacturing process development approaches: traditional vs.~enhanced
\item
  Control strategy for drug substance
\end{itemize}

\hypertarget{ich-q12-lifecycle-management}{%
\subsection{ICH Q12: Lifecycle
Management}\label{ich-q12-lifecycle-management}}

\begin{itemize}
\tightlist
\item
  \textbf{Established Conditions (ECs)}: Legally binding commitments ---
  only changes to ECs require regulatory notification
\item
  \textbf{Post-Approval Change Management Protocols (PACMPs)}:
  Pre-agreed change management for anticipated changes
\item
  \textbf{Reporting categories}: Prior approval, notify-and-wait, annual
  report --- harmonized globally
\item
  The regulatory payoff for investment in process understanding
\end{itemize}

\hypertarget{ich-q13-continuous-manufacturing}{%
\subsection{ICH Q13: Continuous
Manufacturing}\label{ich-q13-continuous-manufacturing}}

\begin{itemize}
\tightlist
\item
  Harmonized guidance for continuous manufacturing of drug substances
  and drug products
\item
  Addresses batch definition, process monitoring, control strategy,
  startup/shutdown
\item
  Covers both small molecules and biologics
\end{itemize}

\hypertarget{ich-q14-analytical-procedure-development}{%
\subsection{ICH Q14: Analytical Procedure
Development}\label{ich-q14-analytical-procedure-development}}

\begin{itemize}
\tightlist
\item
  Complements Q2(R2) on validation
\item
  Introduces Analytical QbD, Analytical Design Space, and Proven
  Acceptable Range (PAR)
\item
  Supports lifecycle management of analytical procedures
\end{itemize}

\begin{center}\rule{0.5\linewidth}{0.5pt}\end{center}

\hypertarget{the-common-technical-document-ctd-module-3}{%
\section{4. The Common Technical Document (CTD) Module
3}\label{the-common-technical-document-ctd-module-3}}

Module 3 (Quality) is the CMC section of every regulatory submission
worldwide. Structured per ICH M4Q(R1).

\hypertarget{module-3-structure}{%
\subsection{Module 3 Structure}\label{module-3-structure}}

\begin{verbatim}
Module 3: Quality
├── 3.2.S — Drug Substance (repeated per substance)
│   ├── 3.2.S.1  General Information (nomenclature, structure, properties)
│   ├── 3.2.S.2  Manufacture (process, controls, validation)
│   ├── 3.2.S.3  Characterisation (structure elucidation, impurities)
│   ├── 3.2.S.4  Control of Drug Substance (specs, methods, validation, batch analyses)
│   ├── 3.2.S.5  Reference Standards
│   ├── 3.2.S.6  Container Closure System
│   └── 3.2.S.7  Stability
│
├── 3.2.P — Drug Product (repeated per product/dosage form)
│   ├── 3.2.P.1  Description and Composition
│   ├── 3.2.P.2  Pharmaceutical Development (QbD, design space)
│   ├── 3.2.P.3  Manufacture (batch formula, process, validation)
│   ├── 3.2.P.4  Control of Excipients
│   ├── 3.2.P.5  Control of Drug Product (specs, methods, validation, batch analyses)
│   ├── 3.2.P.6  Reference Standards
│   ├── 3.2.P.7  Container Closure System (E&L data)
│   └── 3.2.P.8  Stability
│
├── 3.2.A — Appendices
│   ├── 3.2.A.1  Facilities and Equipment (esp. biologics)
│   ├── 3.2.A.2  Adventitious Agents Safety Evaluation (biologics)
│   └── 3.2.A.3  Novel Excipients (if applicable)
│
└── 3.2.R — Regional Information
    └── 3.2.R.1-R.6  Region-specific (e.g., executed batch records for FDA)
\end{verbatim}

\begin{center}\rule{0.5\linewidth}{0.5pt}\end{center}

\hypertarget{drug-substance-api-deep-dive}{%
\section{5. Drug Substance (API) --- Deep
Dive}\label{drug-substance-api-deep-dive}}

\hypertarget{synthetic-route-development}{%
\subsection{5.1 Synthetic Route
Development}\label{synthetic-route-development}}

\textbf{Progression from discovery to commercial:}

\begin{longtable}[]{@{}lll@{}}
\toprule
\begin{minipage}[b]{0.21\columnwidth}\raggedright
Stage\strut
\end{minipage} & \begin{minipage}[b]{0.21\columnwidth}\raggedright
Scale\strut
\end{minipage} & \begin{minipage}[b]{0.49\columnwidth}\raggedright
Characteristics\strut
\end{minipage}\tabularnewline
\midrule
\endhead
\begin{minipage}[t]{0.21\columnwidth}\raggedright
Medicinal Chemistry Route\strut
\end{minipage} & \begin{minipage}[t]{0.21\columnwidth}\raggedright
mg--g\strut
\end{minipage} & \begin{minipage}[t]{0.49\columnwidth}\raggedright
8-15 steps, expensive reagents, column chromatography\strut
\end{minipage}\tabularnewline
\begin{minipage}[t]{0.21\columnwidth}\raggedright
First-Generation Process Route\strut
\end{minipage} & \begin{minipage}[t]{0.21\columnwidth}\raggedright
kg\strut
\end{minipage} & \begin{minipage}[t]{0.49\columnwidth}\raggedright
5-10 steps, eliminate chromatography, replace exotic reagents\strut
\end{minipage}\tabularnewline
\begin{minipage}[t]{0.21\columnwidth}\raggedright
Commercial Route\strut
\end{minipage} & \begin{minipage}[t]{0.21\columnwidth}\raggedright
100 kg--multi-ton\strut
\end{minipage} & \begin{minipage}[t]{0.49\columnwidth}\raggedright
3-7 steps, crystallization purifications, telescoped steps, optimized
yield\strut
\end{minipage}\tabularnewline
\bottomrule
\end{longtable}

\textbf{Key optimization activities:} - \textbf{Reaction screening}:
High-throughput experimentation (HTE) platforms, DoE to map parameter
space - \textbf{Solvent selection}: Replace Class 1/2 solvents with
greener alternatives (2-MeTHF, ethyl acetate, iPrOH) - \textbf{Catalyst
optimization}: Reduce Pd loading to \textless100 ppm; develop ligand
systems enabling milder conditions - \textbf{Telescoping}: Combine
consecutive steps without intermediate isolation -
\textbf{Polymorph-directed crystallization}: Seeding strategy, cooling
rate, anti-solvent addition control

\hypertarget{starting-materials-ich-q11}{%
\subsection{5.2 Starting Materials (ICH
Q11)}\label{starting-materials-ich-q11}}

Starting materials define the regulatory GMP boundary: - Must be
well-defined substances with established specifications - Typically 3-5
synthetic steps from final API - Commercial availability from multiple
suppliers preferred - FDA routinely pushes back on starting materials
too close to the API

\textbf{Starting material specifications:} Identity (IR, NMR, or MS),
assay (\textgreater95\% by HPLC), specified impurities
(\textgreater0.10\%), residual solvents, chiral purity (if applicable,
\textgreater99\% ee).

\hypertarget{critical-process-parameters-cpps-and-critical-quality-attributes-cqas}{%
\subsection{5.3 Critical Process Parameters (CPPs) and Critical Quality
Attributes
(CQAs)}\label{critical-process-parameters-cpps-and-critical-quality-attributes-cqas}}

\textbf{Typical Drug Substance CQAs:}

\begin{longtable}[]{@{}lll@{}}
\toprule
CQA & Typical Target & Test Method\tabularnewline
\midrule
\endhead
Assay & 98.0-102.0\% & HPLC\tabularnewline
Specified impurities & Each NMT 0.15\% & HPLC-UV/MS\tabularnewline
Total impurities & NMT 1.0-2.0\% & HPLC\tabularnewline
Residual solvents & Per ICH Q3C & GC-HS\tabularnewline
Water content & NMT 0.5\% (typical) & Karl Fischer\tabularnewline
Particle size & D90 \textless{} specified & Laser
diffraction\tabularnewline
Polymorphic form & Confirmed Form I & XRPD, DSC\tabularnewline
Elemental impurities & Per ICH Q3D & ICP-MS\tabularnewline
Residual catalysts & Pd \textless{} 10 ppm (oral) &
ICP-MS\tabularnewline
Chiral purity & \textgreater99.5\% ee & Chiral HPLC/SFC\tabularnewline
Microbial limits & Per USP
\textless61\textgreater/\textless62\textgreater{} & Membrane
filtration\tabularnewline
\bottomrule
\end{longtable}

\textbf{CPP examples by reaction type:} - \textbf{Coupling reactions}:
Temperature, catalyst loading, equivalents, reaction time, inert
atmosphere - \textbf{Crystallization}: Cooling rate (°C/hour), seed
loading (0.1-1.0\% w/w), anti-solvent addition rate, agitation speed -
\textbf{Salt formation}: Stoichiometry, addition order, solvent system,
temperature profile

\hypertarget{process-characterization}{%
\subsection{5.4 Process
Characterization}\label{process-characterization}}

\begin{enumerate}
\def\labelenumi{\arabic{enumi}.}
\tightlist
\item
  \textbf{Risk Assessment}: Ishikawa diagrams → FMEA scoring (severity ×
  probability × detectability) → prioritize high-RPN parameters
\item
  \textbf{Screening DoE}: Fractional factorial or Plackett-Burman, 12-20
  experiments per unit operation
\item
  \textbf{Optimization DoE}: Full factorial, central composite, or
  Box-Behnken to build response surface models
\item
  \textbf{Establish ranges}: Proven Acceptable Range (PAR) → Normal
  Operating Range (NOR, narrower) → Design Space (multidimensional, ICH
  Q8)
\end{enumerate}

\hypertarget{impurity-profiling}{%
\subsection{5.5 Impurity Profiling}\label{impurity-profiling}}

\textbf{Classification:} - \textbf{Process-related}: Starting material
carryover, reagent/catalyst residues, by-products, intermediates -
\textbf{Degradation products}: Hydrolysis, oxidation, photodegradation,
thermal degradation - \textbf{Elemental impurities}: Per ICH Q3D class
system - \textbf{Mutagenic impurities}: Per ICH M7 --- TTC of 1.5 µg/day
(lifetime)

\textbf{Forced Degradation (Stress Testing) Protocol:}

\begin{longtable}[]{@{}lll@{}}
\toprule
Stress Condition & Parameters & Duration\tabularnewline
\midrule
\endhead
Acid hydrolysis & 0.1-1N HCl, 60-80°C & 1-7 days\tabularnewline
Base hydrolysis & 0.1-1N NaOH, 60-80°C & 1-7 days\tabularnewline
Oxidation & 3\% H₂O₂, RT-40°C & 1-7 days\tabularnewline
Thermal & 60-80°C, solid state & 1-4 weeks\tabularnewline
Photolytic & ICH Q1B (1.2M lux-hours, 200 Wh/m²) & Per
protocol\tabularnewline
Humidity & 75\% RH, 40°C & 1-4 weeks\tabularnewline
\bottomrule
\end{longtable}

Target: 5-20\% degradation to demonstrate method specificity.

\hypertarget{polymorphism-and-solid-state-characterization}{%
\subsection{5.6 Polymorphism and Solid-State
Characterization}\label{polymorphism-and-solid-state-characterization}}

\textbf{Screening campaign:} Crystallization from 20-40 diverse solvents
(cooling, anti-solvent, evaporation, slurry conversion), mechanical
screening (grinding, ball milling), thermal screening, high-throughput
automated platforms.

\textbf{Characterization techniques:}

\begin{longtable}[]{@{}ll@{}}
\toprule
\begin{minipage}[b]{0.52\columnwidth}\raggedright
Technique\strut
\end{minipage} & \begin{minipage}[b]{0.42\columnwidth}\raggedright
Purpose\strut
\end{minipage}\tabularnewline
\midrule
\endhead
\begin{minipage}[t]{0.52\columnwidth}\raggedright
XRPD\strut
\end{minipage} & \begin{minipage}[t]{0.42\columnwidth}\raggedright
Primary polymorph identification (unique diffraction pattern per
form)\strut
\end{minipage}\tabularnewline
\begin{minipage}[t]{0.52\columnwidth}\raggedright
DSC\strut
\end{minipage} & \begin{minipage}[t]{0.42\columnwidth}\raggedright
Melting points, polymorphic transitions, Tg (amorphous)\strut
\end{minipage}\tabularnewline
\begin{minipage}[t]{0.52\columnwidth}\raggedright
TGA\strut
\end{minipage} & \begin{minipage}[t]{0.42\columnwidth}\raggedright
Desolvation/decomposition weight loss\strut
\end{minipage}\tabularnewline
\begin{minipage}[t]{0.52\columnwidth}\raggedright
DVS\strut
\end{minipage} & \begin{minipage}[t]{0.42\columnwidth}\raggedright
Hygroscopicity profiling\strut
\end{minipage}\tabularnewline
\begin{minipage}[t]{0.52\columnwidth}\raggedright
Single Crystal XRD\strut
\end{minipage} & \begin{minipage}[t]{0.42\columnwidth}\raggedright
Definitive structure determination\strut
\end{minipage}\tabularnewline
\begin{minipage}[t]{0.52\columnwidth}\raggedright
Solid-state NMR\strut
\end{minipage} & \begin{minipage}[t]{0.42\columnwidth}\raggedright
Polymorph distinction, disorder, amorphous content\strut
\end{minipage}\tabularnewline
\begin{minipage}[t]{0.52\columnwidth}\raggedright
Hot Stage Microscopy\strut
\end{minipage} & \begin{minipage}[t]{0.42\columnwidth}\raggedright
Visual phase transition observation\strut
\end{minipage}\tabularnewline
\bottomrule
\end{longtable}

\textbf{Regulatory expectation (ICH Q6A Decision Tree \#4):} If
polymorphism exists, drug substance specification must include polymorph
identity test.

\hypertarget{reference-standards}{%
\subsection{5.7 Reference Standards}\label{reference-standards}}

\textbf{Primary reference standard} --- characterized by: - Quantitative
NMR (qNMR) for absolute purity (traceable to SI units) - Mass balance:
Purity = 100\% − impurities − water − residual solvents − inorganic
content - Multiple orthogonal techniques: HPLC, chiral HPLC, elemental
analysis, spectroscopic identity - Storage under controlled conditions
(typically −20°C or 2-8°C, desiccated) - Annual retest with defined
criteria

\textbf{Working/secondary reference standards} --- qualified against
primary by HPLC assay, used for routine testing.

\begin{center}\rule{0.5\linewidth}{0.5pt}\end{center}

\hypertarget{drug-product-deep-dive}{%
\section{6. Drug Product --- Deep Dive}\label{drug-product-deep-dive}}

\hypertarget{formulation-development-strategies}{%
\subsection{6.1 Formulation Development
Strategies}\label{formulation-development-strategies}}

\textbf{Pre-formulation characterization:} - Solubility: pH-solubility
profile (pH 1-8), biorelevant media (FaSSIF, FeSSIF) - BCS
Classification: Class I (high sol/high perm), II (low sol/high perm),
III (high sol/low perm), IV (low sol/low perm) - pKa, Log P/Log D,
intrinsic dissolution rate - Hygroscopicity (DVS), mechanical properties
(compressibility, flowability) - Excipient compatibility (binary
mixtures at 40°C/75\% RH, 2-4 weeks)

\textbf{Formulation approach by BCS Class:}

\begin{longtable}[]{@{}lll@{}}
\toprule
\begin{minipage}[b]{0.31\columnwidth}\raggedright
BCS Class\strut
\end{minipage} & \begin{minipage}[b]{0.31\columnwidth}\raggedright
Challenge\strut
\end{minipage} & \begin{minipage}[b]{0.29\columnwidth}\raggedright
Strategy\strut
\end{minipage}\tabularnewline
\midrule
\endhead
\begin{minipage}[t]{0.31\columnwidth}\raggedright
I\strut
\end{minipage} & \begin{minipage}[t]{0.31\columnwidth}\raggedright
None major\strut
\end{minipage} & \begin{minipage}[t]{0.29\columnwidth}\raggedright
Standard IR tablets, direct compression\strut
\end{minipage}\tabularnewline
\begin{minipage}[t]{0.31\columnwidth}\raggedright
II\strut
\end{minipage} & \begin{minipage}[t]{0.31\columnwidth}\raggedright
Low solubility\strut
\end{minipage} & \begin{minipage}[t]{0.29\columnwidth}\raggedright
Amorphous solid dispersions (HME, spray drying), nanocrystals, lipid
formulations, salt/cocrystal screening, cyclodextrin complexation\strut
\end{minipage}\tabularnewline
\begin{minipage}[t]{0.31\columnwidth}\raggedright
III\strut
\end{minipage} & \begin{minipage}[t]{0.31\columnwidth}\raggedright
Low permeability\strut
\end{minipage} & \begin{minipage}[t]{0.29\columnwidth}\raggedright
Surfactant addition, tight disintegration control\strut
\end{minipage}\tabularnewline
\begin{minipage}[t]{0.31\columnwidth}\raggedright
IV\strut
\end{minipage} & \begin{minipage}[t]{0.31\columnwidth}\raggedright
Both\strut
\end{minipage} & \begin{minipage}[t]{0.29\columnwidth}\raggedright
Combination of solubility + permeability enhancement\strut
\end{minipage}\tabularnewline
\bottomrule
\end{longtable}

\textbf{Key excipient compatibility issues to watch:} - Amine drugs +
lactose = Maillard reaction (browning) - Acid-sensitive APIs + alkaline
excipients - Oxidation-sensitive APIs + peroxide-containing excipients
(PEG, povidone)

\hypertarget{manufacturing-processes-oral-solids}{%
\subsection{6.2 Manufacturing Processes --- Oral
Solids}\label{manufacturing-processes-oral-solids}}

\textbf{Direct Compression} --- preferred when API has adequate flow
(Carr index \textless25): - Blend uniformity is primary challenge -
Stratified sampling: 10 locations × 3 samples, RSD \textless{} 5\%,
individuals within 90-110\%

\textbf{Wet Granulation:} - High-shear: Impeller speed, chopper speed,
binder addition rate, endpoint by power consumption/torque - Fluid bed:
Inlet air temperature, spray rate, droplet size, air flow rate - Drying:
Fluid bed to LOD 1-3\% w/w - Milling: Comil screen size, impeller speed

\textbf{Dry Granulation (Roller Compaction):} - For
moisture/heat-sensitive APIs - Roll force, gap, speed; target ribbon
solid fraction 0.6-0.8

\textbf{Tableting (Rotary Press):} - Parameters: Compression force,
pre-compression, press speed (dwell time), tooling - CQAs: Weight,
hardness, thickness, friability (\textless1.0\%), disintegration,
dissolution, content uniformity - Common defects: Capping, lamination,
sticking/picking, weight variation

\textbf{Film Coating:} - Pan coating: Inlet air temp, spray rate, pan
speed, atomization pressure - Typical weight gain: 2-4\%
(aesthetics/moisture), 8-15\% (functional coatings)

\hypertarget{injectable-formulations}{%
\subsection{6.3 Injectable Formulations}\label{injectable-formulations}}

\textbf{Key considerations:} - Sterility assurance: Terminal
sterilization (121°C/15 min, preferred) vs.~aseptic processing -
Tonicity: 280-310 mOsm/kg for IV - pH: Target 3-9 for IV, ideally 5-8 -
Viscosity: \textless20 cP for IV push, \textless50 cP for SC

\textbf{Formulation types:} - \textbf{Solutions}: Co-solvents (PEG
300/400, propylene glycol), cyclodextrin complexation (Captisol,
HP-β-CD) - \textbf{Lyophilized}: Bulking agents (mannitol, glycine),
cryoprotectants (sucrose, trehalose). Cycle: freezing → primary drying
(below collapse temperature) → secondary drying (25-40°C, target
residual moisture \textless1\% for proteins) - \textbf{Suspensions}: For
poorly soluble IM/SC depots

\hypertarget{biologic-drug-product-formulation}{%
\subsection{6.4 Biologic Drug Product
Formulation}\label{biologic-drug-product-formulation}}

\begin{itemize}
\tightlist
\item
  Buffer selection: Histidine (10-50 mM, preferred for mAbs), phosphate,
  citrate
\item
  Stabilizers: Sucrose (1-10\%), trehalose
\item
  Surfactants: Polysorbate 20/80 (0.01-0.1\%)
\item
  High-concentration formulations (\textgreater100 mg/mL for SC):
  Viscosity reduction with arginine, proline
\item
  Stress studies: Thermal, agitation, freeze-thaw, light, oxidative ---
  monitor by SEC, CEX, CE-SDS, turbidity, sub-visible particles
\end{itemize}

\hypertarget{scale-up}{%
\subsection{6.5 Scale-Up}\label{scale-up}}

\begin{longtable}[]{@{}lll@{}}
\toprule
\begin{minipage}[b]{0.32\columnwidth}\raggedright
Unit Operation\strut
\end{minipage} & \begin{minipage}[b]{0.39\columnwidth}\raggedright
Scale-Up Parameter\strut
\end{minipage} & \begin{minipage}[b]{0.20\columnwidth}\raggedright
Guidance\strut
\end{minipage}\tabularnewline
\midrule
\endhead
\begin{minipage}[t]{0.32\columnwidth}\raggedright
Blending\strut
\end{minipage} & \begin{minipage}[t]{0.39\columnwidth}\raggedright
Froude number\strut
\end{minipage} & \begin{minipage}[t]{0.20\columnwidth}\raggedright
Keep constant; adjust RPM by √(D\_small/D\_large)\strut
\end{minipage}\tabularnewline
\begin{minipage}[t]{0.32\columnwidth}\raggedright
High-shear granulation\strut
\end{minipage} & \begin{minipage}[t]{0.39\columnwidth}\raggedright
Impeller tip speed\strut
\end{minipage} & \begin{minipage}[t]{0.20\columnwidth}\raggedright
Keep constant (5-10 m/s)\strut
\end{minipage}\tabularnewline
\begin{minipage}[t]{0.32\columnwidth}\raggedright
Fluid bed drying\strut
\end{minipage} & \begin{minipage}[t]{0.39\columnwidth}\raggedright
Air velocity\strut
\end{minipage} & \begin{minipage}[t]{0.20\columnwidth}\raggedright
Scale by cross-sectional area\strut
\end{minipage}\tabularnewline
\begin{minipage}[t]{0.32\columnwidth}\raggedright
Roller compaction\strut
\end{minipage} & \begin{minipage}[t]{0.39\columnwidth}\raggedright
Specific roll force (kN/cm)\strut
\end{minipage} & \begin{minipage}[t]{0.20\columnwidth}\raggedright
Keep constant\strut
\end{minipage}\tabularnewline
\begin{minipage}[t]{0.32\columnwidth}\raggedright
Tableting\strut
\end{minipage} & \begin{minipage}[t]{0.39\columnwidth}\raggedright
Compression force \& dwell time\strut
\end{minipage} & \begin{minipage}[t]{0.20\columnwidth}\raggedright
Match dwell time at higher speeds\strut
\end{minipage}\tabularnewline
\begin{minipage}[t]{0.32\columnwidth}\raggedright
Film coating\strut
\end{minipage} & \begin{minipage}[t]{0.39\columnwidth}\raggedright
Spray rate per surface area\strut
\end{minipage} & \begin{minipage}[t]{0.20\columnwidth}\raggedright
Scale by pan loading\strut
\end{minipage}\tabularnewline
\begin{minipage}[t]{0.32\columnwidth}\raggedright
Lyophilization\strut
\end{minipage} & \begin{minipage}[t]{0.39\columnwidth}\raggedright
Shelf temp \& chamber pressure\strut
\end{minipage} & \begin{minipage}[t]{0.20\columnwidth}\raggedright
Same parameters; edge/center effects increase\strut
\end{minipage}\tabularnewline
\bottomrule
\end{longtable}

\textbf{Scale stages:} Lab (0.5-2 kg) → Pilot (10-50 kg) → Commercial
(100-500+ kg)

\hypertarget{quality-by-design-qbd-workflow-ich-q8}{%
\subsection{6.6 Quality by Design (QbD) Workflow --- ICH
Q8}\label{quality-by-design-qbd-workflow-ich-q8}}

\begin{enumerate}
\def\labelenumi{\arabic{enumi}.}
\tightlist
\item
  Define \textbf{QTPP} (dosage form, route, strength, pharmacokinetic
  target, stability)
\item
  Identify \textbf{CQAs} from QTPP (dissolution, assay, CU, degradants)
\item
  \textbf{Risk Assessment} linking material attributes and process
  parameters to CQAs
\item
  \textbf{DoE} to systematically study impacts
\item
  Establish \textbf{Design Space} (multidimensional region where quality
  is assured)
\item
  Define \textbf{Control Strategy} (input controls + process controls +
  testing)
\item
  \textbf{Lifecycle management} per ICH Q10/Q12
\end{enumerate}

\hypertarget{process-analytical-technology-pat}{%
\subsection{6.7 Process Analytical Technology
(PAT)}\label{process-analytical-technology-pat}}

\begin{longtable}[]{@{}lll@{}}
\toprule
\begin{minipage}[b]{0.23\columnwidth}\raggedright
PAT Tool\strut
\end{minipage} & \begin{minipage}[b]{0.30\columnwidth}\raggedright
Application\strut
\end{minipage} & \begin{minipage}[b]{0.39\columnwidth}\raggedright
What It Measures\strut
\end{minipage}\tabularnewline
\midrule
\endhead
\begin{minipage}[t]{0.23\columnwidth}\raggedright
NIR spectroscopy\strut
\end{minipage} & \begin{minipage}[t]{0.30\columnwidth}\raggedright
Blend uniformity, moisture, coating thickness\strut
\end{minipage} & \begin{minipage}[t]{0.39\columnwidth}\raggedright
Chemical and physical attributes\strut
\end{minipage}\tabularnewline
\begin{minipage}[t]{0.23\columnwidth}\raggedright
Raman spectroscopy\strut
\end{minipage} & \begin{minipage}[t]{0.30\columnwidth}\raggedright
Polymorph ID, API content\strut
\end{minipage} & \begin{minipage}[t]{0.39\columnwidth}\raggedright
Molecular vibrations\strut
\end{minipage}\tabularnewline
\begin{minipage}[t]{0.23\columnwidth}\raggedright
In-line HPLC/UPLC\strut
\end{minipage} & \begin{minipage}[t]{0.30\columnwidth}\raggedright
Reaction monitoring\strut
\end{minipage} & \begin{minipage}[t]{0.39\columnwidth}\raggedright
Chemical composition\strut
\end{minipage}\tabularnewline
\begin{minipage}[t]{0.23\columnwidth}\raggedright
FBRM\strut
\end{minipage} & \begin{minipage}[t]{0.30\columnwidth}\raggedright
Crystallization, granulation\strut
\end{minipage} & \begin{minipage}[t]{0.39\columnwidth}\raggedright
Chord length distribution\strut
\end{minipage}\tabularnewline
\begin{minipage}[t]{0.23\columnwidth}\raggedright
Acoustic emission\strut
\end{minipage} & \begin{minipage}[t]{0.30\columnwidth}\raggedright
Granulation endpoint\strut
\end{minipage} & \begin{minipage}[t]{0.39\columnwidth}\raggedright
Sound energy\strut
\end{minipage}\tabularnewline
\begin{minipage}[t]{0.23\columnwidth}\raggedright
Microwave resonance\strut
\end{minipage} & \begin{minipage}[t]{0.30\columnwidth}\raggedright
Moisture content\strut
\end{minipage} & \begin{minipage}[t]{0.39\columnwidth}\raggedright
Dielectric properties\strut
\end{minipage}\tabularnewline
\bottomrule
\end{longtable}

\begin{center}\rule{0.5\linewidth}{0.5pt}\end{center}

\hypertarget{analytical-methods-deep-dive}{%
\section{7. Analytical Methods --- Deep
Dive}\label{analytical-methods-deep-dive}}

\hypertarget{types-of-analytical-methods-in-cmc}{%
\subsection{7.1 Types of Analytical Methods in
CMC}\label{types-of-analytical-methods-in-cmc}}

\begin{longtable}[]{@{}lll@{}}
\toprule
\begin{minipage}[b]{0.26\columnwidth}\raggedright
Category\strut
\end{minipage} & \begin{minipage}[b]{0.23\columnwidth}\raggedright
Purpose\strut
\end{minipage} & \begin{minipage}[b]{0.42\columnwidth}\raggedright
Common Methods\strut
\end{minipage}\tabularnewline
\midrule
\endhead
\begin{minipage}[t]{0.26\columnwidth}\raggedright
\textbf{Identity}\strut
\end{minipage} & \begin{minipage}[t]{0.23\columnwidth}\raggedright
Confirm correct molecule\strut
\end{minipage} & \begin{minipage}[t]{0.42\columnwidth}\raggedright
IR, UV, HPLC retention time, MS, specific rotation\strut
\end{minipage}\tabularnewline
\begin{minipage}[t]{0.26\columnwidth}\raggedright
\textbf{Assay (Strength)}\strut
\end{minipage} & \begin{minipage}[t]{0.23\columnwidth}\raggedright
Quantitate active ingredient\strut
\end{minipage} & \begin{minipage}[t]{0.42\columnwidth}\raggedright
HPLC (UV, RI), UV spectrophotometry\strut
\end{minipage}\tabularnewline
\begin{minipage}[t]{0.26\columnwidth}\raggedright
\textbf{Purity}\strut
\end{minipage} & \begin{minipage}[t]{0.23\columnwidth}\raggedright
Detect/quantitate impurities\strut
\end{minipage} & \begin{minipage}[t]{0.42\columnwidth}\raggedright
HPLC, GC, CE, ICP-MS, LC-MS/MS\strut
\end{minipage}\tabularnewline
\begin{minipage}[t]{0.26\columnwidth}\raggedright
\textbf{Potency}\strut
\end{minipage} & \begin{minipage}[t]{0.23\columnwidth}\raggedright
Measure biological activity\strut
\end{minipage} & \begin{minipage}[t]{0.42\columnwidth}\raggedright
Bioassays, cell-based assays (biologics)\strut
\end{minipage}\tabularnewline
\begin{minipage}[t]{0.26\columnwidth}\raggedright
\textbf{Physical}\strut
\end{minipage} & \begin{minipage}[t]{0.23\columnwidth}\raggedright
Characterize physical properties\strut
\end{minipage} & \begin{minipage}[t]{0.42\columnwidth}\raggedright
Dissolution, particle size, hardness, friability\strut
\end{minipage}\tabularnewline
\begin{minipage}[t]{0.26\columnwidth}\raggedright
\textbf{Microbiology}\strut
\end{minipage} & \begin{minipage}[t]{0.23\columnwidth}\raggedright
Ensure microbial quality\strut
\end{minipage} & \begin{minipage}[t]{0.42\columnwidth}\raggedright
Sterility (USP \textless71\textgreater), endotoxins (USP
\textless85\textgreater), bioburden (USP
\textless61\textgreater/\textless62\textgreater)\strut
\end{minipage}\tabularnewline
\begin{minipage}[t]{0.26\columnwidth}\raggedright
\textbf{Content Uniformity}\strut
\end{minipage} & \begin{minipage}[t]{0.23\columnwidth}\raggedright
Dose-to-dose consistency\strut
\end{minipage} & \begin{minipage}[t]{0.42\columnwidth}\raggedright
HPLC per USP \textless905\textgreater{}\strut
\end{minipage}\tabularnewline
\bottomrule
\end{longtable}

\hypertarget{method-development-and-validation-ich-q2r2}{%
\subsection{7.2 Method Development and Validation (ICH
Q2(R2))}\label{method-development-and-validation-ich-q2r2}}

\textbf{Method development phases:} 1. \textbf{Method scouting}: Screen
columns, mobile phases, detectors 2. \textbf{Method optimization}:
Systematic optimization of critical method parameters (DoE approach
encouraged in Q14) 3. \textbf{Method qualification}: Demonstrate fitness
for purpose (sufficient for early clinical phases) 4. \textbf{Method
validation}: Full ICH Q2 validation (required for NDA/BLA) 5.
\textbf{Method transfer}: Transfer to QC labs at manufacturing sites 6.
\textbf{Method verification}: Confirm transferred method performs
adequately at receiving site

\hypertarget{compendial-vs.-non-compendial-methods}{%
\subsection{7.3 Compendial vs.~Non-Compendial
Methods}\label{compendial-vs.-non-compendial-methods}}

\begin{itemize}
\tightlist
\item
  \textbf{Compendial} (USP/Ph.Eur./JP): Standard methods in
  pharmacopeial monographs. Require verification, not full validation,
  when used as published.
\item
  \textbf{Non-compendial}: Custom methods developed in-house. Require
  full ICH Q2 validation.
\item
  \textbf{Stability-indicating methods} must resolve drug from all
  degradation products --- demonstrated by forced degradation studies
  and peak purity analysis (PDA).
\end{itemize}

\hypertarget{specifications-setting-ich-q6aq6b}{%
\subsection{7.4 Specifications Setting (ICH
Q6A/Q6B)}\label{specifications-setting-ich-q6aq6b}}

Specifications are justified by: - Pharmacopeial standards - Batch data
(clinical and registration batches) - Stability data - Development data
(process capability) - Toxicological qualification of impurities -
Clinical relevance (e.g., dissolution linked to bioavailability)

\textbf{Release vs.~stability specifications:} Release specifications
are typically tighter, allowing for degradation over the shelf life. For
example, assay at release might be 98.0-102.0\% while the stability
(shelf-life) specification is 95.0-105.0\%.

\hypertarget{analytical-method-lifecycle-management-ich-q14}{%
\subsection{7.5 Analytical Method Lifecycle Management (ICH
Q14)}\label{analytical-method-lifecycle-management-ich-q14}}

\begin{itemize}
\tightlist
\item
  Introduces \textbf{Analytical QbD}: Define Analytical Target Profile
  (ATP), identify critical method parameters, establish Analytical
  Design Space
\item
  \textbf{Proven Acceptable Range (PAR)}: Operating range within which
  method performance is demonstrated
\item
  Supports continuous improvement of methods within the ATP/design space
  without regulatory filing
\end{itemize}

\hypertarget{method-transfer-and-technology-transfer}{%
\subsection{7.6 Method Transfer and Technology
Transfer}\label{method-transfer-and-technology-transfer}}

\textbf{Method transfer approaches:} 1. \textbf{Comparative testing}:
Sending and receiving labs analyze same samples; statistical comparison
(equivalence testing or t-test) 2. \textbf{Co-validation}: Both labs
validate independently using same protocol 3. \textbf{Revalidation}:
Receiving lab performs full validation 4. \textbf{Transfer waiver}: For
compendial methods or simple methods (e.g., pH, KF)

\textbf{Documentation}: Transfer protocol, transfer report, acceptance
criteria (typically receiving lab within ±2-3\% of sending lab mean, or
equivalent precision).

\begin{center}\rule{0.5\linewidth}{0.5pt}\end{center}

\hypertarget{container-closure-systems}{%
\section{8. Container Closure Systems}\label{container-closure-systems}}

\hypertarget{primary-packaging-by-dosage-form}{%
\subsection{8.1 Primary Packaging by Dosage
Form}\label{primary-packaging-by-dosage-form}}

\begin{longtable}[]{@{}llll@{}}
\toprule
\begin{minipage}[b]{0.22\columnwidth}\raggedright
Dosage Form\strut
\end{minipage} & \begin{minipage}[b]{0.19\columnwidth}\raggedright
Container\strut
\end{minipage} & \begin{minipage}[b]{0.15\columnwidth}\raggedright
Closure\strut
\end{minipage} & \begin{minipage}[b]{0.32\columnwidth}\raggedright
Key Considerations\strut
\end{minipage}\tabularnewline
\midrule
\endhead
\begin{minipage}[t]{0.22\columnwidth}\raggedright
Oral solids\strut
\end{minipage} & \begin{minipage}[t]{0.19\columnwidth}\raggedright
HDPE bottles, PVC/Alu or Alu/Alu blisters\strut
\end{minipage} & \begin{minipage}[t]{0.15\columnwidth}\raggedright
Induction-sealed caps, push-through foil\strut
\end{minipage} & \begin{minipage}[t]{0.32\columnwidth}\raggedright
MVTR, child resistance\strut
\end{minipage}\tabularnewline
\begin{minipage}[t]{0.22\columnwidth}\raggedright
IV solutions\strut
\end{minipage} & \begin{minipage}[t]{0.19\columnwidth}\raggedright
Type I borosilicate glass vials, BFS\strut
\end{minipage} & \begin{minipage}[t]{0.15\columnwidth}\raggedright
Bromobutyl stoppers + Al crimp\strut
\end{minipage} & \begin{minipage}[t]{0.32\columnwidth}\raggedright
Extractables, delamination\strut
\end{minipage}\tabularnewline
\begin{minipage}[t]{0.22\columnwidth}\raggedright
Pre-filled syringes\strut
\end{minipage} & \begin{minipage}[t]{0.19\columnwidth}\raggedright
Type I glass, COP\strut
\end{minipage} & \begin{minipage}[t]{0.15\columnwidth}\raggedright
Bromobutyl/chlorobutyl plungers\strut
\end{minipage} & \begin{minipage}[t]{0.32\columnwidth}\raggedright
Silicone oil, tungsten, glue\strut
\end{minipage}\tabularnewline
\begin{minipage}[t]{0.22\columnwidth}\raggedright
Lyophilized\strut
\end{minipage} & \begin{minipage}[t]{0.19\columnwidth}\raggedright
Type I glass vials\strut
\end{minipage} & \begin{minipage}[t]{0.15\columnwidth}\raggedright
Lyophilization stoppers\strut
\end{minipage} & \begin{minipage}[t]{0.32\columnwidth}\raggedright
Moisture ingress, CCI\strut
\end{minipage}\tabularnewline
\begin{minipage}[t]{0.22\columnwidth}\raggedright
Biologics (mAbs)\strut
\end{minipage} & \begin{minipage}[t]{0.19\columnwidth}\raggedright
Type I glass, COP vials/PFS\strut
\end{minipage} & \begin{minipage}[t]{0.15\columnwidth}\raggedright
Fluoropolymer-coated stoppers\strut
\end{minipage} & \begin{minipage}[t]{0.32\columnwidth}\raggedright
Protein adsorption, particles\strut
\end{minipage}\tabularnewline
\bottomrule
\end{longtable}

\hypertarget{extractables-and-leachables-el}{%
\subsection{8.2 Extractables and Leachables
(E\&L)}\label{extractables-and-leachables-el}}

\textbf{Extractables studies} (on components under exaggerated
conditions): - Solvents: Water, saline, acidic, basic, organic (IPA,
hexane) - Methods: GC-MS (volatile organics), LC-MS/UV (non-volatiles),
ICP-MS (elemental), IC (inorganic), TOC - Analytical Evaluation
Threshold (AET): Based on Safety Concern Threshold --- 0.15 µg/day
(parenteral), 1.5 µg/day (oral solid)

\textbf{Leachables studies} (on drug product in actual container over
shelf life): - Leachables should be a subset of extractables - Monitored
at 0, 3, 6, 12, 18, 24, 36 months on stability

\hypertarget{glass-vs.-polymer-for-biologics}{%
\subsection{8.3 Glass vs.~Polymer for
Biologics}\label{glass-vs.-polymer-for-biologics}}

\begin{longtable}[]{@{}lll@{}}
\toprule
Property & Type I Glass & COP/COC\tabularnewline
\midrule
\endhead
Breakage & Risk & Resistant\tabularnewline
Delamination & Possible (FDA advisory 2011) & None\tabularnewline
Protein adsorption & Higher & Lower\tabularnewline
Moisture barrier & Excellent & Lower (mitigated by SiOx
coating)\tabularnewline
O₂ permeability & Very low & Higher\tabularnewline
Weight & Heavy & Light\tabularnewline
Cost & Lower & Higher\tabularnewline
Extractables & Ions & Organic compounds\tabularnewline
\bottomrule
\end{longtable}

\hypertarget{container-closure-integrity-cci}{%
\subsection{8.4 Container Closure Integrity
(CCI)}\label{container-closure-integrity-cci}}

\begin{itemize}
\tightlist
\item
  Critical for sterile products --- maintains sterility through shelf
  life
\item
  \textbf{Deterministic methods preferred} (per USP
  \textless1207\textgreater): Laser-based headspace, HVLD, vacuum decay,
  pressure decay
\item
  \textbf{Probabilistic methods} (dye/microbial ingress): Being phased
  out
\item
  CCI tested at initial and throughout stability
\end{itemize}

\begin{center}\rule{0.5\linewidth}{0.5pt}\end{center}

\hypertarget{stability-studies}{%
\section{9. Stability Studies}\label{stability-studies}}

\hypertarget{study-design-ich-q1a-q1e}{%
\subsection{9.1 Study Design (ICH
Q1A-Q1E)}\label{study-design-ich-q1a-q1e}}

\textbf{Batches:} Minimum 3 batches, pilot scale or larger,
representative of commercial process.

\textbf{Stability-indicating test panel (oral solid example):}

\begin{longtable}[]{@{}ll@{}}
\toprule
Test & Frequency\tabularnewline
\midrule
\endhead
Appearance & Every time point\tabularnewline
Assay (HPLC) & Every time point\tabularnewline
Degradation products & Every time point\tabularnewline
Dissolution & Every time point\tabularnewline
Water content (KF) & Every time point\tabularnewline
Microbial limits & Selected time points\tabularnewline
\bottomrule
\end{longtable}

\textbf{For biologics, additional tests:} SEC (aggregation), CE-SDS
(fragmentation), charge variants (CEX/icIEF), potency (bioassay),
sub-visible particles (MFI).

\hypertarget{shelf-life-determination}{%
\subsection{9.2 Shelf Life
Determination}\label{shelf-life-determination}}

\begin{enumerate}
\def\labelenumi{\arabic{enumi}.}
\tightlist
\item
  \textbf{At filing}: 24-36 months proposed based on 12-18 months
  real-time + accelerated + statistical extrapolation (ICH Q1E)
\item
  \textbf{Extrapolation rules}: Up to 2× long-term data period or 12
  months beyond, whichever is shorter --- only if accelerated shows no
  significant change
\item
  \textbf{Post-approval}: Continue stability to cover proposed shelf
  life; annual batches (≥1/year/strength)
\end{enumerate}

\hypertarget{stability-commitments-in-submission}{%
\subsection{9.3 Stability Commitments in
Submission}\label{stability-commitments-in-submission}}

\begin{itemize}
\tightlist
\item
  Complete ongoing studies on registration batches
\item
  Place first 3 commercial batches on stability
\item
  Annual stability: 1 batch/year/strength/pack configuration
\item
  Report confirmed OOS results to agency
\item
  Shelf life extension via supplement with supporting long-term data
\end{itemize}

\begin{center}\rule{0.5\linewidth}{0.5pt}\end{center}

\hypertarget{process-validation}{%
\section{10. Process Validation}\label{process-validation}}

\hypertarget{fda-2011-lifecycle-approach-three-stages}{%
\subsection{10.1 FDA 2011 Lifecycle Approach --- Three
Stages}\label{fda-2011-lifecycle-approach-three-stages}}

\textbf{Stage 1 --- Process Design:} - DoE to characterize CPP--CQA
relationships - Establish PAR (Proven Acceptable Range) and NOR (Normal
Operating Range) - Define design space and control strategy - Technology
transfer package - Deliverables: Process development report, risk
assessments, design space documentation

\textbf{Stage 2 --- Process Qualification:}

\emph{Equipment Qualification:}

\begin{longtable}[]{@{}ll@{}}
\toprule
\begin{minipage}[b]{0.41\columnwidth}\raggedright
Phase\strut
\end{minipage} & \begin{minipage}[b]{0.53\columnwidth}\raggedright
Purpose\strut
\end{minipage}\tabularnewline
\midrule
\endhead
\begin{minipage}[t]{0.41\columnwidth}\raggedright
DQ (Design Qualification)\strut
\end{minipage} & \begin{minipage}[t]{0.53\columnwidth}\raggedright
Verify proposed design meets URS\strut
\end{minipage}\tabularnewline
\begin{minipage}[t]{0.41\columnwidth}\raggedright
IQ (Installation Qualification)\strut
\end{minipage} & \begin{minipage}[t]{0.53\columnwidth}\raggedright
Equipment installed per specs, utilities connected, calibrations
current\strut
\end{minipage}\tabularnewline
\begin{minipage}[t]{0.41\columnwidth}\raggedright
OQ (Operational Qualification)\strut
\end{minipage} & \begin{minipage}[t]{0.53\columnwidth}\raggedright
Equipment operates correctly across range under no-load conditions\strut
\end{minipage}\tabularnewline
\begin{minipage}[t]{0.41\columnwidth}\raggedright
PQ (Performance Qualification)\strut
\end{minipage} & \begin{minipage}[t]{0.53\columnwidth}\raggedright
Equipment performs reproducibly under production conditions\strut
\end{minipage}\tabularnewline
\bottomrule
\end{longtable}

\emph{Process Performance Qualification (PPQ):} - Typically 3
consecutive successful commercial-scale batches (scientific
justification for number) - Enhanced sampling (more locations, more
tests, more in-process data) - Prospective protocol with pre-defined
acceptance criteria - Statistical analysis: Process capability Ppk ≥
1.33 for critical attributes - All batches must meet all criteria; PPQ
report with deviations documented

\emph{Cleaning Validation:} - Demonstrate removal of product residues,
cleaning agents, and microorganisms - Acceptance criteria: (1) 0.1\%
therapeutic dose carryover, (2) 10 ppm, or (3) visually clean --- most
conservative applied - For potent APIs: Health-Based Exposure Limits
(HBELs) using PDE values - Sampling: Swab (preferred, quantitative) +
rinse (hard-to-reach areas) - 3 consecutive successful cycles;
dirty/clean hold time studies

\textbf{Stage 3 --- Continued Process Verification:} - SPC charts for
CPPs and CQAs (X-bar/R, individuals/moving range) - Ongoing Cpk trending
- Annual Product Quality Review (APQR): batch data, deviations, OOS/OOT,
CAPA, complaints, stability - Revalidation triggers: equipment change,
site change, adverse trends

\hypertarget{sterile-product-validation}{%
\subsection{10.2 Sterile Product
Validation}\label{sterile-product-validation}}

\begin{itemize}
\tightlist
\item
  \textbf{Media fills}: Fill with TSB under production conditions;
  minimum 3 consecutive successful fills; ≥5,000 units; repeat
  semi-annually; acceptance: 0 contaminated units for \textless5,000
  units
\item
  \textbf{Filter validation}: Bacterial challenge (B. diminuta ≥10⁷
  CFU/cm²) for 0.2 µm filters; integrity testing pre- and post-use
\item
  \textbf{Environmental monitoring}: Viable and non-viable particles
  during fills
\item
  \textbf{Hold time studies}: Bioburden/endotoxin limits for bulk
  solution, filtered solution, filled product
\end{itemize}

\begin{center}\rule{0.5\linewidth}{0.5pt}\end{center}

\hypertarget{clinical-trial-phases-as-cmc-triggers}{%
\section{11. Clinical Trial Phases as CMC
Triggers}\label{clinical-trial-phases-as-cmc-triggers}}

This is the critical section: \textbf{how clinical trial data triggers
CMC activities throughout the drug development lifecycle.}

\hypertarget{master-trigger-map}{%
\subsection{11.1 Master Trigger Map}\label{master-trigger-map}}

\begin{verbatim}
┌─────────────────────────────────────────────────────────────────────────────────┐
│                    CLINICAL-TO-CMC TRIGGER TIMELINE                             │
│                                                                                 │
│  DISCOVERY ──► PRE-IND ──► PHASE I ──► PHASE II ──► PHASE III ──► NDA ──► POST │
│                                                                                 │
│  CMC Activities:                                                                │
│  ├── Initial route ─► Process optimization ─► Process lock ─► Validation ─► CPV │
│  ├── Tox batch ──► Clinical supply ──► Scale-up ──► Registration batches ──►    │
│  ├── Pre-form ──► Formulation dev ──► Commercial form ──► Stability ──────►     │
│  ├── Method dev ──► Qualification ──► Validation ──► Transfer ──► Lifecycle     │
│  └── Preliminary specs ──► Refined specs ──► Final specs ──► Justified ────►    │
└─────────────────────────────────────────────────────────────────────────────────┘
\end{verbatim}

\hypertarget{pre-clinical-pre-ind-cmc-activities}{%
\subsection{11.2 Pre-Clinical / Pre-IND CMC
Activities}\label{pre-clinical-pre-ind-cmc-activities}}

\textbf{Regulatory framework:} 21 CFR 312.23(a)(7), FDA guidance ``INDs
for Phase 1 Studies'' (1995).

\textbf{What must be done before any clinical trial begins:}

\begin{longtable}[]{@{}ll@{}}
\toprule
\begin{minipage}[b]{0.32\columnwidth}\raggedright
CMC Area\strut
\end{minipage} & \begin{minipage}[b]{0.62\columnwidth}\raggedright
Pre-IND Requirement\strut
\end{minipage}\tabularnewline
\midrule
\endhead
\begin{minipage}[t]{0.32\columnwidth}\raggedright
Drug Substance\strut
\end{minipage} & \begin{minipage}[t]{0.62\columnwidth}\raggedright
Defined reproducible process (not necessarily optimized), structural
characterization (NMR, MS, IR, UV, XRPD), preliminary impurity profile
qualified against tox batch\strut
\end{minipage}\tabularnewline
\begin{minipage}[t]{0.32\columnwidth}\raggedright
Drug Product\strut
\end{minipage} & \begin{minipage}[t]{0.62\columnwidth}\raggedright
Simple preliminary formulation (powder-in-capsule, simple solution),
adequate for intended route\strut
\end{minipage}\tabularnewline
\begin{minipage}[t]{0.32\columnwidth}\raggedright
Analytical\strut
\end{minipage} & \begin{minipage}[t]{0.62\columnwidth}\raggedright
Fit-for-purpose methods --- suitable but not fully validated; identity,
purity, potency, key safety impurities\strut
\end{minipage}\tabularnewline
\begin{minipage}[t]{0.32\columnwidth}\raggedright
Specifications\strut
\end{minipage} & \begin{minipage}[t]{0.62\columnwidth}\raggedright
Wide, preliminary, justified based on tox batch data\strut
\end{minipage}\tabularnewline
\begin{minipage}[t]{0.32\columnwidth}\raggedright
Stability\strut
\end{minipage} & \begin{minipage}[t]{0.62\columnwidth}\raggedright
1-3 months accelerated + any available real-time data --- enough to
cover planned Phase I duration\strut
\end{minipage}\tabularnewline
\begin{minipage}[t]{0.32\columnwidth}\raggedright
Reference Standard\strut
\end{minipage} & \begin{minipage}[t]{0.62\columnwidth}\raggedright
Characterized sufficiently to serve as reference for
release/stability\strut
\end{minipage}\tabularnewline
\bottomrule
\end{longtable}

\textbf{The Tox Batch Bridge:} Clinical material must be comparable to
(or higher purity than) the GLP toxicology batch. Key comparability:
impurity profiles, potency, physical form.

\textbf{Timeline:} 12-24 months from candidate selection to IND filing.

\hypertarget{phase-i-cmc-triggers}{%
\subsection{11.3 Phase I → CMC Triggers}\label{phase-i-cmc-triggers}}

\textbf{GMP requirements (FDA guidance ``CGMP for Phase 1
Investigational Drugs,'' 2008):} - Risk-proportionate, reduced GMP
vs.~commercial - Process need not be validated; must be defined,
documented, and controlled - Methods need not be fully validated ---
``scientifically sound and suitable'' - Quality unit must be in place;
personnel trained - Batch records required; deviation management, though
less formal

\textbf{Phase I clinical data triggers:}

\begin{longtable}[]{@{}ll@{}}
\toprule
\begin{minipage}[b]{0.39\columnwidth}\raggedright
Clinical Data\strut
\end{minipage} & \begin{minipage}[b]{0.55\columnwidth}\raggedright
CMC Action Triggered\strut
\end{minipage}\tabularnewline
\midrule
\endhead
\begin{minipage}[t]{0.39\columnwidth}\raggedright
\textbf{Unexpected toxicity}\strut
\end{minipage} & \begin{minipage}[t]{0.55\columnwidth}\raggedright
Reformulation (excipient switch, pH adjustment, salt form change),
reassess impurity profile\strut
\end{minipage}\tabularnewline
\begin{minipage}[t]{0.39\columnwidth}\raggedright
\textbf{Injection site reactions}\strut
\end{minipage} & \begin{minipage}[t]{0.55\columnwidth}\raggedright
Modify DP formulation (buffer, tonicity, surfactant)\strut
\end{minipage}\tabularnewline
\begin{minipage}[t]{0.39\columnwidth}\raggedright
\textbf{Poor oral bioavailability}\strut
\end{minipage} & \begin{minipage}[t]{0.55\columnwidth}\raggedright
Enable formulation: particle size reduction, amorphous solid dispersion,
lipid formulation, salt/cocrystal screening\strut
\end{minipage}\tabularnewline
\begin{minipage}[t]{0.39\columnwidth}\raggedright
\textbf{High PK variability}\strut
\end{minipage} & \begin{minipage}[t]{0.55\columnwidth}\raggedright
Evaluate food effects → may trigger enteric coating or
modified-release\strut
\end{minipage}\tabularnewline
\begin{minipage}[t]{0.39\columnwidth}\raggedright
\textbf{Half-life data}\strut
\end{minipage} & \begin{minipage}[t]{0.55\columnwidth}\raggedright
Decision on immediate-release vs.~modified-release\strut
\end{minipage}\tabularnewline
\begin{minipage}[t]{0.39\columnwidth}\raggedright
\textbf{Dose-proportionality}\strut
\end{minipage} & \begin{minipage}[t]{0.55\columnwidth}\raggedright
Determines range of dosage strengths needed\strut
\end{minipage}\tabularnewline
\begin{minipage}[t]{0.39\columnwidth}\raggedright
\textbf{Dose range for Phase II}\strut
\end{minipage} & \begin{minipage}[t]{0.55\columnwidth}\raggedright
Determines batch size requirements and number of strengths\strut
\end{minipage}\tabularnewline
\begin{minipage}[t]{0.39\columnwidth}\raggedright
\textbf{Narrow therapeutic window}\strut
\end{minipage} & \begin{minipage}[t]{0.55\columnwidth}\raggedright
Tighter DP specifications (content uniformity, dissolution)\strut
\end{minipage}\tabularnewline
\bottomrule
\end{longtable}

\textbf{Concurrent CMC activities during Phase I:} - Drug substance
process optimization (yield improvement, problematic reagent removal) -
Salt/polymorph screening and selection - Pre-formulation studies
(excipient compatibility) - Analytical method development and initial
validation - Initiation of long-term stability studies

\hypertarget{phase-ii-cmc-triggers-the-pivotal-investment-decision}{%
\subsection{11.4 Phase II → CMC Triggers (The Pivotal Investment
Decision)}\label{phase-ii-cmc-triggers-the-pivotal-investment-decision}}

\textbf{Phase IIa (Proof of Concept) --- the single most important CMC
trigger:}

Positive proof-of-concept triggers massive CMC investment: - Scale-up
from lab/pilot (1-10 kg) to intermediate/commercial scale (10-100+ kg) -
Technology transfer from development to commercial manufacturing site -
Investment in dedicated manufacturing equipment - Process
characterization studies initiated - Commercial site selection and
qualification begins

\textbf{Phase IIb (Dose Optimization) --- the most critical
clinical-to-CMC trigger:}

\begin{longtable}[]{@{}ll@{}}
\toprule
\begin{minipage}[b]{0.39\columnwidth}\raggedright
Clinical Data\strut
\end{minipage} & \begin{minipage}[b]{0.55\columnwidth}\raggedright
CMC Action Triggered\strut
\end{minipage}\tabularnewline
\midrule
\endhead
\begin{minipage}[t]{0.39\columnwidth}\raggedright
\textbf{Final dose selection}\strut
\end{minipage} & \begin{minipage}[t]{0.55\columnwidth}\raggedright
Determines commercial dosage strength(s) --- arguably the most impactful
single trigger\strut
\end{minipage}\tabularnewline
\begin{minipage}[t]{0.39\columnwidth}\raggedright
\textbf{Dose-response data}\strut
\end{minipage} & \begin{minipage}[t]{0.55\columnwidth}\raggedright
Single vs.~multiple strengths decision\strut
\end{minipage}\tabularnewline
\begin{minipage}[t]{0.39\columnwidth}\raggedright
\textbf{Dosing regimen} (QD, BID, weekly)\strut
\end{minipage} & \begin{minipage}[t]{0.55\columnwidth}\raggedright
Drives formulation strategy (IR, CR, depot)\strut
\end{minipage}\tabularnewline
\begin{minipage}[t]{0.39\columnwidth}\raggedright
\textbf{Route of administration confirmed}\strut
\end{minipage} & \begin{minipage}[t]{0.55\columnwidth}\raggedright
Locks in formulation approach\strut
\end{minipage}\tabularnewline
\bottomrule
\end{longtable}

\textbf{Process development triggered by Phase II:} - Drug substance:
Final route selection, starting material selection, CPP identification -
Scale-up: Often from 1-10 kg to 10-100 kg (new unit operations may be
introduced) - QbD/DoE: Design space mapping per ICH Q8(R2) - PAT
evaluation initiated - Formulation transitions from ``fit-for-purpose''
to ``intended commercial formulation'' - Bioequivalence bridging may be
needed if formulation changed significantly from Phase I

\textbf{Regulatory considerations:} - IND amendments (21 CFR 312.31)
required for significant CMC changes - FDA guidance ``IND Applications
--- Preparing for Phase 2/3 Studies'' (2003) sets expectations: more
complete characterization, updated impurity profiles, preliminary method
validation, expanded stability

\textbf{Timeline:} Phase II CMC activities typically span 2-3 years ---
the most intensive period of CMC development.

\hypertarget{phase-iii-cmc-triggers-commercial-readiness}{%
\subsection{11.5 Phase III → CMC Triggers (Commercial
Readiness)}\label{phase-iii-cmc-triggers-commercial-readiness}}

\textbf{Critical principle:} Sponsors aim to lock the commercial process
BEFORE pivotal Phase III trials begin, because any significant
post-Phase III process change may require bridging studies.

\textbf{Phase III triggers:}

\begin{longtable}[]{@{}lll@{}}
\toprule
\begin{minipage}[b]{0.29\columnwidth}\raggedright
Activity\strut
\end{minipage} & \begin{minipage}[b]{0.26\columnwidth}\raggedright
Trigger\strut
\end{minipage} & \begin{minipage}[b]{0.37\columnwidth}\raggedright
Deliverable\strut
\end{minipage}\tabularnewline
\midrule
\endhead
\begin{minipage}[t]{0.29\columnwidth}\raggedright
Process Design completion (Stage 1)\strut
\end{minipage} & \begin{minipage}[t]{0.26\columnwidth}\raggedright
Phase II data + risk assessment\strut
\end{minipage} & \begin{minipage}[t]{0.37\columnwidth}\raggedright
Design space, CPP/CQA linkage\strut
\end{minipage}\tabularnewline
\begin{minipage}[t]{0.29\columnwidth}\raggedright
Equipment/Facility qualification\strut
\end{minipage} & \begin{minipage}[t]{0.26\columnwidth}\raggedright
Commercial site selected\strut
\end{minipage} & \begin{minipage}[t]{0.37\columnwidth}\raggedright
Full IQ/OQ/PQ\strut
\end{minipage}\tabularnewline
\begin{minipage}[t]{0.29\columnwidth}\raggedright
Process Performance Qualification (PPQ)\strut
\end{minipage} & \begin{minipage}[t]{0.26\columnwidth}\raggedright
Process locked\strut
\end{minipage} & \begin{minipage}[t]{0.37\columnwidth}\raggedright
≥3 consecutive commercial-scale batches\strut
\end{minipage}\tabularnewline
\begin{minipage}[t]{0.29\columnwidth}\raggedright
Registration batches\strut
\end{minipage} & \begin{minipage}[t]{0.26\columnwidth}\raggedright
PPQ complete\strut
\end{minipage} & \begin{minipage}[t]{0.37\columnwidth}\raggedright
Batches for NDA Module 3 + stability\strut
\end{minipage}\tabularnewline
\begin{minipage}[t]{0.29\columnwidth}\raggedright
Full method validation\strut
\end{minipage} & \begin{minipage}[t]{0.26\columnwidth}\raggedright
Commercial methods ready\strut
\end{minipage} & \begin{minipage}[t]{0.37\columnwidth}\raggedright
ICH Q2 validation reports\strut
\end{minipage}\tabularnewline
\begin{minipage}[t]{0.29\columnwidth}\raggedright
Cleaning validation\strut
\end{minipage} & \begin{minipage}[t]{0.26\columnwidth}\raggedright
Multi-product facility\strut
\end{minipage} & \begin{minipage}[t]{0.37\columnwidth}\raggedright
3 consecutive cleaning cycles\strut
\end{minipage}\tabularnewline
\begin{minipage}[t]{0.29\columnwidth}\raggedright
Stability program\strut
\end{minipage} & \begin{minipage}[t]{0.26\columnwidth}\raggedright
Registration batches\strut
\end{minipage} & \begin{minipage}[t]{0.37\columnwidth}\raggedright
ICH Q1A protocol initiated\strut
\end{minipage}\tabularnewline
\begin{minipage}[t]{0.29\columnwidth}\raggedright
Supply chain\strut
\end{minipage} & \begin{minipage}[t]{0.26\columnwidth}\raggedright
Commercial scale confirmed\strut
\end{minipage} & \begin{minipage}[t]{0.37\columnwidth}\raggedright
Multi-source vendor qualification\strut
\end{minipage}\tabularnewline
\begin{minipage}[t]{0.29\columnwidth}\raggedright
Commercial specifications\strut
\end{minipage} & \begin{minipage}[t]{0.26\columnwidth}\raggedright
Batch data + development data\strut
\end{minipage} & \begin{minipage}[t]{0.37\columnwidth}\raggedright
Justified per ICH Q6A/Q6B\strut
\end{minipage}\tabularnewline
\bottomrule
\end{longtable}

\textbf{Registration batch requirements:} - Minimum 1, typically 3
batches at proposed commercial scale (or justified pilot ≥1/10th
commercial) - Final commercial process, final site, commercial-grade
materials, final container closure - Stability initiated per ICH
Q1A(R2): 25°C/60\% RH and 40°C/75\% RH minimum

\textbf{Stability data at NDA filing:} 6 months accelerated + 12 months
long-term (minimum from 3 batches).

\hypertarget{ndabla-submission-cmc-requirements}{%
\subsection{11.6 NDA/BLA Submission --- CMC
Requirements}\label{ndabla-submission-cmc-requirements}}

The complete CMC package in CTD Module 3 includes: - Full process
description with validation data - Full analytical method descriptions
with validation - Batch analyses for clinical, registration, and any
commercial batches - Complete stability data with shelf-life proposal -
Specifications with scientific justification - Container closure system
with E\&L data - Post-approval stability commitment - For biologics:
cell bank characterization, viral safety package, comparability data

\hypertarget{post-approval-cmc-activities-triggered-by-market-data}{%
\subsection{11.7 Post-Approval --- CMC Activities Triggered by Market
Data}\label{post-approval-cmc-activities-triggered-by-market-data}}

\begin{longtable}[]{@{}ll@{}}
\toprule
\begin{minipage}[b]{0.44\columnwidth}\raggedright
Market Data/Event\strut
\end{minipage} & \begin{minipage}[b]{0.50\columnwidth}\raggedright
CMC Activity Triggered\strut
\end{minipage}\tabularnewline
\midrule
\endhead
\begin{minipage}[t]{0.44\columnwidth}\raggedright
Commercial manufacturing experience (CPV data)\strut
\end{minipage} & \begin{minipage}[t]{0.50\columnwidth}\raggedright
SPC trending, process improvements, APQR\strut
\end{minipage}\tabularnewline
\begin{minipage}[t]{0.44\columnwidth}\raggedright
Stability failures/OOS\strut
\end{minipage} & \begin{minipage}[t]{0.50\columnwidth}\raggedright
Investigation, potential recall, reformulation\strut
\end{minipage}\tabularnewline
\begin{minipage}[t]{0.44\columnwidth}\raggedright
Patient feedback (e.g., swallowing difficulty)\strut
\end{minipage} & \begin{minipage}[t]{0.50\columnwidth}\raggedright
New dosage form development (pediatric, ODT)\strut
\end{minipage}\tabularnewline
\begin{minipage}[t]{0.44\columnwidth}\raggedright
Supply chain disruption\strut
\end{minipage} & \begin{minipage}[t]{0.50\columnwidth}\raggedright
Alternate supplier qualification, site transfer\strut
\end{minipage}\tabularnewline
\begin{minipage}[t]{0.44\columnwidth}\raggedright
Cost optimization\strut
\end{minipage} & \begin{minipage}[t]{0.50\columnwidth}\raggedright
Process improvements, yield optimization, site changes\strut
\end{minipage}\tabularnewline
\begin{minipage}[t]{0.44\columnwidth}\raggedright
New indication (different dose/route)\strut
\end{minipage} & \begin{minipage}[t]{0.50\columnwidth}\raggedright
New formulation development, new stability studies\strut
\end{minipage}\tabularnewline
\begin{minipage}[t]{0.44\columnwidth}\raggedright
Patent expiry preparation\strut
\end{minipage} & \begin{minipage}[t]{0.50\columnwidth}\raggedright
Technology transfer to new sites\strut
\end{minipage}\tabularnewline
\begin{minipage}[t]{0.44\columnwidth}\raggedright
Regulatory requirement changes\strut
\end{minipage} & \begin{minipage}[t]{0.50\columnwidth}\raggedright
Nitrosamine assessment, elemental impurity testing\strut
\end{minipage}\tabularnewline
\bottomrule
\end{longtable}

\begin{center}\rule{0.5\linewidth}{0.5pt}\end{center}

\hypertarget{cmc-strategy-and-planning}{%
\section{12. CMC Strategy and
Planning}\label{cmc-strategy-and-planning}}

\hypertarget{cmc-development-plan-phase-aligned}{%
\subsection{12.1 CMC Development Plan ---
Phase-Aligned}\label{cmc-development-plan-phase-aligned}}

\begin{longtable}[]{@{}lll@{}}
\toprule
\begin{minipage}[b]{0.13\columnwidth}\raggedright
Gate\strut
\end{minipage} & \begin{minipage}[b]{0.38\columnwidth}\raggedright
Key CMC Decision\strut
\end{minipage} & \begin{minipage}[b]{0.40\columnwidth}\raggedright
Go/No-Go Criteria\strut
\end{minipage}\tabularnewline
\midrule
\endhead
\begin{minipage}[t]{0.13\columnwidth}\raggedright
Pre-IND\strut
\end{minipage} & \begin{minipage}[t]{0.38\columnwidth}\raggedright
Synthetic route feasible? Formulation viable? GLP material
available?\strut
\end{minipage} & \begin{minipage}[t]{0.40\columnwidth}\raggedright
Process produces quality material; tox batch acceptable\strut
\end{minipage}\tabularnewline
\begin{minipage}[t]{0.13\columnwidth}\raggedright
Phase I Start\strut
\end{minipage} & \begin{minipage}[t]{0.38\columnwidth}\raggedright
GMP material manufactured? IND CMC complete?\strut
\end{minipage} & \begin{minipage}[t]{0.40\columnwidth}\raggedright
Release specs met; stability supports trial duration\strut
\end{minipage}\tabularnewline
\begin{minipage}[t]{0.13\columnwidth}\raggedright
Phase I → Phase II\strut
\end{minipage} & \begin{minipage}[t]{0.38\columnwidth}\raggedright
Process scalable? Formulation suitable?\strut
\end{minipage} & \begin{minipage}[t]{0.40\columnwidth}\raggedright
No showstopper impurities; bioavailability adequate\strut
\end{minipage}\tabularnewline
\begin{minipage}[t]{0.13\columnwidth}\raggedright
Phase II → Phase III\strut
\end{minipage} & \begin{minipage}[t]{0.38\columnwidth}\raggedright
Process locked? Commercial form selected? Methods validation-ready? Site
identified?\strut
\end{minipage} & \begin{minipage}[t]{0.40\columnwidth}\raggedright
Process characterized; design space defined; commercial formulation
stable\strut
\end{minipage}\tabularnewline
\begin{minipage}[t]{0.13\columnwidth}\raggedright
Pre-NDA/BLA\strut
\end{minipage} & \begin{minipage}[t]{0.38\columnwidth}\raggedright
Process validated? Stability supports shelf life? Module 3
complete?\strut
\end{minipage} & \begin{minipage}[t]{0.40\columnwidth}\raggedright
PPQ successful; ≥12 months stability; specs justified\strut
\end{minipage}\tabularnewline
\bottomrule
\end{longtable}

\hypertarget{cmc-investment-and-cost}{%
\subsection{12.2 CMC Investment and
Cost}\label{cmc-investment-and-cost}}

\begin{longtable}[]{@{}lll@{}}
\toprule
Development Phase & Small Molecule CMC Cost & Biologic CMC
Cost\tabularnewline
\midrule
\endhead
IND-enabling (Pre-Phase I) & \$2-5M & \$5-15M\tabularnewline
Phase I-II & \$5-15M & \$15-50M\tabularnewline
Phase III to NDA & \$10-30M & \$50-200M\tabularnewline
Commercial facility & \$50-200M & \$200M-\$2B\tabularnewline
\textbf{Total CMC as \% of development} & \textbf{15-25\%} &
\textbf{30-50\%}\tabularnewline
\bottomrule
\end{longtable}

Cell/gene therapy CMC can be 40-60\% of total development cost.

\hypertarget{parallel-vs.-sequential-activities}{%
\subsection{12.3 Parallel vs.~Sequential
Activities}\label{parallel-vs.-sequential-activities}}

\textbf{Parallel (higher cost, faster):} - Formulation development
simultaneous with API process optimization - Long-term stability while
continuing process characterization - Building commercial suites during
Phase II (high-risk, done for breakthrough therapies) - Qualifying
backup manufacturing sites during pivotal trials

\textbf{Sequential (lower cost, slower):} - Wait for final API process
before formulation development - Complete process validation only after
Phase III feedback - Single-site development → then technology transfer

\textbf{COVID-19 example:} Moderna and Pfizer/BioNTech ran CMC
development almost entirely in parallel with clinical trials ---
manufacturing at-risk at commercial scale during Phase III. Operation
Warp Speed funded this parallel approach, compressing the typical 10+
year timeline to under 1 year.

\hypertarget{platform-approaches}{%
\subsection{12.4 Platform Approaches}\label{platform-approaches}}

For well-established modalities (monoclonal antibodies, ADCs), platform
manufacturing processes and analytical methods can save 1-2 years of
development time: - Standardized cell culture conditions - Standard
purification trains (Protein A → low pH viral inactivation → CEX/AEX →
viral filtration → UF/DF) - Pre-validated analytical method suites -
Pre-qualified container closure systems

\begin{center}\rule{0.5\linewidth}{0.5pt}\end{center}

\hypertarget{regulatory-interactions-for-cmc}{%
\section{13. Regulatory Interactions for
CMC}\label{regulatory-interactions-for-cmc}}

\hypertarget{fda-meeting-types}{%
\subsection{13.1 FDA Meeting Types}\label{fda-meeting-types}}

\begin{longtable}[]{@{}llll@{}}
\toprule
Meeting Type & Purpose & Response Timeline & Meeting
Timeline\tabularnewline
\midrule
\endhead
Type A & CRL resolution, clinical holds & 30 days & 60
days\tabularnewline
Type B & Pre-IND, EOP2, Pre-Phase 3, Pre-NDA/BLA & 60 days & 75
days\tabularnewline
Type C & All other meetings & 60 days & 75 days\tabularnewline
Type D & Written response only & Variable & No meeting\tabularnewline
\bottomrule
\end{longtable}

\hypertarget{key-cmc-topics-by-meeting}{%
\subsection{13.2 Key CMC Topics by
Meeting}\label{key-cmc-topics-by-meeting}}

\textbf{Pre-IND Meeting:} - Synthetic route acceptability, impurity
control strategy - Formulation suitability for FIH - Specifications
adequacy - Stability data requirements - For biologics: cell bank
characterization, viral clearance strategy

\textbf{End-of-Phase 2 Meeting (most important CMC regulatory
interaction):} - Process development status --- ready for registration
batches? - Commercial formulation acceptability; bridging studies
needed? - Commercial specifications justified? - Process validation
approach - Facility readiness - Comparability for process changes (Phase
I/II → Phase III)

\textbf{Pre-NDA/BLA Meeting:} - Module 3 completeness and format -
Outstanding CMC deficiencies - Process validation status - Pre-approval
inspection (PAI) timing - Comparability protocols for anticipated
post-approval changes

\hypertarget{accelerated-pathways-cmc-implications}{%
\subsection{13.3 Accelerated Pathways --- CMC
Implications}\label{accelerated-pathways-cmc-implications}}

\begin{longtable}[]{@{}ll@{}}
\toprule
\begin{minipage}[b]{0.51\columnwidth}\raggedright
Designation\strut
\end{minipage} & \begin{minipage}[b]{0.43\columnwidth}\raggedright
CMC Impact\strut
\end{minipage}\tabularnewline
\midrule
\endhead
\begin{minipage}[t]{0.51\columnwidth}\raggedright
\textbf{Breakthrough Therapy}\strut
\end{minipage} & \begin{minipage}[t]{0.43\columnwidth}\raggedright
More frequent CMC meetings, rolling submission, concurrent validation
may be accepted, less stability data at filing with commitments\strut
\end{minipage}\tabularnewline
\begin{minipage}[t]{0.51\columnwidth}\raggedright
\textbf{Accelerated Approval}\strut
\end{minipage} & \begin{minipage}[t]{0.43\columnwidth}\raggedright
CMC requirements generally not reduced; some flexibility on stability
data\strut
\end{minipage}\tabularnewline
\begin{minipage}[t]{0.51\columnwidth}\raggedright
\textbf{Priority Review}\strut
\end{minipage} & \begin{minipage}[t]{0.43\columnwidth}\raggedright
6-month review clock; PAI must complete within shorter window; CMC
readiness critical at filing\strut
\end{minipage}\tabularnewline
\begin{minipage}[t]{0.51\columnwidth}\raggedright
\textbf{Fast Track}\strut
\end{minipage} & \begin{minipage}[t]{0.43\columnwidth}\raggedright
Rolling review; more CMC interactions\strut
\end{minipage}\tabularnewline
\bottomrule
\end{longtable}

\begin{center}\rule{0.5\linewidth}{0.5pt}\end{center}

\hypertarget{post-approval-cmc-changes-and-lifecycle-management}{%
\section{14. Post-Approval CMC Changes and Lifecycle
Management}\label{post-approval-cmc-changes-and-lifecycle-management}}

\hypertarget{fda-post-approval-change-framework}{%
\subsection{14.1 FDA Post-Approval Change
Framework}\label{fda-post-approval-change-framework}}

\textbf{SUPAC (Scale-Up and Post-Approval Changes) --- Change Levels:}

\begin{longtable}[]{@{}llll@{}}
\toprule
\begin{minipage}[b]{0.17\columnwidth}\raggedright
Level\strut
\end{minipage} & \begin{minipage}[b]{0.20\columnwidth}\raggedright
Impact\strut
\end{minipage} & \begin{minipage}[b]{0.27\columnwidth}\raggedright
Reporting\strut
\end{minipage} & \begin{minipage}[b]{0.25\columnwidth}\raggedright
Examples\strut
\end{minipage}\tabularnewline
\midrule
\endhead
\begin{minipage}[t]{0.17\columnwidth}\raggedright
\textbf{Level 1}\strut
\end{minipage} & \begin{minipage}[t]{0.20\columnwidth}\raggedright
Unlikely to affect quality\strut
\end{minipage} & \begin{minipage}[t]{0.27\columnwidth}\raggedright
\textbf{Annual Report}\strut
\end{minipage} & \begin{minipage}[t]{0.25\columnwidth}\raggedright
≤10× batch size (non-sterile), same-design equipment, editorial batch
record changes\strut
\end{minipage}\tabularnewline
\begin{minipage}[t]{0.17\columnwidth}\raggedright
\textbf{Level 2}\strut
\end{minipage} & \begin{minipage}[t]{0.20\columnwidth}\raggedright
Could have significant impact\strut
\end{minipage} & \begin{minipage}[t]{0.27\columnwidth}\raggedright
\textbf{CBE-30}\strut
\end{minipage} & \begin{minipage}[t]{0.25\columnwidth}\raggedright
Moderate process changes, moderate batch size changes, equipment changes
(different operating principle), same-campus site change\strut
\end{minipage}\tabularnewline
\begin{minipage}[t]{0.17\columnwidth}\raggedright
\textbf{Level 3}\strut
\end{minipage} & \begin{minipage}[t]{0.20\columnwidth}\raggedright
Likely to adversely affect quality\strut
\end{minipage} & \begin{minipage}[t]{0.27\columnwidth}\raggedright
\textbf{PAS (Prior Approval Supplement)}\strut
\end{minipage} & \begin{minipage}[t]{0.25\columnwidth}\raggedright
New site, new synthetic route, new container closure material,
qualitative excipient changes, widening specifications\strut
\end{minipage}\tabularnewline
\bottomrule
\end{longtable}

\hypertarget{supplement-types}{%
\subsection{14.2 Supplement Types}\label{supplement-types}}

\begin{longtable}[]{@{}lll@{}}
\toprule
\begin{minipage}[b]{0.24\columnwidth}\raggedright
Type\strut
\end{minipage} & \begin{minipage}[b]{0.32\columnwidth}\raggedright
Timing\strut
\end{minipage} & \begin{minipage}[b]{0.36\columnwidth}\raggedright
Examples\strut
\end{minipage}\tabularnewline
\midrule
\endhead
\begin{minipage}[t]{0.24\columnwidth}\raggedright
\textbf{PAS}\strut
\end{minipage} & \begin{minipage}[t]{0.32\columnwidth}\raggedright
Cannot implement until FDA approves (4-6 month review)\strut
\end{minipage} & \begin{minipage}[t]{0.36\columnwidth}\raggedright
New manufacturing site, major process change, specification
widening\strut
\end{minipage}\tabularnewline
\begin{minipage}[t]{0.24\columnwidth}\raggedright
\textbf{CBE-30}\strut
\end{minipage} & \begin{minipage}[t]{0.32\columnwidth}\raggedright
Implement 30 days after FDA receipt (unless objection)\strut
\end{minipage} & \begin{minipage}[t]{0.36\columnwidth}\raggedright
Moderate process changes, tightening specs, moderate container
changes\strut
\end{minipage}\tabularnewline
\begin{minipage}[t]{0.24\columnwidth}\raggedright
\textbf{CBE-0}\strut
\end{minipage} & \begin{minipage}[t]{0.32\columnwidth}\raggedright
Implement immediately upon filing\strut
\end{minipage} & \begin{minipage}[t]{0.36\columnwidth}\raggedright
Strengthening warnings, compendial compliance changes\strut
\end{minipage}\tabularnewline
\begin{minipage}[t]{0.24\columnwidth}\raggedright
\textbf{Annual Report}\strut
\end{minipage} & \begin{minipage}[t]{0.32\columnwidth}\raggedright
Within 60 days of NDA anniversary\strut
\end{minipage} & \begin{minipage}[t]{0.36\columnwidth}\raggedright
Minor changes (Level 1), ongoing stability data\strut
\end{minipage}\tabularnewline
\bottomrule
\end{longtable}

\hypertarget{comparability-protocols-21-cfr-314.70a}{%
\subsection{14.3 Comparability Protocols (21 CFR
314.70(a))}\label{comparability-protocols-21-cfr-314.70a}}

Pre-agreed protocols describing: - Specific tests and acceptance
criteria for anticipated changes - If results meet criteria → change
reported at lower category (PAS → CBE-30) - ICH Q12 expands this concept
via PACMPs (Post-Approval Change Management Protocols)

\hypertarget{ich-q12-established-conditions-ecs}{%
\subsection{14.4 ICH Q12 --- Established Conditions
(ECs)}\label{ich-q12-established-conditions-ecs}}

\begin{itemize}
\tightlist
\item
  \textbf{ECs are the legally binding commitments} in the marketing
  authorization
\item
  Only changes to ECs require regulatory notification
\item
  Non-EC parameters managed within the quality system (no filing
  required)
\item
  Promotes risk-based post-approval change management
\item
  Categorizes Module 3 information by regulatory impact
\end{itemize}

\hypertarget{continued-process-verification-stage-3}{%
\subsection{14.5 Continued Process Verification (Stage
3)}\label{continued-process-verification-stage-3}}

\begin{itemize}
\tightlist
\item
  Statistical Process Control (SPC) charts for CPPs and CQAs
\item
  Annual Product Quality Review (APQR): all batches, deviations,
  OOS/OOT, complaints, stability
\item
  Identifies adverse trends early → triggers investigation or
  revalidation
\end{itemize}

\hypertarget{post-approval-commitments-pacs}{%
\subsection{14.6 Post-Approval Commitments
(PACs)}\label{post-approval-commitments-pacs}}

Common CMC post-approval commitments: - Complete process validation (if
concurrent) - Additional stability studies (shelf-life extension) -
Method validation for new methods - Pediatric formulation development -
Enhanced process controls based on CPV data

\begin{center}\rule{0.5\linewidth}{0.5pt}\end{center}

\hypertarget{biologics-specific-cmc}{%
\section{15. Biologics-Specific CMC}\label{biologics-specific-cmc}}

\hypertarget{the-process-is-the-product}{%
\subsection{15.1 ``The Process IS the
Product''}\label{the-process-is-the-product}}

Unlike small molecules (fully characterized by structure), biologics are
defined by their manufacturing process. Even small process changes can
alter the product's: - Glycosylation pattern (affects efficacy,
immunogenicity) - Charge variant profile - Aggregation propensity
(affects safety) - Higher-order structure (affects potency)

\hypertarget{cell-line-and-cell-bank-system}{%
\subsection{15.2 Cell Line and Cell Bank
System}\label{cell-line-and-cell-bank-system}}

\textbf{Development:} Transfection → single-cell cloning → clone
screening → lead clone selection → cell bank generation

\textbf{Cell Bank System:} - \textbf{Master Cell Bank (MCB)}: Generated
from lead clone; extensively characterized; stored in liquid nitrogen;
limited use (generates WCBs) - \textbf{Working Cell Bank (WCB)}:
Generated from MCB; used for production; characterized with confirmatory
testing - \textbf{End-of-Production Cells (EPC)}: Cells harvested at the
limit of in vitro cell age; tested for genetic stability

\textbf{Characterization testing (ICH Q5A, Q5B, Q5D):}

\begin{longtable}[]{@{}ll@{}}
\toprule
\begin{minipage}[b]{0.57\columnwidth}\raggedright
Test Category\strut
\end{minipage} & \begin{minipage}[b]{0.37\columnwidth}\raggedright
Examples\strut
\end{minipage}\tabularnewline
\midrule
\endhead
\begin{minipage}[t]{0.57\columnwidth}\raggedright
Identity\strut
\end{minipage} & \begin{minipage}[t]{0.37\columnwidth}\raggedright
Isoenzyme analysis, DNA fingerprinting, gene sequencing\strut
\end{minipage}\tabularnewline
\begin{minipage}[t]{0.57\columnwidth}\raggedright
Purity\strut
\end{minipage} & \begin{minipage}[t]{0.37\columnwidth}\raggedright
Sterility, mycoplasma, adventitious viruses (in vitro/in vivo),
species-specific viruses\strut
\end{minipage}\tabularnewline
\begin{minipage}[t]{0.57\columnwidth}\raggedright
Genetic stability\strut
\end{minipage} & \begin{minipage}[t]{0.37\columnwidth}\raggedright
Gene copy number, restriction enzyme mapping, sequencing at MCB and
EPC\strut
\end{minipage}\tabularnewline
\begin{minipage}[t]{0.57\columnwidth}\raggedright
Tumorigenicity\strut
\end{minipage} & \begin{minipage}[t]{0.37\columnwidth}\raggedright
Soft agar cloning, nude mouse studies (for some cell lines)\strut
\end{minipage}\tabularnewline
\begin{minipage}[t]{0.57\columnwidth}\raggedright
Retroviral particles\strut
\end{minipage} & \begin{minipage}[t]{0.37\columnwidth}\raggedright
TEM, reverse transcriptase activity, infectivity assays\strut
\end{minipage}\tabularnewline
\bottomrule
\end{longtable}

\hypertarget{viral-safety-ich-q5ar2}{%
\subsection{15.3 Viral Safety (ICH
Q5A(R2))}\label{viral-safety-ich-q5ar2}}

Three complementary approaches: 1. \textbf{Cell line and raw material
testing} for adventitious agents 2. \textbf{Viral clearance validation}:
Demonstrate ≥4 log₁₀ reduction through orthogonal steps (low pH
inactivation, nanofiltration, chromatography) 3. \textbf{Product
testing} at appropriate manufacturing steps

\hypertarget{biologic-characterization}{%
\subsection{15.4 Biologic
Characterization}\label{biologic-characterization}}

\begin{longtable}[]{@{}ll@{}}
\toprule
\begin{minipage}[b]{0.52\columnwidth}\raggedright
Attribute\strut
\end{minipage} & \begin{minipage}[b]{0.42\columnwidth}\raggedright
Methods\strut
\end{minipage}\tabularnewline
\midrule
\endhead
\begin{minipage}[t]{0.52\columnwidth}\raggedright
Primary structure\strut
\end{minipage} & \begin{minipage}[t]{0.42\columnwidth}\raggedright
Peptide mapping (LC-MS/MS), intact mass, amino acid analysis,
N/C-terminal sequencing\strut
\end{minipage}\tabularnewline
\begin{minipage}[t]{0.52\columnwidth}\raggedright
Higher-order structure\strut
\end{minipage} & \begin{minipage}[t]{0.42\columnwidth}\raggedright
CD, DSC, FTIR, HDX-MS, cryo-EM\strut
\end{minipage}\tabularnewline
\begin{minipage}[t]{0.52\columnwidth}\raggedright
Glycosylation\strut
\end{minipage} & \begin{minipage}[t]{0.42\columnwidth}\raggedright
HILIC, CE-LIF, LC-MS, MALDI-TOF\strut
\end{minipage}\tabularnewline
\begin{minipage}[t]{0.52\columnwidth}\raggedright
Charge variants\strut
\end{minipage} & \begin{minipage}[t]{0.42\columnwidth}\raggedright
CEX-HPLC, icIEF, CZE\strut
\end{minipage}\tabularnewline
\begin{minipage}[t]{0.52\columnwidth}\raggedright
Size variants\strut
\end{minipage} & \begin{minipage}[t]{0.42\columnwidth}\raggedright
SEC-HPLC, CE-SDS (reduced/non-reduced), AUC, SEC-MALS\strut
\end{minipage}\tabularnewline
\begin{minipage}[t]{0.52\columnwidth}\raggedright
Potency\strut
\end{minipage} & \begin{minipage}[t]{0.42\columnwidth}\raggedright
Cell-based bioassay, binding assay (SPR, ELISA), ADCC/CDC assays\strut
\end{minipage}\tabularnewline
\begin{minipage}[t]{0.52\columnwidth}\raggedright
Process-related impurities\strut
\end{minipage} & \begin{minipage}[t]{0.42\columnwidth}\raggedright
HCP (ELISA), host cell DNA (qPCR), Protein A (ELISA), residual
media\strut
\end{minipage}\tabularnewline
\bottomrule
\end{longtable}

\hypertarget{comparability-ich-q5e}{%
\subsection{15.5 Comparability (ICH Q5E)}\label{comparability-ich-q5e}}

Required when manufacturing process changes occur during development or
post-approval: - \textbf{Tiered approach}: Analytical (most extensive) →
functional/biological → nonclinical/clinical (if needed) - Analytical
comparability is usually sufficient if comprehensive and well-designed -
Critical for biosimilar development (entire pathway depends on
analytical similarity)

\begin{center}\rule{0.5\linewidth}{0.5pt}\end{center}

\hypertarget{small-molecule-vs.-biologics-cmc-key-differences}{%
\section{16. Small Molecule vs.~Biologics CMC --- Key
Differences}\label{small-molecule-vs.-biologics-cmc-key-differences}}

\begin{longtable}[]{@{}lll@{}}
\toprule
\begin{minipage}[b]{0.22\columnwidth}\raggedright
Aspect\strut
\end{minipage} & \begin{minipage}[b]{0.41\columnwidth}\raggedright
Small Molecule\strut
\end{minipage} & \begin{minipage}[b]{0.28\columnwidth}\raggedright
Biologic\strut
\end{minipage}\tabularnewline
\midrule
\endhead
\begin{minipage}[t]{0.22\columnwidth}\raggedright
Molecular weight\strut
\end{minipage} & \begin{minipage}[t]{0.41\columnwidth}\raggedright
\textless1,000 Da\strut
\end{minipage} & \begin{minipage}[t]{0.28\columnwidth}\raggedright
10,000-150,000+ Da\strut
\end{minipage}\tabularnewline
\begin{minipage}[t]{0.22\columnwidth}\raggedright
Structure\strut
\end{minipage} & \begin{minipage}[t]{0.41\columnwidth}\raggedright
Fully characterized\strut
\end{minipage} & \begin{minipage}[t]{0.28\columnwidth}\raggedright
Complex; defined by process\strut
\end{minipage}\tabularnewline
\begin{minipage}[t]{0.22\columnwidth}\raggedright
Manufacturing\strut
\end{minipage} & \begin{minipage}[t]{0.41\columnwidth}\raggedright
Chemical synthesis\strut
\end{minipage} & \begin{minipage}[t]{0.28\columnwidth}\raggedright
Cell culture/fermentation\strut
\end{minipage}\tabularnewline
\begin{minipage}[t]{0.22\columnwidth}\raggedright
Impurity control\strut
\end{minipage} & \begin{minipage}[t]{0.41\columnwidth}\raggedright
ICH Q3A/B thresholds\strut
\end{minipage} & \begin{minipage}[t]{0.28\columnwidth}\raggedright
Process-specific; HCP, DNA, aggregates\strut
\end{minipage}\tabularnewline
\begin{minipage}[t]{0.22\columnwidth}\raggedright
Analytical methods\strut
\end{minipage} & \begin{minipage}[t]{0.41\columnwidth}\raggedright
HPLC, GC, titration\strut
\end{minipage} & \begin{minipage}[t]{0.28\columnwidth}\raggedright
Bioassays, SEC, CEX, peptide mapping\strut
\end{minipage}\tabularnewline
\begin{minipage}[t]{0.22\columnwidth}\raggedright
Specifications\strut
\end{minipage} & \begin{minipage}[t]{0.41\columnwidth}\raggedright
ICH Q6A\strut
\end{minipage} & \begin{minipage}[t]{0.28\columnwidth}\raggedright
ICH Q6B\strut
\end{minipage}\tabularnewline
\begin{minipage}[t]{0.22\columnwidth}\raggedright
Process changes\strut
\end{minipage} & \begin{minipage}[t]{0.41\columnwidth}\raggedright
Generally lower impact\strut
\end{minipage} & \begin{minipage}[t]{0.28\columnwidth}\raggedright
Can fundamentally alter product\strut
\end{minipage}\tabularnewline
\begin{minipage}[t]{0.22\columnwidth}\raggedright
Stability testing\strut
\end{minipage} & \begin{minipage}[t]{0.41\columnwidth}\raggedright
Standard ICH Q1A\strut
\end{minipage} & \begin{minipage}[t]{0.28\columnwidth}\raggedright
ICH Q5C; potency is primary indicator\strut
\end{minipage}\tabularnewline
\begin{minipage}[t]{0.22\columnwidth}\raggedright
Stability extrapolation\strut
\end{minipage} & \begin{minipage}[t]{0.41\columnwidth}\raggedright
Allowed per Q1E\strut
\end{minipage} & \begin{minipage}[t]{0.28\columnwidth}\raggedright
Generally not accepted for biologics\strut
\end{minipage}\tabularnewline
\begin{minipage}[t]{0.22\columnwidth}\raggedright
Comparability\strut
\end{minipage} & \begin{minipage}[t]{0.41\columnwidth}\raggedright
Rarely needed\strut
\end{minipage} & \begin{minipage}[t]{0.28\columnwidth}\raggedright
Required for any significant change (Q5E)\strut
\end{minipage}\tabularnewline
\begin{minipage}[t]{0.22\columnwidth}\raggedright
GMP standard\strut
\end{minipage} & \begin{minipage}[t]{0.41\columnwidth}\raggedright
21 CFR 210/211, ICH Q7\strut
\end{minipage} & \begin{minipage}[t]{0.28\columnwidth}\raggedright
21 CFR 210/211 + 600-680, ICH Q7\strut
\end{minipage}\tabularnewline
\begin{minipage}[t]{0.22\columnwidth}\raggedright
Typical CMC cost\strut
\end{minipage} & \begin{minipage}[t]{0.41\columnwidth}\raggedright
15-25\% of development\strut
\end{minipage} & \begin{minipage}[t]{0.28\columnwidth}\raggedright
30-50\% of development\strut
\end{minipage}\tabularnewline
\bottomrule
\end{longtable}

\begin{center}\rule{0.5\linewidth}{0.5pt}\end{center}

\hypertarget{advanced-therapies-cell-gene-and-mrna-cmc}{%
\section{17. Advanced Therapies --- Cell, Gene, and mRNA
CMC}\label{advanced-therapies-cell-gene-and-mrna-cmc}}

\hypertarget{cell-therapy-car-t-cmc-challenges}{%
\subsection{17.1 Cell Therapy (CAR-T) CMC
Challenges}\label{cell-therapy-car-t-cmc-challenges}}

\begin{longtable}[]{@{}ll@{}}
\toprule
\begin{minipage}[b]{0.54\columnwidth}\raggedright
Challenge\strut
\end{minipage} & \begin{minipage}[b]{0.40\columnwidth}\raggedright
Detail\strut
\end{minipage}\tabularnewline
\midrule
\endhead
\begin{minipage}[t]{0.54\columnwidth}\raggedright
Autologous manufacturing\strut
\end{minipage} & \begin{minipage}[t]{0.40\columnwidth}\raggedright
Each batch = one patient; no traditional batch statistics\strut
\end{minipage}\tabularnewline
\begin{minipage}[t]{0.54\columnwidth}\raggedright
Starting material variability\strut
\end{minipage} & \begin{minipage}[t]{0.40\columnwidth}\raggedright
Patient cells vary in quality and quantity\strut
\end{minipage}\tabularnewline
\begin{minipage}[t]{0.54\columnwidth}\raggedright
Short shelf life\strut
\end{minipage} & \begin{minipage}[t]{0.40\columnwidth}\raggedright
Some products must be administered within hours\strut
\end{minipage}\tabularnewline
\begin{minipage}[t]{0.54\columnwidth}\raggedright
Chain of identity\strut
\end{minipage} & \begin{minipage}[t]{0.40\columnwidth}\raggedright
Patient-specific material tracking throughout manufacturing\strut
\end{minipage}\tabularnewline
\begin{minipage}[t]{0.54\columnwidth}\raggedright
Potency assays\strut
\end{minipage} & \begin{minipage}[t]{0.40\columnwidth}\raggedright
Functional assays with high variability (cytotoxicity, cytokine
release)\strut
\end{minipage}\tabularnewline
\begin{minipage}[t]{0.54\columnwidth}\raggedright
Sterility testing\strut
\end{minipage} & \begin{minipage}[t]{0.40\columnwidth}\raggedright
14-day test vs.~short shelf life --- release before completion may be
necessary\strut
\end{minipage}\tabularnewline
\begin{minipage}[t]{0.54\columnwidth}\raggedright
Scale-out (not scale-up)\strut
\end{minipage} & \begin{minipage}[t]{0.40\columnwidth}\raggedright
Each batch is small; scale = more parallel batches\strut
\end{minipage}\tabularnewline
\bottomrule
\end{longtable}

\textbf{Example:} Kymriah (tisagenlecleucel, Novartis) --- initial
manufacturing failure rates were high (\textasciitilde8\%);
decentralized model posed quality assurance challenges.

\hypertarget{gene-therapy-aav-vectors-cmc-challenges}{%
\subsection{17.2 Gene Therapy (AAV Vectors) CMC
Challenges}\label{gene-therapy-aav-vectors-cmc-challenges}}

\begin{longtable}[]{@{}ll@{}}
\toprule
\begin{minipage}[b]{0.54\columnwidth}\raggedright
Challenge\strut
\end{minipage} & \begin{minipage}[b]{0.40\columnwidth}\raggedright
Detail\strut
\end{minipage}\tabularnewline
\midrule
\endhead
\begin{minipage}[t]{0.54\columnwidth}\raggedright
Scalability\strut
\end{minipage} & \begin{minipage}[t]{0.40\columnwidth}\raggedright
Transitioning from adherent to suspension cell culture\strut
\end{minipage}\tabularnewline
\begin{minipage}[t]{0.54\columnwidth}\raggedright
Full/empty capsids\strut
\end{minipage} & \begin{minipage}[t]{0.40\columnwidth}\raggedright
Ratio is a CQA; separation by ultracentrifugation or
chromatography\strut
\end{minipage}\tabularnewline
\begin{minipage}[t]{0.54\columnwidth}\raggedright
Vector characterization\strut
\end{minipage} & \begin{minipage}[t]{0.40\columnwidth}\raggedright
Genome integrity, capsid identity, potency (transgene expression)\strut
\end{minipage}\tabularnewline
\begin{minipage}[t]{0.54\columnwidth}\raggedright
Impurity profile\strut
\end{minipage} & \begin{minipage}[t]{0.40\columnwidth}\raggedright
Host cell DNA, HCP, residual plasmid DNA\strut
\end{minipage}\tabularnewline
\begin{minipage}[t]{0.54\columnwidth}\raggedright
Lot-to-lot variability\strut
\end{minipage} & \begin{minipage}[t]{0.40\columnwidth}\raggedright
Greater than traditional biologics\strut
\end{minipage}\tabularnewline
\begin{minipage}[t]{0.54\columnwidth}\raggedright
Manufacturing capacity\strut
\end{minipage} & \begin{minipage}[t]{0.40\columnwidth}\raggedright
Enormous quantities needed; limited global capacity\strut
\end{minipage}\tabularnewline
\bottomrule
\end{longtable}

\textbf{Example:} Zolgensma (onasemnogene abeparvovec, Novartis) ---
manufacturing at scale was extremely challenging; \$2.1M per dose
reflects manufacturing complexity.

\hypertarget{mrna-therapeutics-cmc}{%
\subsection{17.3 mRNA Therapeutics CMC}\label{mrna-therapeutics-cmc}}

\textbf{Drug substance (mRNA):} - In vitro transcription (IVT)
optimization - Purification to remove dsRNA (immune activation trigger)
- CQAs: Capping efficiency, poly(A) tail length, sequence integrity,
truncated species, template DNA removal

\textbf{Drug product (LNP formulation):} - Lipid nanoparticle
composition: ionizable lipid, PEGylated lipid, DSPC, cholesterol - CQAs:
Particle size, polydispersity, encapsulation efficiency, mRNA integrity
within LNP - Cold chain challenges: Original Comirnaty required −70°C;
industry working toward 2-8°C stability through lyophilization

\hypertarget{regulatory-framework-for-advanced-therapies}{%
\subsection{17.4 Regulatory Framework for Advanced
Therapies}\label{regulatory-framework-for-advanced-therapies}}

\begin{itemize}
\tightlist
\item
  FDA: ``CMC Information for Human Gene Therapy INDs'' (2020), potency
  guidance for cell/gene therapy
\item
  EMA: ATMP Regulation (EC) No 1394/2007; Committee for Advanced
  Therapies (CAT)
\item
  ICH: Guidelines for gene therapy and cell therapy vectors under
  development
\end{itemize}

\begin{center}\rule{0.5\linewidth}{0.5pt}\end{center}

\hypertarget{emerging-trends-in-cmc}{%
\section{18. Emerging Trends in CMC}\label{emerging-trends-in-cmc}}

\hypertarget{continuous-manufacturing-ich-q13}{%
\subsection{18.1 Continuous Manufacturing (ICH
Q13)}\label{continuous-manufacturing-ich-q13}}

\begin{longtable}[]{@{}ll@{}}
\toprule
\begin{minipage}[b]{0.47\columnwidth}\raggedright
Aspect\strut
\end{minipage} & \begin{minipage}[b]{0.47\columnwidth}\raggedright
Detail\strut
\end{minipage}\tabularnewline
\midrule
\endhead
\begin{minipage}[t]{0.47\columnwidth}\raggedright
Definition\strut
\end{minipage} & \begin{minipage}[t]{0.47\columnwidth}\raggedright
Material continuously fed and processed; output collected
continuously\strut
\end{minipage}\tabularnewline
\begin{minipage}[t]{0.47\columnwidth}\raggedright
Advantages\strut
\end{minipage} & \begin{minipage}[t]{0.47\columnwidth}\raggedright
80-90\% smaller footprint, hours vs.~weeks, real-time quality, flexible
batch sizes, reduced waste\strut
\end{minipage}\tabularnewline
\begin{minipage}[t]{0.47\columnwidth}\raggedright
Challenges\strut
\end{minipage} & \begin{minipage}[t]{0.47\columnwidth}\raggedright
Batch definition, RTD modeling, cleaning/changeover, capital
investment\strut
\end{minipage}\tabularnewline
\begin{minipage}[t]{0.47\columnwidth}\raggedright
Regulatory\strut
\end{minipage} & \begin{minipage}[t]{0.47\columnwidth}\raggedright
ICH Q13 provides harmonized guidance; FDA strongly supportive\strut
\end{minipage}\tabularnewline
\bottomrule
\end{longtable}

\textbf{Real-world examples:} - Vertex (Orkambi/Symdeko): First
FDA-approved continuous manufacturing (2015) - Janssen (Prezista):
Batch-to-continuous transition approved 2016 - Pfizer: Portable,
continuous, miniature, and modular (PCMM) platform

\hypertarget{real-time-release-testing-rtrt}{%
\subsection{18.2 Real-Time Release Testing
(RTRT)}\label{real-time-release-testing-rtrt}}

\begin{itemize}
\tightlist
\item
  Uses PAT data to release product without traditional end-product
  testing
\item
  Example: NIR model predicts tablet content uniformity during
  compression, replacing HPLC
\item
  Requires robust chemometric model development and validation
\item
  Growing regulatory acceptance
\end{itemize}

\hypertarget{digital-twins}{%
\subsection{18.3 Digital Twins}\label{digital-twins}}

\begin{itemize}
\tightlist
\item
  Digital replica of manufacturing process using real-time data,
  simulation, and ML
\item
  Applications: Virtual experiments, scale-up prediction, real-time
  monitoring, predictive maintenance
\item
  Adopted by large pharma (GSK, Sanofi, Amgen)
\end{itemize}

\hypertarget{aiml-in-cmc}{%
\subsection{18.4 AI/ML in CMC}\label{aiml-in-cmc}}

\begin{longtable}[]{@{}ll@{}}
\toprule
Application & Example\tabularnewline
\midrule
\endhead
Process development & ML models to predict optimal process
parameters\tabularnewline
Process monitoring & AI-driven anomaly detection in real-time
data\tabularnewline
Quality prediction & Predictive CQA models from process
parameters\tabularnewline
Impurity prediction & Computational degradation pathway
prediction\tabularnewline
Document authoring & AI-assisted Module 3 CTD generation\tabularnewline
Batch record review & NLP-based automated review\tabularnewline
\bottomrule
\end{longtable}

\hypertarget{other-trends}{%
\subsection{18.5 Other Trends}\label{other-trends}}

\begin{itemize}
\tightlist
\item
  \textbf{Decentralized/Point-of-Care manufacturing} (cell therapies)
\item
  \textbf{Modular and portable manufacturing} (pandemic preparedness)
\item
  \textbf{3D printing of pharmaceuticals} (Spritam, first FDA-approved
  3D-printed drug, 2015)
\item
  \textbf{FDA Emerging Technology Program} for novel manufacturing
  approaches
\item
  \textbf{FDA KASA} (Knowledge-Aided Assessment and Structured
  Application) for structured CMC review
\end{itemize}

\begin{center}\rule{0.5\linewidth}{0.5pt}\end{center}

\hypertarget{common-cmc-pitfalls-and-how-to-avoid-them}{%
\section{19. Common CMC Pitfalls and How to Avoid
Them}\label{common-cmc-pitfalls-and-how-to-avoid-them}}

\hypertarget{top-causes-of-cmc-related-fda-crls}{%
\subsection{19.1 Top Causes of CMC-Related FDA
CRLs}\label{top-causes-of-cmc-related-fda-crls}}

\begin{longtable}[]{@{}lll@{}}
\toprule
\begin{minipage}[b]{0.26\columnwidth}\raggedright
Rank\strut
\end{minipage} & \begin{minipage}[b]{0.30\columnwidth}\raggedright
Cause\strut
\end{minipage} & \begin{minipage}[b]{0.35\columnwidth}\raggedright
Detail\strut
\end{minipage}\tabularnewline
\midrule
\endhead
\begin{minipage}[t]{0.26\columnwidth}\raggedright
1\strut
\end{minipage} & \begin{minipage}[t]{0.30\columnwidth}\raggedright
\textbf{Manufacturing facility deficiencies}\strut
\end{minipage} & \begin{minipage}[t]{0.35\columnwidth}\raggedright
PAI failures, GMP violations, data integrity issues\strut
\end{minipage}\tabularnewline
\begin{minipage}[t]{0.26\columnwidth}\raggedright
2\strut
\end{minipage} & \begin{minipage}[t]{0.30\columnwidth}\raggedright
\textbf{Insufficient process validation}\strut
\end{minipage} & \begin{minipage}[t]{0.35\columnwidth}\raggedright
Incomplete PPQ, validation not supporting commercial process\strut
\end{minipage}\tabularnewline
\begin{minipage}[t]{0.26\columnwidth}\raggedright
3\strut
\end{minipage} & \begin{minipage}[t]{0.30\columnwidth}\raggedright
\textbf{Stability failures}\strut
\end{minipage} & \begin{minipage}[t]{0.35\columnwidth}\raggedright
OOS results during review, insufficient data for proposed shelf
life\strut
\end{minipage}\tabularnewline
\begin{minipage}[t]{0.26\columnwidth}\raggedright
4\strut
\end{minipage} & \begin{minipage}[t]{0.30\columnwidth}\raggedright
\textbf{Specifications issues}\strut
\end{minipage} & \begin{minipage}[t]{0.35\columnwidth}\raggedright
Unjustified specs, missing critical tests, inadequate impurity
control\strut
\end{minipage}\tabularnewline
\begin{minipage}[t]{0.26\columnwidth}\raggedright
5\strut
\end{minipage} & \begin{minipage}[t]{0.30\columnwidth}\raggedright
\textbf{Analytical method deficiencies}\strut
\end{minipage} & \begin{minipage}[t]{0.35\columnwidth}\raggedright
Unvalidated methods, inadequate potency assay precision
(biologics)\strut
\end{minipage}\tabularnewline
\begin{minipage}[t]{0.26\columnwidth}\raggedright
6\strut
\end{minipage} & \begin{minipage}[t]{0.30\columnwidth}\raggedright
\textbf{Incomplete Module 3}\strut
\end{minipage} & \begin{minipage}[t]{0.35\columnwidth}\raggedright
Missing sections, internal inconsistencies\strut
\end{minipage}\tabularnewline
\bottomrule
\end{longtable}

\hypertarget{cmc-related-clinical-holds}{%
\subsection{19.2 CMC-Related Clinical
Holds}\label{cmc-related-clinical-holds}}

Triggers: - Safety-related impurities (especially genotoxic per ICH M7)
- Sterility or endotoxin failures (injectables) - Stability failures
affecting safety/efficacy - Significant GMP deviations - Inadequate CMC
information in IND - Potency assay issues (biologics --- uncertain
dosing) - Container closure integrity failures (sterile products)

Resolution: 30 days for FDA to respond after sponsor addresses all
deficiencies.

\hypertarget{top-10-best-practices}{%
\subsection{19.3 Top 10 Best Practices}\label{top-10-best-practices}}

\begin{enumerate}
\def\labelenumi{\arabic{enumi}.}
\tightlist
\item
  \textbf{Engage CMC early} --- CMC scientists at the table from
  candidate selection
\item
  \textbf{Plan manufacturing site readiness} --- begin site
  qualification 12-18 months before NDA filing
\item
  \textbf{Start stability studies early} --- cannot accelerate time;
  lock commercial formulation early
\item
  \textbf{De-risk supply chain} --- qualify multiple suppliers for
  critical materials
\item
  \textbf{Invest in analytical development} --- under-resourcing is a
  common cause of delays
\item
  \textbf{Robust change management} --- track all process changes and
  their impact
\item
  \textbf{Use platform approaches} where possible --- saves 1-2 years
  for established modalities
\item
  \textbf{Conduct mock PAI audits} --- before filing, identify and
  remediate gaps
\item
  \textbf{Prepare for information requests} --- have CMC team ready for
  rapid FDA response
\item
  \textbf{Align globally} --- if multi-region filing, harmonize CMC
  strategy across FDA/EMA/PMDA from the start
\end{enumerate}

\hypertarget{real-world-cautionary-examples}{%
\subsection{19.4 Real-World Cautionary
Examples}\label{real-world-cautionary-examples}}

\begin{itemize}
\tightlist
\item
  \textbf{Dendreon (Provenge):} Manufacturing delays and cost overruns
  for autologous cell therapy → bankruptcy
\item
  \textbf{Emergent BioSolutions (COVID):} Cross-contamination between
  AstraZeneca and J\&J vaccine substances → millions of doses destroyed,
  FDA facility shutdown
\item
  \textbf{Sarepta (Elevidys):} Gene therapy manufacturing consistency
  and potency assay variability → initial CMC challenges
\item
  \textbf{Multiple gene therapy programs:} Delays from insufficient
  viral vector manufacturing capacity
\end{itemize}

\begin{center}\rule{0.5\linewidth}{0.5pt}\end{center}

\hypertarget{fda-vs.-ema-cmc-comparison}{%
\section{20. FDA vs.~EMA --- CMC
Comparison}\label{fda-vs.-ema-cmc-comparison}}

\begin{longtable}[]{@{}lll@{}}
\toprule
\begin{minipage}[b]{0.26\columnwidth}\raggedright
Aspect\strut
\end{minipage} & \begin{minipage}[b]{0.33\columnwidth}\raggedright
FDA (US)\strut
\end{minipage} & \begin{minipage}[b]{0.33\columnwidth}\raggedright
EMA (EU)\strut
\end{minipage}\tabularnewline
\midrule
\endhead
\begin{minipage}[t]{0.26\columnwidth}\raggedright
Filing format\strut
\end{minipage} & \begin{minipage}[t]{0.33\columnwidth}\raggedright
CTD Module 3 (identical)\strut
\end{minipage} & \begin{minipage}[t]{0.33\columnwidth}\raggedright
CTD Module 3 (identical)\strut
\end{minipage}\tabularnewline
\begin{minipage}[t]{0.26\columnwidth}\raggedright
API documentation\strut
\end{minipage} & \begin{minipage}[t]{0.33\columnwidth}\raggedright
DMF (Drug Master File)\strut
\end{minipage} & \begin{minipage}[t]{0.33\columnwidth}\raggedright
ASMF or CEP (Certificate of Suitability from EDQM)\strut
\end{minipage}\tabularnewline
\begin{minipage}[t]{0.26\columnwidth}\raggedright
Post-approval changes\strut
\end{minipage} & \begin{minipage}[t]{0.33\columnwidth}\raggedright
SUPAC framework (PAS, CBE-30, CBE-0, Annual Report)\strut
\end{minipage} & \begin{minipage}[t]{0.33\columnwidth}\raggedright
Variations Regulation (Type IA, IAIN, IB, Type II, Extensions)\strut
\end{minipage}\tabularnewline
\begin{minipage}[t]{0.26\columnwidth}\raggedright
GMP inspections\strut
\end{minipage} & \begin{minipage}[t]{0.33\columnwidth}\raggedright
FDA ORA directly inspects\strut
\end{minipage} & \begin{minipage}[t]{0.33\columnwidth}\raggedright
National competent authorities (MHRA, ANSM, etc.)\strut
\end{minipage}\tabularnewline
\begin{minipage}[t]{0.26\columnwidth}\raggedright
Batch release\strut
\end{minipage} & \begin{minipage}[t]{0.33\columnwidth}\raggedright
QC release by manufacturer\strut
\end{minipage} & \begin{minipage}[t]{0.33\columnwidth}\raggedright
Qualified Person (QP) certification legally required\strut
\end{minipage}\tabularnewline
\begin{minipage}[t]{0.26\columnwidth}\raggedright
Compendial standard\strut
\end{minipage} & \begin{minipage}[t]{0.33\columnwidth}\raggedright
USP-NF\strut
\end{minipage} & \begin{minipage}[t]{0.33\columnwidth}\raggedright
European Pharmacopoeia (Ph. Eur.)\strut
\end{minipage}\tabularnewline
\begin{minipage}[t]{0.26\columnwidth}\raggedright
Sterile manufacturing\strut
\end{minipage} & \begin{minipage}[t]{0.33\columnwidth}\raggedright
FDA Aseptic Processing Guidance (2004)\strut
\end{minipage} & \begin{minipage}[t]{0.33\columnwidth}\raggedright
EU GMP Annex 1 (2022 --- significantly more detailed)\strut
\end{minipage}\tabularnewline
\begin{minipage}[t]{0.26\columnwidth}\raggedright
Assessment transparency\strut
\end{minipage} & \begin{minipage}[t]{0.33\columnwidth}\raggedright
FDA review documents\strut
\end{minipage} & \begin{minipage}[t]{0.33\columnwidth}\raggedright
EPAR with detailed quality assessment\strut
\end{minipage}\tabularnewline
\begin{minipage}[t]{0.26\columnwidth}\raggedright
Assessment process\strut
\end{minipage} & \begin{minipage}[t]{0.33\columnwidth}\raggedright
Single FDA division review\strut
\end{minipage} & \begin{minipage}[t]{0.33\columnwidth}\raggedright
Rapporteur + Co-Rapporteur from 2 member states\strut
\end{minipage}\tabularnewline
\begin{minipage}[t]{0.26\columnwidth}\raggedright
Biosimilar CMC\strut
\end{minipage} & \begin{minipage}[t]{0.33\columnwidth}\raggedright
FDA framework (2012+)\strut
\end{minipage} & \begin{minipage}[t]{0.33\columnwidth}\raggedright
EMA pathway (established earlier, more extensive guidelines)\strut
\end{minipage}\tabularnewline
\begin{minipage}[t]{0.26\columnwidth}\raggedright
ATMPs\strut
\end{minipage} & \begin{minipage}[t]{0.33\columnwidth}\raggedright
FDA gene/cell therapy guidances\strut
\end{minipage} & \begin{minipage}[t]{0.33\columnwidth}\raggedright
Regulation (EC) 1394/2007, CAT committee\strut
\end{minipage}\tabularnewline
\begin{minipage}[t]{0.26\columnwidth}\raggedright
TSE/BSE risk\strut
\end{minipage} & \begin{minipage}[t]{0.33\columnwidth}\raggedright
Adventitious agent assessment\strut
\end{minipage} & \begin{minipage}[t]{0.33\columnwidth}\raggedright
Specific formalized TSE guidance (EMA/410/01 Rev.3)\strut
\end{minipage}\tabularnewline
\bottomrule
\end{longtable}

\textbf{EMA Variation Types:}

\begin{longtable}[]{@{}lll@{}}
\toprule
\begin{minipage}[b]{0.21\columnwidth}\raggedright
Type\strut
\end{minipage} & \begin{minipage}[b]{0.42\columnwidth}\raggedright
Description\strut
\end{minipage} & \begin{minipage}[b]{0.28\columnwidth}\raggedright
Timing\strut
\end{minipage}\tabularnewline
\midrule
\endhead
\begin{minipage}[t]{0.21\columnwidth}\raggedright
Type IA\strut
\end{minipage} & \begin{minipage}[t]{0.42\columnwidth}\raggedright
Minor, ``do and tell''\strut
\end{minipage} & \begin{minipage}[t]{0.28\columnwidth}\raggedright
Implement immediately, notify within 12 months\strut
\end{minipage}\tabularnewline
\begin{minipage}[t]{0.21\columnwidth}\raggedright
Type IAIN\strut
\end{minipage} & \begin{minipage}[t]{0.42\columnwidth}\raggedright
Minor, immediate notification\strut
\end{minipage} & \begin{minipage}[t]{0.28\columnwidth}\raggedright
Implement and notify immediately\strut
\end{minipage}\tabularnewline
\begin{minipage}[t]{0.21\columnwidth}\raggedright
Type IB\strut
\end{minipage} & \begin{minipage}[t]{0.42\columnwidth}\raggedright
Minor, ``tell, wait, and do''\strut
\end{minipage} & \begin{minipage}[t]{0.28\columnwidth}\raggedright
Notify; implement after 30 days if no objection\strut
\end{minipage}\tabularnewline
\begin{minipage}[t]{0.21\columnwidth}\raggedright
Type II\strut
\end{minipage} & \begin{minipage}[t]{0.42\columnwidth}\raggedright
Major\strut
\end{minipage} & \begin{minipage}[t]{0.28\columnwidth}\raggedright
Requires prior approval (60-day review)\strut
\end{minipage}\tabularnewline
\begin{minipage}[t]{0.21\columnwidth}\raggedright
Extension\strut
\end{minipage} & \begin{minipage}[t]{0.42\columnwidth}\raggedright
Fundamental change\strut
\end{minipage} & \begin{minipage}[t]{0.28\columnwidth}\raggedright
New application-level review\strut
\end{minipage}\tabularnewline
\bottomrule
\end{longtable}

\begin{center}\rule{0.5\linewidth}{0.5pt}\end{center}

\hypertarget{technology-transfer}{%
\section{21. Technology Transfer}\label{technology-transfer}}

\hypertarget{what-technology-transfer-means-in-cmc}{%
\subsection{21.1 What Technology Transfer Means in
CMC}\label{what-technology-transfer-means-in-cmc}}

Technology transfer is the systematic, documented process of
transferring process knowledge, analytical methods, and manufacturing
procedures from one site (the ``sending unit'') to another (the
``receiving unit''). It is one of the most practically important and
frequently encountered CMC activities.

\textbf{Common transfer scenarios:} - Development lab → pilot plant -
Pilot plant → commercial manufacturing site - CDMO → in-house
manufacturing (or reverse) - Site A → Site B (commercial to commercial,
for capacity, cost, or risk mitigation) - Originator → licensee
(in-licensing/out-licensing deals)

\hypertarget{technology-transfer-package}{%
\subsection{21.2 Technology Transfer
Package}\label{technology-transfer-package}}

A complete technology transfer includes:

\begin{longtable}[]{@{}ll@{}}
\toprule
\begin{minipage}[b]{0.49\columnwidth}\raggedright
Component\strut
\end{minipage} & \begin{minipage}[b]{0.45\columnwidth}\raggedright
Contents\strut
\end{minipage}\tabularnewline
\midrule
\endhead
\begin{minipage}[t]{0.49\columnwidth}\raggedright
\textbf{Process description}\strut
\end{minipage} & \begin{minipage}[t]{0.45\columnwidth}\raggedright
Detailed process flow, unit operations, CPPs, CQAs, design space\strut
\end{minipage}\tabularnewline
\begin{minipage}[t]{0.49\columnwidth}\raggedright
\textbf{Master batch record}\strut
\end{minipage} & \begin{minipage}[t]{0.45\columnwidth}\raggedright
Step-by-step manufacturing instructions, IPCs, sampling plans\strut
\end{minipage}\tabularnewline
\begin{minipage}[t]{0.49\columnwidth}\raggedright
\textbf{Analytical methods}\strut
\end{minipage} & \begin{minipage}[t]{0.45\columnwidth}\raggedright
Method descriptions, validation reports, reference standards\strut
\end{minipage}\tabularnewline
\begin{minipage}[t]{0.49\columnwidth}\raggedright
\textbf{Raw material specs}\strut
\end{minipage} & \begin{minipage}[t]{0.45\columnwidth}\raggedright
Specifications for all starting materials, intermediates,
excipients\strut
\end{minipage}\tabularnewline
\begin{minipage}[t]{0.49\columnwidth}\raggedright
\textbf{Equipment requirements}\strut
\end{minipage} & \begin{minipage}[t]{0.45\columnwidth}\raggedright
Equipment specifications, capacity, materials of construction, scale-up
rationale\strut
\end{minipage}\tabularnewline
\begin{minipage}[t]{0.49\columnwidth}\raggedright
\textbf{Cleaning procedures}\strut
\end{minipage} & \begin{minipage}[t]{0.45\columnwidth}\raggedright
Validated cleaning methods, acceptance criteria, MACO/HBEL
calculations\strut
\end{minipage}\tabularnewline
\begin{minipage}[t]{0.49\columnwidth}\raggedright
\textbf{Stability data}\strut
\end{minipage} & \begin{minipage}[t]{0.45\columnwidth}\raggedright
Historical stability data, protocol for new-site batches\strut
\end{minipage}\tabularnewline
\begin{minipage}[t]{0.49\columnwidth}\raggedright
\textbf{Regulatory history}\strut
\end{minipage} & \begin{minipage}[t]{0.45\columnwidth}\raggedright
IND/NDA amendments, process change history, prior regulatory
commitments\strut
\end{minipage}\tabularnewline
\begin{minipage}[t]{0.49\columnwidth}\raggedright
\textbf{Risk assessment}\strut
\end{minipage} & \begin{minipage}[t]{0.45\columnwidth}\raggedright
Site-specific risk assessment, equipment mapping, gap analysis\strut
\end{minipage}\tabularnewline
\bottomrule
\end{longtable}

\hypertarget{transfer-execution}{%
\subsection{21.3 Transfer Execution}\label{transfer-execution}}

\textbf{Transfer batches:} - Typically 1-3 batches at receiving site
(prior to formal PPQ) - Comparison to sending site batch data on all
CQAs - Acceptance criteria defined in transfer protocol (e.g., CQAs
within ±X\% of sending site mean) - Enhanced sampling (similar to PPQ)

\textbf{Equipment mapping:} - Different equipment brands/models at
sending vs.~receiving site require parameter translation - Dimensional
analysis, geometric similarity, scale-up parameters (tip speed, Froude
number, Reynolds number) - Example: Different granulator → match
impeller tip speed, not RPM

\textbf{Common transfer failures and root causes:} - Dissolution shift
(different granulator design, compression profile) - Yield drop
(material handling losses at different scale) - Impurity profile changes
(different heating/cooling rates, hold times) - Content uniformity
issues (different blender geometry)

\hypertarget{analytical-method-transfer}{%
\subsection{21.4 Analytical Method
Transfer}\label{analytical-method-transfer}}

\textbf{Approaches (per USP \textless1224\textgreater):}

\begin{longtable}[]{@{}ll@{}}
\toprule
\begin{minipage}[b]{0.45\columnwidth}\raggedright
Approach\strut
\end{minipage} & \begin{minipage}[b]{0.49\columnwidth}\raggedright
When Used\strut
\end{minipage}\tabularnewline
\midrule
\endhead
\begin{minipage}[t]{0.45\columnwidth}\raggedright
Comparative testing\strut
\end{minipage} & \begin{minipage}[t]{0.49\columnwidth}\raggedright
Most common; both labs analyze same samples; statistical
acceptance\strut
\end{minipage}\tabularnewline
\begin{minipage}[t]{0.45\columnwidth}\raggedright
Co-validation\strut
\end{minipage} & \begin{minipage}[t]{0.49\columnwidth}\raggedright
Both labs validate independently with same protocol\strut
\end{minipage}\tabularnewline
\begin{minipage}[t]{0.45\columnwidth}\raggedright
Revalidation\strut
\end{minipage} & \begin{minipage}[t]{0.49\columnwidth}\raggedright
Receiving lab performs full validation\strut
\end{minipage}\tabularnewline
\begin{minipage}[t]{0.45\columnwidth}\raggedright
Transfer waiver\strut
\end{minipage} & \begin{minipage}[t]{0.49\columnwidth}\raggedright
Simple methods (pH, KF) or compendial methods\strut
\end{minipage}\tabularnewline
\bottomrule
\end{longtable}

\textbf{Acceptance criteria:} Receiving lab typically within ±2-3\% of
sending lab mean, equivalent precision (F-test for variances, t-test for
means).

\hypertarget{regulatory-implications}{%
\subsection{21.5 Regulatory
Implications}\label{regulatory-implications}}

\begin{itemize}
\tightlist
\item
  \textbf{New commercial manufacturing site} = PAS (Prior Approval
  Supplement) in US; Type II Variation in EU
\item
  PAI at the receiving site typically required
\item
  New stability data from receiving site batches required
\item
  If transfer occurs during development, IND amendment required for
  significant site changes
\end{itemize}

\begin{center}\rule{0.5\linewidth}{0.5pt}\end{center}

\hypertarget{cmocdmo-management-and-quality-agreements}{%
\section{22. CMO/CDMO Management and Quality
Agreements}\label{cmocdmo-management-and-quality-agreements}}

\hypertarget{why-cdmo-management-matters}{%
\subsection{22.1 Why CDMO Management
Matters}\label{why-cdmo-management-matters}}

The majority of small and mid-size pharma companies use CDMOs (Contract
Development and Manufacturing Organizations) for some or all
manufacturing. The marketing authorization holder retains full
regulatory responsibility regardless of outsourcing.

\hypertarget{cdmo-selection-criteria}{%
\subsection{22.2 CDMO Selection
Criteria}\label{cdmo-selection-criteria}}

\begin{longtable}[]{@{}ll@{}}
\toprule
\begin{minipage}[b]{0.47\columnwidth}\raggedright
Criterion\strut
\end{minipage} & \begin{minipage}[b]{0.47\columnwidth}\raggedright
Assessment\strut
\end{minipage}\tabularnewline
\midrule
\endhead
\begin{minipage}[t]{0.47\columnwidth}\raggedright
Technical capability\strut
\end{minipage} & \begin{minipage}[t]{0.47\columnwidth}\raggedright
Equipment, scale, dosage form expertise, modality experience\strut
\end{minipage}\tabularnewline
\begin{minipage}[t]{0.47\columnwidth}\raggedright
Regulatory track record\strut
\end{minipage} & \begin{minipage}[t]{0.47\columnwidth}\raggedright
FDA/EMA inspection history, 483s, Warning Letters\strut
\end{minipage}\tabularnewline
\begin{minipage}[t]{0.47\columnwidth}\raggedright
Capacity and availability\strut
\end{minipage} & \begin{minipage}[t]{0.47\columnwidth}\raggedright
Current utilization, ability to meet timeline, dedicated vs.~shared
suites\strut
\end{minipage}\tabularnewline
\begin{minipage}[t]{0.47\columnwidth}\raggedright
Quality systems maturity\strut
\end{minipage} & \begin{minipage}[t]{0.47\columnwidth}\raggedright
QMS, deviation rate, CAPA effectiveness, change control\strut
\end{minipage}\tabularnewline
\begin{minipage}[t]{0.47\columnwidth}\raggedright
Geographic location\strut
\end{minipage} & \begin{minipage}[t]{0.47\columnwidth}\raggedright
Proximity to sponsor, regulatory risk (Import Alerts), logistics\strut
\end{minipage}\tabularnewline
\begin{minipage}[t]{0.47\columnwidth}\raggedright
Financial stability\strut
\end{minipage} & \begin{minipage}[t]{0.47\columnwidth}\raggedright
Solvency, long-term viability\strut
\end{minipage}\tabularnewline
\begin{minipage}[t]{0.47\columnwidth}\raggedright
IP protection\strut
\end{minipage} & \begin{minipage}[t]{0.47\columnwidth}\raggedright
Confidentiality practices, data security, non-compete clauses\strut
\end{minipage}\tabularnewline
\begin{minipage}[t]{0.47\columnwidth}\raggedright
References\strut
\end{minipage} & \begin{minipage}[t]{0.47\columnwidth}\raggedright
Track record with similar molecules/products\strut
\end{minipage}\tabularnewline
\bottomrule
\end{longtable}

\hypertarget{quality-agreements}{%
\subsection{22.3 Quality Agreements}\label{quality-agreements}}

\textbf{Required by:} 21 CFR 211.22 (US), EU GMP Chapter 7, ICH Q7
Section 16

\textbf{Essential content:}

\begin{longtable}[]{@{}ll@{}}
\toprule
\begin{minipage}[b]{0.38\columnwidth}\raggedright
Section\strut
\end{minipage} & \begin{minipage}[b]{0.56\columnwidth}\raggedright
Key Elements\strut
\end{minipage}\tabularnewline
\midrule
\endhead
\begin{minipage}[t]{0.38\columnwidth}\raggedright
\textbf{Roles and responsibilities}\strut
\end{minipage} & \begin{minipage}[t]{0.56\columnwidth}\raggedright
Who approves batch records, who investigates deviations, who releases
batches\strut
\end{minipage}\tabularnewline
\begin{minipage}[t]{0.38\columnwidth}\raggedright
\textbf{Change control}\strut
\end{minipage} & \begin{minipage}[t]{0.56\columnwidth}\raggedright
Mutual notification requirements, approval process for changes affecting
product quality\strut
\end{minipage}\tabularnewline
\begin{minipage}[t]{0.38\columnwidth}\raggedright
\textbf{Deviation management}\strut
\end{minipage} & \begin{minipage}[t]{0.56\columnwidth}\raggedright
Notification timelines, investigation responsibilities, CAPA
requirements\strut
\end{minipage}\tabularnewline
\begin{minipage}[t]{0.38\columnwidth}\raggedright
\textbf{Batch release}\strut
\end{minipage} & \begin{minipage}[t]{0.56\columnwidth}\raggedright
Release criteria, documentation required, QP responsibilities (EU)\strut
\end{minipage}\tabularnewline
\begin{minipage}[t]{0.38\columnwidth}\raggedright
\textbf{Audit rights}\strut
\end{minipage} & \begin{minipage}[t]{0.56\columnwidth}\raggedright
Frequency, scope, notice period, regulatory inspection support\strut
\end{minipage}\tabularnewline
\begin{minipage}[t]{0.38\columnwidth}\raggedright
\textbf{Regulatory obligations}\strut
\end{minipage} & \begin{minipage}[t]{0.56\columnwidth}\raggedright
Filing responsibilities, inspection support, annual report
contributions\strut
\end{minipage}\tabularnewline
\begin{minipage}[t]{0.38\columnwidth}\raggedright
\textbf{Specifications and testing}\strut
\end{minipage} & \begin{minipage}[t]{0.56\columnwidth}\raggedright
Who performs release testing, who holds reference standards\strut
\end{minipage}\tabularnewline
\begin{minipage}[t]{0.38\columnwidth}\raggedright
\textbf{Stability}\strut
\end{minipage} & \begin{minipage}[t]{0.56\columnwidth}\raggedright
Who conducts stability studies, data sharing\strut
\end{minipage}\tabularnewline
\begin{minipage}[t]{0.38\columnwidth}\raggedright
\textbf{Complaints and recalls}\strut
\end{minipage} & \begin{minipage}[t]{0.56\columnwidth}\raggedright
Investigation responsibilities, communication protocols\strut
\end{minipage}\tabularnewline
\begin{minipage}[t]{0.38\columnwidth}\raggedright
\textbf{Material management}\strut
\end{minipage} & \begin{minipage}[t]{0.56\columnwidth}\raggedright
Raw material sourcing, approved suppliers, incoming testing\strut
\end{minipage}\tabularnewline
\begin{minipage}[t]{0.38\columnwidth}\raggedright
\textbf{Subcontracting}\strut
\end{minipage} & \begin{minipage}[t]{0.56\columnwidth}\raggedright
Approval requirements for further outsourcing\strut
\end{minipage}\tabularnewline
\bottomrule
\end{longtable}

\hypertarget{sponsor-oversight-of-cdmos}{%
\subsection{22.4 Sponsor Oversight of
CDMOs}\label{sponsor-oversight-of-cdmos}}

\begin{itemize}
\tightlist
\item
  \textbf{Regular audits}: Initial qualification audit, annual/biannual
  surveillance audits, for-cause audits
\item
  \textbf{Batch record review}: Real-time or retrospective review of all
  batch production records
\item
  \textbf{Deviation trending}: Monitor CDMO deviation rates and types
  over time
\item
  \textbf{KPI tracking}: Right-first-time rate, batch cycle time,
  deviation closure time, OOS rate
\item
  \textbf{On-site presence}: Embedded quality representative for
  critical programs
\item
  \textbf{Remote monitoring}: Digital batch record systems enabling
  real-time visibility
\end{itemize}

\hypertarget{multi-cdmo-strategy}{%
\subsection{22.5 Multi-CDMO Strategy}\label{multi-cdmo-strategy}}

\begin{itemize}
\tightlist
\item
  Dual-source critical manufacturing steps for supply chain resilience
\item
  Regulatory implications: each site requires its own
  supplement/variation
\item
  Tech transfer between CDMOs requires full transfer package
\item
  Comparability data between sites required
\end{itemize}

\begin{center}\rule{0.5\linewidth}{0.5pt}\end{center}

\hypertarget{supply-chain-and-vendor-qualification}{%
\section{23. Supply Chain and Vendor
Qualification}\label{supply-chain-and-vendor-qualification}}

\hypertarget{vendor-qualification-program}{%
\subsection{23.1 Vendor Qualification
Program}\label{vendor-qualification-program}}

\textbf{Risk-based tiering:}

\begin{longtable}[]{@{}lll@{}}
\toprule
\begin{minipage}[b]{0.14\columnwidth}\raggedright
Tier\strut
\end{minipage} & \begin{minipage}[b]{0.33\columnwidth}\raggedright
Material Type\strut
\end{minipage} & \begin{minipage}[b]{0.44\columnwidth}\raggedright
Qualification Level\strut
\end{minipage}\tabularnewline
\midrule
\endhead
\begin{minipage}[t]{0.14\columnwidth}\raggedright
Critical\strut
\end{minipage} & \begin{minipage}[t]{0.33\columnwidth}\raggedright
API starting materials, key intermediates, critical excipients\strut
\end{minipage} & \begin{minipage}[t]{0.44\columnwidth}\raggedright
Full on-site audit, qualification testing, annual review\strut
\end{minipage}\tabularnewline
\begin{minipage}[t]{0.14\columnwidth}\raggedright
Major\strut
\end{minipage} & \begin{minipage}[t]{0.33\columnwidth}\raggedright
Excipients, processing aids, primary packaging\strut
\end{minipage} & \begin{minipage}[t]{0.44\columnwidth}\raggedright
Desktop audit + questionnaire, CoA verification + periodic testing\strut
\end{minipage}\tabularnewline
\begin{minipage}[t]{0.14\columnwidth}\raggedright
Minor\strut
\end{minipage} & \begin{minipage}[t]{0.33\columnwidth}\raggedright
Secondary packaging, non-product contact materials\strut
\end{minipage} & \begin{minipage}[t]{0.44\columnwidth}\raggedright
Questionnaire, CoA review\strut
\end{minipage}\tabularnewline
\bottomrule
\end{longtable}

\hypertarget{supplier-qualification-activities}{%
\subsection{23.2 Supplier Qualification
Activities}\label{supplier-qualification-activities}}

\begin{itemize}
\tightlist
\item
  \textbf{Audit program}: Initial qualification audit, routine
  surveillance (every 2-3 years for critical), for-cause audits
\item
  \textbf{Qualification testing}: Compare supplier CoA results to
  in-house testing (minimum 3 lots)
\item
  \textbf{Change notification agreements}: Supplier must notify of any
  changes to manufacturing process, specifications, or site
\item
  \textbf{Approved supplier list}: Maintained by quality unit; changes
  require formal change control
\end{itemize}

\hypertarget{incoming-material-testing}{%
\subsection{23.3 Incoming Material
Testing}\label{incoming-material-testing}}

\begin{itemize}
\tightlist
\item
  \textbf{Full testing}: Per specification for each lot received
  (required for APIs)
\item
  \textbf{Reduced testing}: Identity + CoA verification after
  qualification (permitted for excipients under risk-based
  justification, per 21 CFR 211.84)
\item
  \textbf{Skip-lot testing}: Testing every Nth lot after extensive
  qualification and trending (requires scientific justification)
\end{itemize}

\hypertarget{supply-chain-risk-management}{%
\subsection{23.4 Supply Chain Risk
Management}\label{supply-chain-risk-management}}

\begin{longtable}[]{@{}ll@{}}
\toprule
\begin{minipage}[b]{0.33\columnwidth}\raggedright
Risk\strut
\end{minipage} & \begin{minipage}[b]{0.61\columnwidth}\raggedright
Mitigation\strut
\end{minipage}\tabularnewline
\midrule
\endhead
\begin{minipage}[t]{0.33\columnwidth}\raggedright
Single-source dependency\strut
\end{minipage} & \begin{minipage}[t]{0.61\columnwidth}\raggedright
Qualify ≥2 suppliers for critical materials\strut
\end{minipage}\tabularnewline
\begin{minipage}[t]{0.33\columnwidth}\raggedright
Geographic concentration\strut
\end{minipage} & \begin{minipage}[t]{0.61\columnwidth}\raggedright
Diversify supplier locations; avoid sole reliance on one country\strut
\end{minipage}\tabularnewline
\begin{minipage}[t]{0.33\columnwidth}\raggedright
Regulatory action (Import Alert, Warning Letter)\strut
\end{minipage} & \begin{minipage}[t]{0.61\columnwidth}\raggedright
Monitoring of FDA enforcement actions against suppliers\strut
\end{minipage}\tabularnewline
\begin{minipage}[t]{0.33\columnwidth}\raggedright
Geopolitical disruption\strut
\end{minipage} & \begin{minipage}[t]{0.61\columnwidth}\raggedright
Safety stock policy (3-6 months for critical materials)\strut
\end{minipage}\tabularnewline
\begin{minipage}[t]{0.33\columnwidth}\raggedright
Quality drift\strut
\end{minipage} & \begin{minipage}[t]{0.61\columnwidth}\raggedright
Trending of incoming material test results; annual quality review\strut
\end{minipage}\tabularnewline
\begin{minipage}[t]{0.33\columnwidth}\raggedright
Raw material discontinuation\strut
\end{minipage} & \begin{minipage}[t]{0.61\columnwidth}\raggedright
Early notification agreements; contingency qualification of
alternatives\strut
\end{minipage}\tabularnewline
\bottomrule
\end{longtable}

\hypertarget{water-systems}{%
\subsection{23.5 Water Systems}\label{water-systems}}

\begin{itemize}
\tightlist
\item
  \textbf{Purified Water (PW)}: USP requirements, for non-sterile
  manufacturing
\item
  \textbf{Water for Injection (WFI)}: By distillation (traditional) or
  membrane-based (now accepted by Ph. Eur. since 2017); required for
  parenteral products
\item
  \textbf{System qualification}: IQ/OQ/PQ, microbial and chemical
  monitoring, action/alert limits
\item
  \textbf{Continuous monitoring}: TOC, conductivity at point-of-use,
  microbial sampling per schedule
\end{itemize}

\begin{center}\rule{0.5\linewidth}{0.5pt}\end{center}

\hypertarget{biosimilar-cmc-pathway}{%
\section{24. Biosimilar CMC Pathway}\label{biosimilar-cmc-pathway}}

\hypertarget{biosimilar-development-is-fundamentally-a-cmc-exercise}{%
\subsection{24.1 Biosimilar Development is Fundamentally a CMC
Exercise}\label{biosimilar-development-is-fundamentally-a-cmc-exercise}}

Unlike innovator biologics where clinical efficacy drives the program,
biosimilar development is anchored in \textbf{analytical similarity}.
The entire regulatory pathway (351(k) in US, Article 10(4) in EU)
depends on demonstrating that the proposed biosimilar is highly similar
to the reference product with no clinically meaningful differences.

\hypertarget{stepwise-development}{%
\subsection{24.2 Stepwise Development}\label{stepwise-development}}

\begin{verbatim}
Analytical Similarity (foundation)
    ↓
Functional Similarity (binding, bioassay)
    ↓
Clinical PK/PD (bridging)
    ↓
Clinical Efficacy/Safety (abbreviated, in one indication)
    ↓
Extrapolation to other indications (based on totality of evidence)
\end{verbatim}

\hypertarget{reference-product-characterization}{%
\subsection{24.3 Reference Product
Characterization}\label{reference-product-characterization}}

\begin{itemize}
\tightlist
\item
  \textbf{Sourcing}: Purchase multiple lots of the US-licensed reference
  product (and EU reference for global filing)
\item
  \textbf{Lot-to-lot variability}: Characterize 10-15+ lots to establish
  the quality range for each CQA
\item
  \textbf{Comprehensive analytical characterization}: Primary structure,
  higher-order structure, glycosylation, charge variants, size variants,
  potency, process-related impurities
\end{itemize}

\hypertarget{analytical-similarity-assessment}{%
\subsection{24.4 Analytical Similarity
Assessment}\label{analytical-similarity-assessment}}

\textbf{Statistical approaches:}

\begin{longtable}[]{@{}lll@{}}
\toprule
\begin{minipage}[b]{0.15\columnwidth}\raggedright
Tier\strut
\end{minipage} & \begin{minipage}[b]{0.28\columnwidth}\raggedright
Attributes\strut
\end{minipage} & \begin{minipage}[b]{0.48\columnwidth}\raggedright
Statistical Method\strut
\end{minipage}\tabularnewline
\midrule
\endhead
\begin{minipage}[t]{0.15\columnwidth}\raggedright
Tier 1 (most critical)\strut
\end{minipage} & \begin{minipage}[t]{0.28\columnwidth}\raggedright
Potency, glycosylation, binding\strut
\end{minipage} & \begin{minipage}[t]{0.48\columnwidth}\raggedright
Equivalence testing (90\% CI within quality range)\strut
\end{minipage}\tabularnewline
\begin{minipage}[t]{0.15\columnwidth}\raggedright
Tier 2 (major)\strut
\end{minipage} & \begin{minipage}[t]{0.28\columnwidth}\raggedright
Charge variants, oxidation, aggregation\strut
\end{minipage} & \begin{minipage}[t]{0.48\columnwidth}\raggedright
Quality range comparison\strut
\end{minipage}\tabularnewline
\begin{minipage}[t]{0.15\columnwidth}\raggedright
Tier 3 (minor)\strut
\end{minipage} & \begin{minipage}[t]{0.28\columnwidth}\raggedright
Visual appearance, pH, osmolality\strut
\end{minipage} & \begin{minipage}[t]{0.48\columnwidth}\raggedright
Raw data comparison\strut
\end{minipage}\tabularnewline
\bottomrule
\end{longtable}

\textbf{FDA Guidance:} ``Development of Therapeutic Protein Biosimilars:
Comparative Analytical Assessment and Other Quality-Related
Considerations'' (2019)

\hypertarget{process-development-for-biosimilars}{%
\subsection{24.5 Process Development for
Biosimilars}\label{process-development-for-biosimilars}}

\begin{itemize}
\tightlist
\item
  \textbf{Reverse engineering}: Develop process to match the target
  product profile derived from reference product characterization (not
  from a known originator process)
\item
  \textbf{Cell line development}: New cell line producing a product with
  matching CQA profile
\item
  \textbf{Process optimization}: Tune upstream/downstream to achieve
  desired glycosylation, charge variant, and purity profiles
\item
  \textbf{Iterative}: Multiple rounds of process modification to
  converge on analytical similarity
\end{itemize}

\hypertarget{interchangeability}{%
\subsection{24.6 Interchangeability}\label{interchangeability}}

\begin{itemize}
\tightlist
\item
  Additional requirement beyond biosimilarity (US only)
\item
  Requires switching studies demonstrating no increased risk from
  alternating between reference and biosimilar
\item
  Designated as interchangeable by FDA → allows pharmacy-level
  substitution without prescriber intervention
\end{itemize}

\begin{center}\rule{0.5\linewidth}{0.5pt}\end{center}

\hypertarget{combination-products-and-fixed-dose-combinations}{%
\section{25. Combination Products and Fixed-Dose
Combinations}\label{combination-products-and-fixed-dose-combinations}}

\hypertarget{drug-device-combination-products}{%
\subsection{25.1 Drug-Device Combination
Products}\label{drug-device-combination-products}}

\textbf{Regulatory framework:} 21 CFR Part 4; jurisdiction determined by
primary mode of action (CDER vs.~CDRH)

\textbf{Common combination products:}

\begin{longtable}[]{@{}lll@{}}
\toprule
\begin{minipage}[b]{0.24\columnwidth}\raggedright
Product Type\strut
\end{minipage} & \begin{minipage}[b]{0.32\columnwidth}\raggedright
Device Component\strut
\end{minipage} & \begin{minipage}[b]{0.35\columnwidth}\raggedright
CMC Considerations\strut
\end{minipage}\tabularnewline
\midrule
\endhead
\begin{minipage}[t]{0.24\columnwidth}\raggedright
Pre-filled syringes with autoinjector\strut
\end{minipage} & \begin{minipage}[t]{0.32\columnwidth}\raggedright
Autoinjector device\strut
\end{minipage} & \begin{minipage}[t]{0.35\columnwidth}\raggedright
Human factors, device reliability, extractables/leachables from
device\strut
\end{minipage}\tabularnewline
\begin{minipage}[t]{0.24\columnwidth}\raggedright
Metered-dose inhalers (MDI)\strut
\end{minipage} & \begin{minipage}[t]{0.32\columnwidth}\raggedright
Valve, actuator\strut
\end{minipage} & \begin{minipage}[t]{0.35\columnwidth}\raggedright
Delivered dose uniformity, aerodynamic particle size distribution
(APSD), propellant compatibility\strut
\end{minipage}\tabularnewline
\begin{minipage}[t]{0.24\columnwidth}\raggedright
Dry powder inhalers (DPI)\strut
\end{minipage} & \begin{minipage}[t]{0.32\columnwidth}\raggedright
Inhaler device\strut
\end{minipage} & \begin{minipage}[t]{0.35\columnwidth}\raggedright
Device-formulation interaction, moisture sensitivity, device
robustness\strut
\end{minipage}\tabularnewline
\begin{minipage}[t]{0.24\columnwidth}\raggedright
Transdermal patches\strut
\end{minipage} & \begin{minipage}[t]{0.32\columnwidth}\raggedright
Adhesive system, backing membrane\strut
\end{minipage} & \begin{minipage}[t]{0.35\columnwidth}\raggedright
Adhesion testing, drug release rate, residual drug, skin
irritation\strut
\end{minipage}\tabularnewline
\begin{minipage}[t]{0.24\columnwidth}\raggedright
Drug-eluting stents\strut
\end{minipage} & \begin{minipage}[t]{0.32\columnwidth}\raggedright
Stent platform\strut
\end{minipage} & \begin{minipage}[t]{0.35\columnwidth}\raggedright
Drug loading, elution kinetics, coating integrity, sterile
processing\strut
\end{minipage}\tabularnewline
\begin{minipage}[t]{0.24\columnwidth}\raggedright
Intravitreal implants\strut
\end{minipage} & \begin{minipage}[t]{0.32\columnwidth}\raggedright
Implant device\strut
\end{minipage} & \begin{minipage}[t]{0.35\columnwidth}\raggedright
Biodegradable polymer, drug release over months, sterility\strut
\end{minipage}\tabularnewline
\bottomrule
\end{longtable}

\textbf{Additional CMC requirements for combination products:} - Design
History File (per 21 CFR 820 for device constituent) - Human
factors/usability engineering (FDA guidance) - Device-specific stability
testing (mechanical, functional) - Biocompatibility testing (ISO 10993
series)

\hypertarget{fixed-dose-combinations-fdcs}{%
\subsection{25.2 Fixed-Dose Combinations
(FDCs)}\label{fixed-dose-combinations-fdcs}}

\textbf{CMC challenges unique to FDCs:} - \textbf{Drug-drug
compatibility}: Binary/ternary excipient compatibility studies between
APIs - \textbf{Dissolution requirements}: Discriminating dissolution
conditions for each API individually - \textbf{Content uniformity}: Must
meet USP \textless905\textgreater{} for each API - \textbf{Stability}:
Monitor degradation of each API and potential cross-degradation (API A
catalyzing degradation of API B) - \textbf{Manufacturing}: May require
separate granulations blended together, bilayer tablets, or coated
pellets in capsule

\textbf{Example (Trikafta):} Triple-combination CFTR modulator --- three
distinct APIs with different physicochemical properties, each requiring
its own impurity control and dissolution specification within a single
fixed-dose tablet.

\hypertarget{orally-inhaled-and-nasal-drug-products-oindp}{%
\subsection{25.3 Orally Inhaled and Nasal Drug Products
(OINDP)}\label{orally-inhaled-and-nasal-drug-products-oindp}}

\begin{itemize}
\tightlist
\item
  \textbf{Unique CQAs}: Delivered dose uniformity (DDU), aerodynamic
  particle size distribution (APSD by cascade impaction), moisture
  content, device actuation force
\item
  \textbf{FDA-specific guidance}: ``Metered Dose Inhaler (MDI) and Dry
  Powder Inhaler (DPI) Drug Products'' --- extensive CMC and
  bioequivalence guidance
\item
  \textbf{Analytical methods}: Andersen Cascade Impactor (ACI), Next
  Generation Impactor (NGI), dose content uniformity apparatus
\item
  \textbf{Device variability}: Must validate across multiple device lots
\end{itemize}

\begin{center}\rule{0.5\linewidth}{0.5pt}\end{center}

\hypertarget{nitrosamine-impurity-assessment}{%
\section{26. Nitrosamine Impurity
Assessment}\label{nitrosamine-impurity-assessment}}

\hypertarget{background}{%
\subsection{26.1 Background}\label{background}}

In 2018, NDMA (N-nitrosodimethylamine) contamination was discovered in
valsartan API from a specific manufacturer, triggering a global crisis.
This expanded to multiple sartan drugs and then to ranitidine,
metformin, and other products.

\hypertarget{regulatory-response}{%
\subsection{26.2 Regulatory Response}\label{regulatory-response}}

\begin{itemize}
\tightlist
\item
  \textbf{FDA}: ``Control of Nitrosamine Impurities in Human Drugs''
  guidance (2021, updated 2023)
\item
  \textbf{EMA}: Article 5(3) referral --- mandatory risk assessments for
  all MAH
\item
  \textbf{Health Canada, PMDA, TGA}: Harmonized requirements
\item
  \textbf{Scope}: All chemically synthesized drug products, biologics,
  and OTC products
\end{itemize}

\hypertarget{risk-assessment-methodology}{%
\subsection{26.3 Risk Assessment
Methodology}\label{risk-assessment-methodology}}

\textbf{Root causes for nitrosamine formation:} 1. Use of nitrite salts
or nitrosating agents in API synthesis 2. Presence of secondary/tertiary
amines + trace nitrites 3. Contaminated raw materials or solvents 4.
Recycled materials/solvents carrying over nitrosating agents 5. Drug
substance degradation forming nitrosamine under acidic conditions 6.
Interaction between drug and excipients or packaging materials
(especially for NDSRIs)

\hypertarget{acceptable-intake-ai-limits}{%
\subsection{26.4 Acceptable Intake (AI)
Limits}\label{acceptable-intake-ai-limits}}

\begin{longtable}[]{@{}ll@{}}
\toprule
Nitrosamine & AI (ng/day)\tabularnewline
\midrule
\endhead
NDMA & 96\tabularnewline
NDEA & 26.5\tabularnewline
NMBA & 96\tabularnewline
NIPEA & 26.5\tabularnewline
NDIPA & 26.5\tabularnewline
NDBA & 26.5\tabularnewline
\bottomrule
\end{longtable}

For NDSRIs (Nitrosamine Drug Substance Related Impurities), AI is
determined from compound-specific carcinogenicity data or read-across.

\hypertarget{analytical-methods}{%
\subsection{26.5 Analytical Methods}\label{analytical-methods}}

\begin{itemize}
\tightlist
\item
  \textbf{LC-MS/MS}: Required sensitivity at ng/level (ppb in drug
  substance, sub-ppm in drug product)
\item
  \textbf{GC-MS/MS}: For volatile nitrosamines
\item
  \textbf{Method challenges}: Matrix effects, low levels, multiple
  analytes in single method
\item
  \textbf{Confirmatory testing}: Required when risk assessment
  identifies potential for nitrosamine formation
\end{itemize}

\hypertarget{regulatory-timelines}{%
\subsection{26.6 Regulatory Timelines}\label{regulatory-timelines}}

\begin{itemize}
\tightlist
\item
  Initial risk assessment: Completed (deadlines passed for most
  products)
\item
  Confirmatory testing: If risk identified, testing must be completed
  and results shared with agencies
\item
  \textbf{Ongoing obligation}: Any manufacturing change must be
  re-evaluated for nitrosamine risk
\end{itemize}

\begin{center}\rule{0.5\linewidth}{0.5pt}\end{center}

\hypertarget{emerging-modalities-adcs-bispecifics-oligonucleotides}{%
\section{27. Emerging Modalities --- ADCs, Bispecifics,
Oligonucleotides}\label{emerging-modalities-adcs-bispecifics-oligonucleotides}}

\hypertarget{antibody-drug-conjugates-adcs}{%
\subsection{27.1 Antibody-Drug Conjugates
(ADCs)}\label{antibody-drug-conjugates-adcs}}

\textbf{Unique CMC challenges:}

\begin{longtable}[]{@{}ll@{}}
\toprule
\begin{minipage}[b]{0.47\columnwidth}\raggedright
Aspect\strut
\end{minipage} & \begin{minipage}[b]{0.47\columnwidth}\raggedright
Detail\strut
\end{minipage}\tabularnewline
\midrule
\endhead
\begin{minipage}[t]{0.47\columnwidth}\raggedright
Conjugation chemistry\strut
\end{minipage} & \begin{minipage}[t]{0.47\columnwidth}\raggedright
Site-specific (engineered cysteine, enzymatic) vs.~stochastic (lysine,
interchain cysteine)\strut
\end{minipage}\tabularnewline
\begin{minipage}[t]{0.47\columnwidth}\raggedright
Drug-Antibody Ratio (DAR)\strut
\end{minipage} & \begin{minipage}[t]{0.47\columnwidth}\raggedright
Critical CQA --- target typically 2-4; measured by HIC, RP-HPLC\strut
\end{minipage}\tabularnewline
\begin{minipage}[t]{0.47\columnwidth}\raggedright
Linker stability\strut
\end{minipage} & \begin{minipage}[t]{0.47\columnwidth}\raggedright
Cleavable (pH-sensitive, protease) vs.~non-cleavable; in vivo stability
is key\strut
\end{minipage}\tabularnewline
\begin{minipage}[t]{0.47\columnwidth}\raggedright
Payload potency\strut
\end{minipage} & \begin{minipage}[t]{0.47\columnwidth}\raggedright
Cytotoxic (auristatins, maytansinoids, camptothecins) --- OEB5
containment required\strut
\end{minipage}\tabularnewline
\begin{minipage}[t]{0.47\columnwidth}\raggedright
Free drug control\strut
\end{minipage} & \begin{minipage}[t]{0.47\columnwidth}\raggedright
Unconjugated payload is a critical impurity --- must be minimized and
controlled\strut
\end{minipage}\tabularnewline
\begin{minipage}[t]{0.47\columnwidth}\raggedright
Aggregate control\strut
\end{minipage} & \begin{minipage}[t]{0.47\columnwidth}\raggedright
Conjugation can increase aggregation propensity\strut
\end{minipage}\tabularnewline
\begin{minipage}[t]{0.47\columnwidth}\raggedright
Manufacturing flow\strut
\end{minipage} & \begin{minipage}[t]{0.47\columnwidth}\raggedright
mAb production → conjugation → purification → fill/finish (multiple CMC
dossiers)\strut
\end{minipage}\tabularnewline
\begin{minipage}[t]{0.47\columnwidth}\raggedright
Analytical methods\strut
\end{minipage} & \begin{minipage}[t]{0.47\columnwidth}\raggedright
HIC (DAR distribution), RP-HPLC (drug loading), SEC (aggregates),
potency (cytotoxicity assay)\strut
\end{minipage}\tabularnewline
\bottomrule
\end{longtable}

\hypertarget{bispecific-antibodies}{%
\subsection{27.2 Bispecific Antibodies}\label{bispecific-antibodies}}

\textbf{Diverse formats with different CMC implications:}

\begin{longtable}[]{@{}ll@{}}
\toprule
\begin{minipage}[b]{0.34\columnwidth}\raggedright
Format\strut
\end{minipage} & \begin{minipage}[b]{0.60\columnwidth}\raggedright
CMC Challenge\strut
\end{minipage}\tabularnewline
\midrule
\endhead
\begin{minipage}[t]{0.34\columnwidth}\raggedright
BiTE (bispecific T-cell engager)\strut
\end{minipage} & \begin{minipage}[t]{0.60\columnwidth}\raggedright
Short half-life → continuous IV infusion or half-life extension\strut
\end{minipage}\tabularnewline
\begin{minipage}[t]{0.34\columnwidth}\raggedright
Knobs-into-holes / CrossMAb\strut
\end{minipage} & \begin{minipage}[t]{0.60\columnwidth}\raggedright
Correct heterodimerization vs.~homodimer; analytical
discrimination\strut
\end{minipage}\tabularnewline
\begin{minipage}[t]{0.34\columnwidth}\raggedright
DVD-Ig (dual variable domain)\strut
\end{minipage} & \begin{minipage}[t]{0.60\columnwidth}\raggedright
Larger molecule, aggregation risk\strut
\end{minipage}\tabularnewline
\begin{minipage}[t]{0.34\columnwidth}\raggedright
Trispecific and beyond\strut
\end{minipage} & \begin{minipage}[t]{0.60\columnwidth}\raggedright
Increasing complexity, more mispaired species\strut
\end{minipage}\tabularnewline
\bottomrule
\end{longtable}

\textbf{Key CQAs unique to bispecifics:} - Correct chain pairing (\%
heterodimer vs.~homodimer/mispaired) - Dual-target binding activity (two
potency assays may be needed) - Charge and size variant analysis
distinguishing heterodimer from homodimer

\hypertarget{oligonucleotide-therapeutics-sirna-aso}{%
\subsection{27.3 Oligonucleotide Therapeutics (siRNA,
ASO)}\label{oligonucleotide-therapeutics-sirna-aso}}

\textbf{Drug substance CMC:}

\begin{longtable}[]{@{}ll@{}}
\toprule
\begin{minipage}[b]{0.47\columnwidth}\raggedright
Aspect\strut
\end{minipage} & \begin{minipage}[b]{0.47\columnwidth}\raggedright
Detail\strut
\end{minipage}\tabularnewline
\midrule
\endhead
\begin{minipage}[t]{0.47\columnwidth}\raggedright
Synthesis\strut
\end{minipage} & \begin{minipage}[t]{0.47\columnwidth}\raggedright
Solid-phase synthesis (phosphoramidite chemistry); 20-25mer
typical\strut
\end{minipage}\tabularnewline
\begin{minipage}[t]{0.47\columnwidth}\raggedright
Purification\strut
\end{minipage} & \begin{minipage}[t]{0.47\columnwidth}\raggedright
IEX chromatography, RP-HPLC\strut
\end{minipage}\tabularnewline
\begin{minipage}[t]{0.47\columnwidth}\raggedright
Key CQAs\strut
\end{minipage} & \begin{minipage}[t]{0.47\columnwidth}\raggedright
Sequence identity, full-length purity (n-1 species are key impurities),
phosphorothioate content, 2'-modification (2'-OMe, 2'-F, 2'-MOE),
counterion (Na⁺)\strut
\end{minipage}\tabularnewline
\begin{minipage}[t]{0.47\columnwidth}\raggedright
GalNAc conjugation\strut
\end{minipage} & \begin{minipage}[t]{0.47\columnwidth}\raggedright
For hepatocyte targeting; conjugation efficiency, linker integrity\strut
\end{minipage}\tabularnewline
\begin{minipage}[t]{0.47\columnwidth}\raggedright
Impurities\strut
\end{minipage} & \begin{minipage}[t]{0.47\columnwidth}\raggedright
n-1/n+1 truncations, depurination products, phosphodiester linkages (in
PS backbone), stereoisomers\strut
\end{minipage}\tabularnewline
\begin{minipage}[t]{0.47\columnwidth}\raggedright
Analytical\strut
\end{minipage} & \begin{minipage}[t]{0.47\columnwidth}\raggedright
LC-MS (intact mass, impurity ID), IEX-HPLC (purity), CGE (length), CD
(secondary structure)\strut
\end{minipage}\tabularnewline
\bottomrule
\end{longtable}

\textbf{Drug product CMC:} - siRNA in lipid nanoparticles: Same LNP
challenges as mRNA (particle size, encapsulation, stability) - ASO
subcutaneous injection: Solution formulation, typically simple phosphate
buffer

\hypertarget{peptide-therapeutics}{%
\subsection{27.4 Peptide Therapeutics}\label{peptide-therapeutics}}

\begin{itemize}
\tightlist
\item
  \textbf{Manufacturing}: Solid-phase peptide synthesis (SPPS) or
  recombinant expression
\item
  \textbf{Key CQAs}: Sequence identity, disulfide bond formation
  (correct pairing), diastereomers (D-amino acid impurities),
  acetylation/truncation impurities, aggregation
\item
  \textbf{Analytical}: RP-HPLC, LC-MS/MS (peptide mapping), CE, amino
  acid analysis
\item
  \textbf{Stability}: Peptides can be prone to deamidation, oxidation,
  and aggregation
\end{itemize}

\hypertarget{radiopharmaceuticals}{%
\subsection{27.5 Radiopharmaceuticals}\label{radiopharmaceuticals}}

\begin{itemize}
\tightlist
\item
  \textbf{Unique challenge}: Short half-lives (⁶⁸Ga: 68 min, ¹⁸F: 110
  min) require rapid manufacturing and QC
\item
  \textbf{Regulatory}: 21 CFR Part 212 (cGMP for PET drugs)
\item
  \textbf{Release testing}: Must be completed before radioactive decay
  renders product unusable → emphasis on rapid methods and parametric
  release
\item
  \textbf{Stability}: Measured in hours, not months
\end{itemize}

\begin{center}\rule{0.5\linewidth}{0.5pt}\end{center}

\hypertarget{drug-master-files-dmfs}{%
\section{28. Drug Master Files (DMFs)}\label{drug-master-files-dmfs}}

\hypertarget{what-is-a-dmf}{%
\subsection{28.1 What is a DMF?}\label{what-is-a-dmf}}

A Drug Master File is a submission to FDA that may be used to provide
confidential detailed information about facilities, processes, or
articles used in the manufacturing, processing, packaging, and storing
of drugs.

\hypertarget{dmf-types}{%
\subsection{28.2 DMF Types}\label{dmf-types}}

\begin{longtable}[]{@{}lll@{}}
\toprule
\begin{minipage}[b]{0.21\columnwidth}\raggedright
Type\strut
\end{minipage} & \begin{minipage}[b]{0.32\columnwidth}\raggedright
Content\strut
\end{minipage} & \begin{minipage}[b]{0.39\columnwidth}\raggedright
Common Use\strut
\end{minipage}\tabularnewline
\midrule
\endhead
\begin{minipage}[t]{0.21\columnwidth}\raggedright
\textbf{Type I}\strut
\end{minipage} & \begin{minipage}[t]{0.32\columnwidth}\raggedright
Manufacturing site, facilities, operating procedures, personnel\strut
\end{minipage} & \begin{minipage}[t]{0.39\columnwidth}\raggedright
Rarely used; largely obsolete\strut
\end{minipage}\tabularnewline
\begin{minipage}[t]{0.21\columnwidth}\raggedright
\textbf{Type II}\strut
\end{minipage} & \begin{minipage}[t]{0.32\columnwidth}\raggedright
Drug substance, drug substance intermediate, drug product\strut
\end{minipage} & \begin{minipage}[t]{0.39\columnwidth}\raggedright
\textbf{Most common} --- used by API manufacturers\strut
\end{minipage}\tabularnewline
\begin{minipage}[t]{0.21\columnwidth}\raggedright
\textbf{Type III}\strut
\end{minipage} & \begin{minipage}[t]{0.32\columnwidth}\raggedright
Packaging material\strut
\end{minipage} & \begin{minipage}[t]{0.39\columnwidth}\raggedright
Container closure systems, particularly for complex packaging\strut
\end{minipage}\tabularnewline
\begin{minipage}[t]{0.21\columnwidth}\raggedright
\textbf{Type IV}\strut
\end{minipage} & \begin{minipage}[t]{0.32\columnwidth}\raggedright
Excipient, colorant, flavor, essence, or material used in
preparation\strut
\end{minipage} & \begin{minipage}[t]{0.39\columnwidth}\raggedright
Novel excipients\strut
\end{minipage}\tabularnewline
\begin{minipage}[t]{0.21\columnwidth}\raggedright
\textbf{Type V}\strut
\end{minipage} & \begin{minipage}[t]{0.32\columnwidth}\raggedright
FDA-accepted reference information\strut
\end{minipage} & \begin{minipage}[t]{0.39\columnwidth}\raggedright
Rarely used\strut
\end{minipage}\tabularnewline
\bottomrule
\end{longtable}

\hypertarget{how-dmfs-work}{%
\subsection{28.3 How DMFs Work}\label{how-dmfs-work}}

\begin{itemize}
\tightlist
\item
  \textbf{DMF holder} (e.g., API manufacturer) files the DMF with FDA
\item
  DMF holder provides a \textbf{Letter of Authorization} to the
  NDA/ANDA/IND applicant
\item
  Applicant references the DMF in their application
\item
  \textbf{FDA reviews the DMF only when referenced} --- DMFs are not
  independently reviewed or approved
\item
  DMF holder maintains the DMF with annual updates
\end{itemize}

\hypertarget{dmf-vs.-asmf-vs.-cep}{%
\subsection{28.4 DMF vs.~ASMF vs.~CEP}\label{dmf-vs.-asmf-vs.-cep}}

\begin{longtable}[]{@{}llll@{}}
\toprule
\begin{minipage}[b]{0.17\columnwidth}\raggedright
Aspect\strut
\end{minipage} & \begin{minipage}[b]{0.24\columnwidth}\raggedright
DMF (FDA)\strut
\end{minipage} & \begin{minipage}[b]{0.24\columnwidth}\raggedright
ASMF (EMA)\strut
\end{minipage} & \begin{minipage}[b]{0.24\columnwidth}\raggedright
CEP (EDQM)\strut
\end{minipage}\tabularnewline
\midrule
\endhead
\begin{minipage}[t]{0.17\columnwidth}\raggedright
Structure\strut
\end{minipage} & \begin{minipage}[t]{0.24\columnwidth}\raggedright
Single confidential document\strut
\end{minipage} & \begin{minipage}[t]{0.24\columnwidth}\raggedright
Split: Applicant's Part (open) + Restricted Part (confidential)\strut
\end{minipage} & \begin{minipage}[t]{0.24\columnwidth}\raggedright
Certificate referencing Ph. Eur. monograph\strut
\end{minipage}\tabularnewline
\begin{minipage}[t]{0.17\columnwidth}\raggedright
Review\strut
\end{minipage} & \begin{minipage}[t]{0.24\columnwidth}\raggedright
Only when referenced in an application\strut
\end{minipage} & \begin{minipage}[t]{0.24\columnwidth}\raggedright
Assessed with each MAA\strut
\end{minipage} & \begin{minipage}[t]{0.24\columnwidth}\raggedright
Assessed independently by EDQM\strut
\end{minipage}\tabularnewline
\begin{minipage}[t]{0.17\columnwidth}\raggedright
Approval status\strut
\end{minipage} & \begin{minipage}[t]{0.24\columnwidth}\raggedright
No standalone approval\strut
\end{minipage} & \begin{minipage}[t]{0.24\columnwidth}\raggedright
No standalone approval\strut
\end{minipage} & \begin{minipage}[t]{0.24\columnwidth}\raggedright
Independent certificate issued\strut
\end{minipage}\tabularnewline
\begin{minipage}[t]{0.17\columnwidth}\raggedright
Strategic advantage\strut
\end{minipage} & \begin{minipage}[t]{0.24\columnwidth}\raggedright
Flexibility\strut
\end{minipage} & \begin{minipage}[t]{0.24\columnwidth}\raggedright
Transparency to applicant\strut
\end{minipage} & \begin{minipage}[t]{0.24\columnwidth}\raggedright
Faster MAA assessment\strut
\end{minipage}\tabularnewline
\bottomrule
\end{longtable}

\hypertarget{strategic-considerations}{%
\subsection{28.5 Strategic
Considerations}\label{strategic-considerations}}

\begin{itemize}
\tightlist
\item
  DMFs protect API manufacturer's proprietary process information while
  allowing drug product applicants to reference it
\item
  Due diligence on DMF quality is essential when sourcing APIs from DMF
  holders
\item
  A poorly maintained DMF can delay NDA/ANDA approval
\end{itemize}

\begin{center}\rule{0.5\linewidth}{0.5pt}\end{center}

\hypertarget{data-integrity-and-electronic-records}{%
\section{29. Data Integrity and Electronic
Records}\label{data-integrity-and-electronic-records}}

\hypertarget{why-data-integrity-matters-in-cmc}{%
\subsection{29.1 Why Data Integrity Matters in
CMC}\label{why-data-integrity-matters-in-cmc}}

Data integrity is a \textbf{top FDA enforcement priority}. Data
integrity failures are a leading cause of Warning Letters, Import
Alerts, and PAI failures, particularly at international manufacturing
sites.

\hypertarget{alcoa-principles}{%
\subsection{29.2 ALCOA+ Principles}\label{alcoa-principles}}

\begin{longtable}[]{@{}ll@{}}
\toprule
Principle & Meaning\tabularnewline
\midrule
\endhead
\textbf{A}ttributable & Data traceable to the person who generated
it\tabularnewline
\textbf{L}egible & Data readable and permanent\tabularnewline
\textbf{C}ontemporaneous & Recorded at the time of the
activity\tabularnewline
\textbf{O}riginal & First recording of the data (or a verified true
copy)\tabularnewline
\textbf{A}ccurate & Data is correct, truthful, and reflects what
actually happened\tabularnewline
+ \textbf{C}omplete & All data included, no selective
reporting\tabularnewline
+ \textbf{C}onsistent & Data follows logical sequence (timestamps, audit
trails)\tabularnewline
+ \textbf{E}nduring & Recorded on durable media, maintained for
retention period\tabularnewline
+ \textbf{A}vailable & Accessible for review throughout retention
period\tabularnewline
\bottomrule
\end{longtable}

\hypertarget{cfr-part-11-eu-gmp-annex-11}{%
\subsection{29.3 21 CFR Part 11 / EU GMP Annex
11}\label{cfr-part-11-eu-gmp-annex-11}}

\textbf{Electronic records and electronic signatures requirements:} -
System validation (computerized system validation --- CSV) - Audit
trail: Secure, computer-generated, time-stamped record of all changes -
Access controls: Unique user IDs, role-based permissions, periodic
access reviews - Electronic signatures: Legally equivalent to
handwritten signatures when compliant - Backup and disaster recovery -
Audit trail review: Must be part of routine batch record review

\hypertarget{common-data-integrity-findings}{%
\subsection{29.4 Common Data Integrity
Findings}\label{common-data-integrity-findings}}

\begin{itemize}
\tightlist
\item
  Shared login credentials
\item
  Disabled audit trails
\item
  Backdating of records
\item
  Selective reporting (deleting or hiding failing test results)
\item
  Uncontrolled use of standalone systems (e.g., laboratory balances, pH
  meters without audit trails)
\item
  Insufficient audit trail review during batch release
\item
  Lack of CSV for laboratory instruments
\end{itemize}

\begin{center}\rule{0.5\linewidth}{0.5pt}\end{center}

\hypertarget{facility-design-environmental-controls-and-utilities}{%
\section{30. Facility Design, Environmental Controls, and
Utilities}\label{facility-design-environmental-controls-and-utilities}}

\hypertarget{cleanroom-classification}{%
\subsection{30.1 Cleanroom
Classification}\label{cleanroom-classification}}

\begin{longtable}[]{@{}lllll@{}}
\toprule
\begin{minipage}[b]{0.08\columnwidth}\raggedright
Standard\strut
\end{minipage} & \begin{minipage}[b]{0.06\columnwidth}\raggedright
Grade\strut
\end{minipage} & \begin{minipage}[b]{0.29\columnwidth}\raggedright
Max Particles ≥0.5µm/m³ (at rest)\strut
\end{minipage} & \begin{minipage}[b]{0.33\columnwidth}\raggedright
Max Particles ≥0.5µm/m³ (in operation)\strut
\end{minipage} & \begin{minipage}[b]{0.10\columnwidth}\raggedright
Typical Use\strut
\end{minipage}\tabularnewline
\midrule
\endhead
\begin{minipage}[t]{0.08\columnwidth}\raggedright
EU GMP\strut
\end{minipage} & \begin{minipage}[t]{0.06\columnwidth}\raggedright
Grade A\strut
\end{minipage} & \begin{minipage}[t]{0.29\columnwidth}\raggedright
3,520\strut
\end{minipage} & \begin{minipage}[t]{0.33\columnwidth}\raggedright
3,520\strut
\end{minipage} & \begin{minipage}[t]{0.10\columnwidth}\raggedright
Aseptic filling, stoppering\strut
\end{minipage}\tabularnewline
\begin{minipage}[t]{0.08\columnwidth}\raggedright
EU GMP\strut
\end{minipage} & \begin{minipage}[t]{0.06\columnwidth}\raggedright
Grade B\strut
\end{minipage} & \begin{minipage}[t]{0.29\columnwidth}\raggedright
3,520\strut
\end{minipage} & \begin{minipage}[t]{0.33\columnwidth}\raggedright
352,000\strut
\end{minipage} & \begin{minipage}[t]{0.10\columnwidth}\raggedright
Background for Grade A\strut
\end{minipage}\tabularnewline
\begin{minipage}[t]{0.08\columnwidth}\raggedright
EU GMP\strut
\end{minipage} & \begin{minipage}[t]{0.06\columnwidth}\raggedright
Grade C\strut
\end{minipage} & \begin{minipage}[t]{0.29\columnwidth}\raggedright
352,000\strut
\end{minipage} & \begin{minipage}[t]{0.33\columnwidth}\raggedright
3,520,000\strut
\end{minipage} & \begin{minipage}[t]{0.10\columnwidth}\raggedright
Less critical processing\strut
\end{minipage}\tabularnewline
\begin{minipage}[t]{0.08\columnwidth}\raggedright
EU GMP\strut
\end{minipage} & \begin{minipage}[t]{0.06\columnwidth}\raggedright
Grade D\strut
\end{minipage} & \begin{minipage}[t]{0.29\columnwidth}\raggedright
3,520,000\strut
\end{minipage} & \begin{minipage}[t]{0.33\columnwidth}\raggedright
Not defined\strut
\end{minipage} & \begin{minipage}[t]{0.10\columnwidth}\raggedright
Preparation, component washing\strut
\end{minipage}\tabularnewline
\begin{minipage}[t]{0.08\columnwidth}\raggedright
ISO 14644\strut
\end{minipage} & \begin{minipage}[t]{0.06\columnwidth}\raggedright
ISO 5\strut
\end{minipage} & \begin{minipage}[t]{0.29\columnwidth}\raggedright
3,520\strut
\end{minipage} & \begin{minipage}[t]{0.33\columnwidth}\raggedright
---\strut
\end{minipage} & \begin{minipage}[t]{0.10\columnwidth}\raggedright
Equivalent to Grade A/B at rest\strut
\end{minipage}\tabularnewline
\begin{minipage}[t]{0.08\columnwidth}\raggedright
ISO 14644\strut
\end{minipage} & \begin{minipage}[t]{0.06\columnwidth}\raggedright
ISO 7\strut
\end{minipage} & \begin{minipage}[t]{0.29\columnwidth}\raggedright
352,000\strut
\end{minipage} & \begin{minipage}[t]{0.33\columnwidth}\raggedright
---\strut
\end{minipage} & \begin{minipage}[t]{0.10\columnwidth}\raggedright
Equivalent to Grade C at rest\strut
\end{minipage}\tabularnewline
\begin{minipage}[t]{0.08\columnwidth}\raggedright
ISO 14644\strut
\end{minipage} & \begin{minipage}[t]{0.06\columnwidth}\raggedright
ISO 8\strut
\end{minipage} & \begin{minipage}[t]{0.29\columnwidth}\raggedright
3,520,000\strut
\end{minipage} & \begin{minipage}[t]{0.33\columnwidth}\raggedright
---\strut
\end{minipage} & \begin{minipage}[t]{0.10\columnwidth}\raggedright
Equivalent to Grade D at rest\strut
\end{minipage}\tabularnewline
\bottomrule
\end{longtable}

\hypertarget{hvac-systems}{%
\subsection{30.2 HVAC Systems}\label{hvac-systems}}

\begin{itemize}
\tightlist
\item
  \textbf{Air changes per hour}: Grade A/B: \textgreater20 ACPH; Grade
  C: \textgreater20 ACPH; Grade D: \textgreater6 ACPH
\item
  \textbf{Pressure differentials}: 10-15 Pa between grades; positive
  pressure in clean areas; negative pressure in potent compound suites
\item
  \textbf{HEPA filtration}: 99.97\% efficiency for particles ≥0.3µm;
  terminal HEPA filters in Grade A/B
\item
  \textbf{Temperature and humidity}: Typically 20±2°C, 45±10\% RH
  (controlled for product and operator comfort)
\item
  \textbf{Qualification}: DQ → IQ → OQ → PQ; HEPA filter integrity
  testing (DOP/PAO challenge), airflow visualization (smoke studies)
\end{itemize}

\hypertarget{environmental-monitoring-program}{%
\subsection{30.3 Environmental Monitoring
Program}\label{environmental-monitoring-program}}

\begin{longtable}[]{@{}lllll@{}}
\toprule
Parameter & Grade A & Grade B & Grade C & Grade D\tabularnewline
\midrule
\endhead
Viable air (CFU/m³) & \textless1 & 10 & 100 & 200\tabularnewline
Settle plates (CFU/4h) & \textless1 & 5 & 50 & 100\tabularnewline
Contact plates (CFU/plate) & \textless1 & 5 & 25 & 50\tabularnewline
Glove prints (CFU/glove) & \textless1 & 5 & --- & ---\tabularnewline
\bottomrule
\end{longtable}

\begin{itemize}
\tightlist
\item
  Trending of EM data with alert and action limits
\item
  Identification of recovered organisms to species level for Grade A/B
\item
  Investigation required for any action limit excursion
\end{itemize}

\hypertarget{eu-gmp-annex-1-2022-revision-key-changes}{%
\subsection{30.4 EU GMP Annex 1 (2022 Revision) --- Key
Changes}\label{eu-gmp-annex-1-2022-revision-key-changes}}

The 2022 revision of Annex 1 is the most significant GMP update for
sterile products in decades: - \textbf{Contamination Control Strategy
(CCS)}: Required holistic document assessing all contamination risks -
\textbf{Expanded scope}: Now covers all sterile products (not just
terminally sterilized and aseptically processed) - \textbf{Barrier
technology preferred}: RABS (Restricted Access Barrier Systems) and
isolators encouraged over conventional cleanrooms - \textbf{Enhanced EM
requirements}: More detailed trending, identification requirements -
\textbf{APS (Aseptic Process Simulation)}: More prescriptive
requirements for media fills - \textbf{Integrity testing}: Enhanced CCI
requirements for all sterile products

\hypertarget{water-systems-1}{%
\subsection{30.5 Water Systems}\label{water-systems-1}}

\begin{longtable}[]{@{}llll@{}}
\toprule
\begin{minipage}[b]{0.14\columnwidth}\raggedright
Type\strut
\end{minipage} & \begin{minipage}[b]{0.25\columnwidth}\raggedright
Generation\strut
\end{minipage} & \begin{minipage}[b]{0.39\columnwidth}\raggedright
Quality Standard\strut
\end{minipage} & \begin{minipage}[b]{0.11\columnwidth}\raggedright
Use\strut
\end{minipage}\tabularnewline
\midrule
\endhead
\begin{minipage}[t]{0.14\columnwidth}\raggedright
Purified Water (PW)\strut
\end{minipage} & \begin{minipage}[t]{0.25\columnwidth}\raggedright
RO + EDI (or distillation)\strut
\end{minipage} & \begin{minipage}[t]{0.39\columnwidth}\raggedright
USP \textless1231\textgreater, conductivity ≤1.3 µS/cm, TOC ≤500
ppb\strut
\end{minipage} & \begin{minipage}[t]{0.11\columnwidth}\raggedright
Non-sterile manufacturing, equipment cleaning\strut
\end{minipage}\tabularnewline
\begin{minipage}[t]{0.14\columnwidth}\raggedright
Water for Injection (WFI)\strut
\end{minipage} & \begin{minipage}[t]{0.25\columnwidth}\raggedright
Distillation or membrane (Ph. Eur. 2017+)\strut
\end{minipage} & \begin{minipage}[t]{0.39\columnwidth}\raggedright
Same as PW + endotoxin ≤0.25 EU/mL\strut
\end{minipage} & \begin{minipage}[t]{0.11\columnwidth}\raggedright
Parenteral manufacturing, final rinse\strut
\end{minipage}\tabularnewline
\begin{minipage}[t]{0.14\columnwidth}\raggedright
Clean Steam\strut
\end{minipage} & \begin{minipage}[t]{0.25\columnwidth}\raggedright
Clean steam generator\strut
\end{minipage} & \begin{minipage}[t]{0.39\columnwidth}\raggedright
Condensate meets WFI specs\strut
\end{minipage} & \begin{minipage}[t]{0.11\columnwidth}\raggedright
Autoclave sterilization\strut
\end{minipage}\tabularnewline
\bottomrule
\end{longtable}

\hypertarget{compressed-gas-systems}{%
\subsection{30.6 Compressed Gas Systems}\label{compressed-gas-systems}}

\begin{itemize}
\tightlist
\item
  Nitrogen, compressed air, CO₂ used in manufacturing
\item
  Must be filtered (0.2µm for sterile applications)
\item
  Qualification: particle count, microbial testing, dew point, oil
  content
\item
  Monitoring: Periodic requalification per schedule
\end{itemize}

\begin{center}\rule{0.5\linewidth}{0.5pt}\end{center}

\hypertarget{cmc-due-diligence-for-licensing-and-ma}{%
\section{31. CMC Due Diligence for Licensing and
M\&A}\label{cmc-due-diligence-for-licensing-and-ma}}

\hypertarget{when-cmc-due-diligence-is-needed}{%
\subsection{31.1 When CMC Due Diligence is
Needed}\label{when-cmc-due-diligence-is-needed}}

\begin{itemize}
\tightlist
\item
  In-licensing a clinical-stage or commercial-stage asset
\item
  Acquiring a company or product line
\item
  Evaluating a CDMO partnership
\item
  Assessing a co-development or co-marketing deal
\end{itemize}

\hypertarget{cmc-due-diligence-assessment-framework}{%
\subsection{31.2 CMC Due Diligence Assessment
Framework}\label{cmc-due-diligence-assessment-framework}}

\begin{longtable}[]{@{}lll@{}}
\toprule
\begin{minipage}[b]{0.18\columnwidth}\raggedright
Area\strut
\end{minipage} & \begin{minipage}[b]{0.40\columnwidth}\raggedright
Key Questions\strut
\end{minipage} & \begin{minipage}[b]{0.33\columnwidth}\raggedright
Red Flags\strut
\end{minipage}\tabularnewline
\midrule
\endhead
\begin{minipage}[t]{0.18\columnwidth}\raggedright
\textbf{Process maturity}\strut
\end{minipage} & \begin{minipage}[t]{0.40\columnwidth}\raggedright
How well-characterized is the process? Is there a design space? Are CPPs
defined?\strut
\end{minipage} & \begin{minipage}[t]{0.33\columnwidth}\raggedright
No DoE data, poorly defined process parameters\strut
\end{minipage}\tabularnewline
\begin{minipage}[t]{0.18\columnwidth}\raggedright
\textbf{Scalability}\strut
\end{minipage} & \begin{minipage}[t]{0.40\columnwidth}\raggedright
Has the process been run at commercial scale? What scale-up challenges
exist?\strut
\end{minipage} & \begin{minipage}[t]{0.33\columnwidth}\raggedright
Only lab-scale data, untested critical steps at scale\strut
\end{minipage}\tabularnewline
\begin{minipage}[t]{0.18\columnwidth}\raggedright
\textbf{Manufacturing site}\strut
\end{minipage} & \begin{minipage}[t]{0.40\columnwidth}\raggedright
GMP-compliant? Recent inspection history? Capacity available?\strut
\end{minipage} & \begin{minipage}[t]{0.33\columnwidth}\raggedright
Warning Letters, 483s, capacity constraints\strut
\end{minipage}\tabularnewline
\begin{minipage}[t]{0.18\columnwidth}\raggedright
\textbf{Supply chain}\strut
\end{minipage} & \begin{minipage}[t]{0.40\columnwidth}\raggedright
Single-source for critical materials? Starting materials
IP-protected?\strut
\end{minipage} & \begin{minipage}[t]{0.33\columnwidth}\raggedright
Sole-source API, no backup suppliers\strut
\end{minipage}\tabularnewline
\begin{minipage}[t]{0.18\columnwidth}\raggedright
\textbf{Analytical methods}\strut
\end{minipage} & \begin{minipage}[t]{0.40\columnwidth}\raggedright
Validated? Stability-indicating? Transfer-ready?\strut
\end{minipage} & \begin{minipage}[t]{0.33\columnwidth}\raggedright
Unvalidated methods, no forced degradation data\strut
\end{minipage}\tabularnewline
\begin{minipage}[t]{0.18\columnwidth}\raggedright
\textbf{Stability}\strut
\end{minipage} & \begin{minipage}[t]{0.40\columnwidth}\raggedright
Adequate data for claimed shelf life? Any adverse trends?\strut
\end{minipage} & \begin{minipage}[t]{0.33\columnwidth}\raggedright
Short stability history, unexplained degradation\strut
\end{minipage}\tabularnewline
\begin{minipage}[t]{0.18\columnwidth}\raggedright
\textbf{IP around manufacturing}\strut
\end{minipage} & \begin{minipage}[t]{0.40\columnwidth}\raggedright
Process patents? Freedom-to-operate?\strut
\end{minipage} & \begin{minipage}[t]{0.33\columnwidth}\raggedright
Originator process patents blocking alternative routes\strut
\end{minipage}\tabularnewline
\begin{minipage}[t]{0.18\columnwidth}\raggedright
\textbf{Regulatory filing}\strut
\end{minipage} & \begin{minipage}[t]{0.40\columnwidth}\raggedright
Module 3 quality? IND/NDA complete and consistent?\strut
\end{minipage} & \begin{minipage}[t]{0.33\columnwidth}\raggedright
Incomplete Module 3, unresolved FDA questions\strut
\end{minipage}\tabularnewline
\begin{minipage}[t]{0.18\columnwidth}\raggedright
\textbf{Comparability}\strut
\end{minipage} & \begin{minipage}[t]{0.40\columnwidth}\raggedright
Process changes during development? Bridging data available?\strut
\end{minipage} & \begin{minipage}[t]{0.33\columnwidth}\raggedright
Multiple process changes without comparability data\strut
\end{minipage}\tabularnewline
\begin{minipage}[t]{0.18\columnwidth}\raggedright
\textbf{CMC team}\strut
\end{minipage} & \begin{minipage}[t]{0.40\columnwidth}\raggedright
Experienced CMC team? Key person dependencies?\strut
\end{minipage} & \begin{minipage}[t]{0.33\columnwidth}\raggedright
CMC knowledge concentrated in 1-2 individuals\strut
\end{minipage}\tabularnewline
\bottomrule
\end{longtable}

\hypertarget{cmc-risk-scoring}{%
\subsection{31.3 CMC Risk Scoring}\label{cmc-risk-scoring}}

\begin{longtable}[]{@{}lll@{}}
\toprule
\begin{minipage}[b]{0.24\columnwidth}\raggedright
Score\strut
\end{minipage} & \begin{minipage}[b]{0.24\columnwidth}\raggedright
Level\strut
\end{minipage} & \begin{minipage}[b]{0.44\columnwidth}\raggedright
Implication\strut
\end{minipage}\tabularnewline
\midrule
\endhead
\begin{minipage}[t]{0.24\columnwidth}\raggedright
Green\strut
\end{minipage} & \begin{minipage}[t]{0.24\columnwidth}\raggedright
Low risk\strut
\end{minipage} & \begin{minipage}[t]{0.44\columnwidth}\raggedright
Standard development, no major concerns\strut
\end{minipage}\tabularnewline
\begin{minipage}[t]{0.24\columnwidth}\raggedright
Yellow\strut
\end{minipage} & \begin{minipage}[t]{0.24\columnwidth}\raggedright
Moderate risk\strut
\end{minipage} & \begin{minipage}[t]{0.44\columnwidth}\raggedright
Addressable with investment and time; factor into deal terms\strut
\end{minipage}\tabularnewline
\begin{minipage}[t]{0.24\columnwidth}\raggedright
Red\strut
\end{minipage} & \begin{minipage}[t]{0.24\columnwidth}\raggedright
High risk\strut
\end{minipage} & \begin{minipage}[t]{0.44\columnwidth}\raggedright
Potential deal-breaker; may require process redesign, site change, or
significant timeline extension\strut
\end{minipage}\tabularnewline
\bottomrule
\end{longtable}

\hypertarget{cmc-related-deal-terms}{%
\subsection{31.4 CMC-Related Deal Terms}\label{cmc-related-deal-terms}}

\begin{itemize}
\tightlist
\item
  \textbf{Tech transfer timeline and milestones} with acceptance
  criteria
\item
  \textbf{Manufacturing supply agreements} with defined quality
  standards
\item
  \textbf{CMC milestone payments} tied to IND filing, process
  validation, NDA submission
\item
  \textbf{Representations and warranties} on GMP compliance, data
  integrity, and regulatory standing
\item
  \textbf{Indemnification} for CMC-related regulatory actions
\end{itemize}

\begin{center}\rule{0.5\linewidth}{0.5pt}\end{center}

\hypertarget{generic-drug-anda-cmc-requirements}{%
\section{32. Generic Drug (ANDA) CMC
Requirements}\label{generic-drug-anda-cmc-requirements}}

\hypertarget{how-anda-cmc-differs-from-nda-cmc}{%
\subsection{32.1 How ANDA CMC Differs from NDA
CMC}\label{how-anda-cmc-differs-from-nda-cmc}}

ANDAs (Abbreviated New Drug Applications) rely on bioequivalence to the
Reference Listed Drug (RLD) rather than independent safety and efficacy
data. CMC requirements are adapted accordingly.

\hypertarget{key-anda-cmc-requirements}{%
\subsection{32.2 Key ANDA CMC
Requirements}\label{key-anda-cmc-requirements}}

\begin{longtable}[]{@{}ll@{}}
\toprule
\begin{minipage}[b]{0.56\columnwidth}\raggedright
Requirement\strut
\end{minipage} & \begin{minipage}[b]{0.38\columnwidth}\raggedright
Detail\strut
\end{minipage}\tabularnewline
\midrule
\endhead
\begin{minipage}[t]{0.56\columnwidth}\raggedright
\textbf{Q1/Q2 Sameness}\strut
\end{minipage} & \begin{minipage}[t]{0.38\columnwidth}\raggedright
Qualitatively (Q1) and quantitatively (Q2) same inactive ingredients as
RLD (for oral solids --- other dosage forms have different rules)\strut
\end{minipage}\tabularnewline
\begin{minipage}[t]{0.56\columnwidth}\raggedright
\textbf{Dissolution}\strut
\end{minipage} & \begin{minipage}[t]{0.38\columnwidth}\raggedright
Comparative dissolution profiles in multiple media; f₂ similarity factor
≥50 (or biowaiver per BCS)\strut
\end{minipage}\tabularnewline
\begin{minipage}[t]{0.56\columnwidth}\raggedright
\textbf{Bioequivalence}\strut
\end{minipage} & \begin{minipage}[t]{0.38\columnwidth}\raggedright
In vivo PK study demonstrating 90\% CI of Cmax and AUC within 80-125\%
of RLD\strut
\end{minipage}\tabularnewline
\begin{minipage}[t]{0.56\columnwidth}\raggedright
\textbf{BCS Biowaivers}\strut
\end{minipage} & \begin{minipage}[t]{0.38\columnwidth}\raggedright
ICH M9 / FDA guidance: BCS Class I drugs may qualify for biowaiver based
on in vitro dissolution data\strut
\end{minipage}\tabularnewline
\begin{minipage}[t]{0.56\columnwidth}\raggedright
\textbf{Stability}\strut
\end{minipage} & \begin{minipage}[t]{0.38\columnwidth}\raggedright
At least 3 batches, 3 months accelerated + 6 months long-term at time of
filing; shelf life based on available data\strut
\end{minipage}\tabularnewline
\begin{minipage}[t]{0.56\columnwidth}\raggedright
\textbf{Process validation}\strut
\end{minipage} & \begin{minipage}[t]{0.38\columnwidth}\raggedright
Required at time of ANDA approval (on-site verification batch) or
committed before commercial distribution\strut
\end{minipage}\tabularnewline
\begin{minipage}[t]{0.56\columnwidth}\raggedright
\textbf{Specifications}\strut
\end{minipage} & \begin{minipage}[t]{0.38\columnwidth}\raggedright
Generally aligned with RLD product; compendial where available\strut
\end{minipage}\tabularnewline
\bottomrule
\end{longtable}

\hypertarget{paragraph-iv-certifications-and-cmc}{%
\subsection{32.3 Paragraph IV Certifications and
CMC}\label{paragraph-iv-certifications-and-cmc}}

When a generic applicant files a Paragraph IV certification challenging
an Orange Book-listed patent: - \textbf{Polymorph patents}: CMC must
demonstrate the generic product uses a non-infringing polymorph (or that
the patent is invalid/not infringed) - \textbf{Formulation patents}:
Must design around patented excipient combinations or ratios -
\textbf{Process patents}: Must use a non-infringing manufacturing
process - CMC strategy directly impacts patent litigation risk and
market exclusivity timing

\begin{center}\rule{0.5\linewidth}{0.5pt}\end{center}

\hypertarget{master-reference-table}{%
\section{33. Master Reference Table}\label{master-reference-table}}

\hypertarget{ich-quality-guidelines}{%
\subsection{ICH Quality Guidelines}\label{ich-quality-guidelines}}

\begin{longtable}[]{@{}lll@{}}
\toprule
\begin{minipage}[b]{0.30\columnwidth}\raggedright
Guideline\strut
\end{minipage} & \begin{minipage}[b]{0.19\columnwidth}\raggedright
Title\strut
\end{minipage} & \begin{minipage}[b]{0.41\columnwidth}\raggedright
CMC Relevance\strut
\end{minipage}\tabularnewline
\midrule
\endhead
\begin{minipage}[t]{0.30\columnwidth}\raggedright
Q1A-Q1E\strut
\end{minipage} & \begin{minipage}[t]{0.19\columnwidth}\raggedright
Stability Testing\strut
\end{minipage} & \begin{minipage}[t]{0.41\columnwidth}\raggedright
Core stability study design and evaluation\strut
\end{minipage}\tabularnewline
\begin{minipage}[t]{0.30\columnwidth}\raggedright
Q2(R2)\strut
\end{minipage} & \begin{minipage}[t]{0.19\columnwidth}\raggedright
Validation of Analytical Procedures\strut
\end{minipage} & \begin{minipage}[t]{0.41\columnwidth}\raggedright
Method validation requirements\strut
\end{minipage}\tabularnewline
\begin{minipage}[t]{0.30\columnwidth}\raggedright
Q3A(R2)\strut
\end{minipage} & \begin{minipage}[t]{0.19\columnwidth}\raggedright
Impurities in New Drug Substances\strut
\end{minipage} & \begin{minipage}[t]{0.41\columnwidth}\raggedright
DS impurity thresholds\strut
\end{minipage}\tabularnewline
\begin{minipage}[t]{0.30\columnwidth}\raggedright
Q3B(R2)\strut
\end{minipage} & \begin{minipage}[t]{0.19\columnwidth}\raggedright
Impurities in New Drug Products\strut
\end{minipage} & \begin{minipage}[t]{0.41\columnwidth}\raggedright
DP degradant thresholds\strut
\end{minipage}\tabularnewline
\begin{minipage}[t]{0.30\columnwidth}\raggedright
Q3C(R8)\strut
\end{minipage} & \begin{minipage}[t]{0.19\columnwidth}\raggedright
Residual Solvents\strut
\end{minipage} & \begin{minipage}[t]{0.41\columnwidth}\raggedright
Solvent classification and limits\strut
\end{minipage}\tabularnewline
\begin{minipage}[t]{0.30\columnwidth}\raggedright
Q3D(R2)\strut
\end{minipage} & \begin{minipage}[t]{0.19\columnwidth}\raggedright
Elemental Impurities\strut
\end{minipage} & \begin{minipage}[t]{0.41\columnwidth}\raggedright
Metal impurity risk assessment\strut
\end{minipage}\tabularnewline
\begin{minipage}[t]{0.30\columnwidth}\raggedright
Q3E\strut
\end{minipage} & \begin{minipage}[t]{0.19\columnwidth}\raggedright
Extractables and Leachables\strut
\end{minipage} & \begin{minipage}[t]{0.41\columnwidth}\raggedright
E\&L assessment framework\strut
\end{minipage}\tabularnewline
\begin{minipage}[t]{0.30\columnwidth}\raggedright
Q5A(R2)\strut
\end{minipage} & \begin{minipage}[t]{0.19\columnwidth}\raggedright
Viral Safety\strut
\end{minipage} & \begin{minipage}[t]{0.41\columnwidth}\raggedright
Viral clearance requirements for biologics\strut
\end{minipage}\tabularnewline
\begin{minipage}[t]{0.30\columnwidth}\raggedright
Q5B\strut
\end{minipage} & \begin{minipage}[t]{0.19\columnwidth}\raggedright
Expression Construct Analysis\strut
\end{minipage} & \begin{minipage}[t]{0.41\columnwidth}\raggedright
Genetic characterization of cell lines\strut
\end{minipage}\tabularnewline
\begin{minipage}[t]{0.30\columnwidth}\raggedright
Q5C\strut
\end{minipage} & \begin{minipage}[t]{0.19\columnwidth}\raggedright
Stability of Biotech Products\strut
\end{minipage} & \begin{minipage}[t]{0.41\columnwidth}\raggedright
Biologics-specific stability considerations\strut
\end{minipage}\tabularnewline
\begin{minipage}[t]{0.30\columnwidth}\raggedright
Q5D\strut
\end{minipage} & \begin{minipage}[t]{0.19\columnwidth}\raggedright
Cell Substrate Characterisation\strut
\end{minipage} & \begin{minipage}[t]{0.41\columnwidth}\raggedright
Cell bank testing requirements\strut
\end{minipage}\tabularnewline
\begin{minipage}[t]{0.30\columnwidth}\raggedright
Q5E\strut
\end{minipage} & \begin{minipage}[t]{0.19\columnwidth}\raggedright
Comparability\strut
\end{minipage} & \begin{minipage}[t]{0.41\columnwidth}\raggedright
Pre/post-change comparability for biologics\strut
\end{minipage}\tabularnewline
\begin{minipage}[t]{0.30\columnwidth}\raggedright
Q6A\strut
\end{minipage} & \begin{minipage}[t]{0.19\columnwidth}\raggedright
Specifications: Chemical Substances\strut
\end{minipage} & \begin{minipage}[t]{0.41\columnwidth}\raggedright
Test requirements and acceptance criteria\strut
\end{minipage}\tabularnewline
\begin{minipage}[t]{0.30\columnwidth}\raggedright
Q6B\strut
\end{minipage} & \begin{minipage}[t]{0.19\columnwidth}\raggedright
Specifications: Biotech Products\strut
\end{minipage} & \begin{minipage}[t]{0.41\columnwidth}\raggedright
Biologic-specific test requirements\strut
\end{minipage}\tabularnewline
\begin{minipage}[t]{0.30\columnwidth}\raggedright
Q7\strut
\end{minipage} & \begin{minipage}[t]{0.19\columnwidth}\raggedright
GMP for APIs\strut
\end{minipage} & \begin{minipage}[t]{0.41\columnwidth}\raggedright
API manufacturing GMP standard\strut
\end{minipage}\tabularnewline
\begin{minipage}[t]{0.30\columnwidth}\raggedright
Q8(R2)\strut
\end{minipage} & \begin{minipage}[t]{0.19\columnwidth}\raggedright
Pharmaceutical Development\strut
\end{minipage} & \begin{minipage}[t]{0.41\columnwidth}\raggedright
QbD, design space, QTPP, CQAs\strut
\end{minipage}\tabularnewline
\begin{minipage}[t]{0.30\columnwidth}\raggedright
Q9(R1)\strut
\end{minipage} & \begin{minipage}[t]{0.19\columnwidth}\raggedright
Quality Risk Management\strut
\end{minipage} & \begin{minipage}[t]{0.41\columnwidth}\raggedright
Risk assessment tools and process\strut
\end{minipage}\tabularnewline
\begin{minipage}[t]{0.30\columnwidth}\raggedright
Q10\strut
\end{minipage} & \begin{minipage}[t]{0.19\columnwidth}\raggedright
Pharmaceutical Quality System\strut
\end{minipage} & \begin{minipage}[t]{0.41\columnwidth}\raggedright
PQS model, CAPA, change management\strut
\end{minipage}\tabularnewline
\begin{minipage}[t]{0.30\columnwidth}\raggedright
Q11\strut
\end{minipage} & \begin{minipage}[t]{0.19\columnwidth}\raggedright
Drug Substance Development\strut
\end{minipage} & \begin{minipage}[t]{0.41\columnwidth}\raggedright
Extends QbD to API, starting material justification\strut
\end{minipage}\tabularnewline
\begin{minipage}[t]{0.30\columnwidth}\raggedright
Q12\strut
\end{minipage} & \begin{minipage}[t]{0.19\columnwidth}\raggedright
Lifecycle Management\strut
\end{minipage} & \begin{minipage}[t]{0.41\columnwidth}\raggedright
Established conditions, PACMPs\strut
\end{minipage}\tabularnewline
\begin{minipage}[t]{0.30\columnwidth}\raggedright
Q13\strut
\end{minipage} & \begin{minipage}[t]{0.19\columnwidth}\raggedright
Continuous Manufacturing\strut
\end{minipage} & \begin{minipage}[t]{0.41\columnwidth}\raggedright
CM guidance for DS and DP\strut
\end{minipage}\tabularnewline
\begin{minipage}[t]{0.30\columnwidth}\raggedright
Q14\strut
\end{minipage} & \begin{minipage}[t]{0.19\columnwidth}\raggedright
Analytical Procedure Development\strut
\end{minipage} & \begin{minipage}[t]{0.41\columnwidth}\raggedright
Analytical QbD, ATP, PAR\strut
\end{minipage}\tabularnewline
\begin{minipage}[t]{0.30\columnwidth}\raggedright
M4Q(R1)\strut
\end{minipage} & \begin{minipage}[t]{0.19\columnwidth}\raggedright
CTD --- Quality\strut
\end{minipage} & \begin{minipage}[t]{0.41\columnwidth}\raggedright
Module 3 structure and content\strut
\end{minipage}\tabularnewline
\begin{minipage}[t]{0.30\columnwidth}\raggedright
M7(R2)\strut
\end{minipage} & \begin{minipage}[t]{0.19\columnwidth}\raggedright
Mutagenic Impurities\strut
\end{minipage} & \begin{minipage}[t]{0.41\columnwidth}\raggedright
Assessment and control of genotoxic impurities\strut
\end{minipage}\tabularnewline
\bottomrule
\end{longtable}

\hypertarget{key-fda-regulations}{%
\subsection{Key FDA Regulations}\label{key-fda-regulations}}

\begin{longtable}[]{@{}ll@{}}
\toprule
\begin{minipage}[b]{0.57\columnwidth}\raggedright
Regulation\strut
\end{minipage} & \begin{minipage}[b]{0.37\columnwidth}\raggedright
Title\strut
\end{minipage}\tabularnewline
\midrule
\endhead
\begin{minipage}[t]{0.57\columnwidth}\raggedright
21 CFR 210/211\strut
\end{minipage} & \begin{minipage}[t]{0.37\columnwidth}\raggedright
cGMP for Finished Pharmaceuticals\strut
\end{minipage}\tabularnewline
\begin{minipage}[t]{0.57\columnwidth}\raggedright
21 CFR 312\strut
\end{minipage} & \begin{minipage}[t]{0.37\columnwidth}\raggedright
IND Regulations (312.23 = CMC content)\strut
\end{minipage}\tabularnewline
\begin{minipage}[t]{0.57\columnwidth}\raggedright
21 CFR 314\strut
\end{minipage} & \begin{minipage}[t]{0.37\columnwidth}\raggedright
NDA Regulations (314.50 = application content; 314.70 = post-approval
changes)\strut
\end{minipage}\tabularnewline
\begin{minipage}[t]{0.57\columnwidth}\raggedright
21 CFR 601\strut
\end{minipage} & \begin{minipage}[t]{0.37\columnwidth}\raggedright
BLA Regulations\strut
\end{minipage}\tabularnewline
\begin{minipage}[t]{0.57\columnwidth}\raggedright
21 CFR 600-680\strut
\end{minipage} & \begin{minipage}[t]{0.37\columnwidth}\raggedright
Biologics-specific Regulations\strut
\end{minipage}\tabularnewline
\bottomrule
\end{longtable}

\hypertarget{key-fda-guidance-documents}{%
\subsection{Key FDA Guidance
Documents}\label{key-fda-guidance-documents}}

\begin{longtable}[]{@{}lll@{}}
\toprule
\begin{minipage}[b]{0.40\columnwidth}\raggedright
Guidance\strut
\end{minipage} & \begin{minipage}[b]{0.24\columnwidth}\raggedright
Year\strut
\end{minipage} & \begin{minipage}[b]{0.28\columnwidth}\raggedright
Topic\strut
\end{minipage}\tabularnewline
\midrule
\endhead
\begin{minipage}[t]{0.40\columnwidth}\raggedright
CGMP for Phase 1 Investigational Drugs\strut
\end{minipage} & \begin{minipage}[t]{0.24\columnwidth}\raggedright
2008\strut
\end{minipage} & \begin{minipage}[t]{0.28\columnwidth}\raggedright
Reduced GMP for Phase 1\strut
\end{minipage}\tabularnewline
\begin{minipage}[t]{0.40\columnwidth}\raggedright
Process Validation: General Principles and Practices\strut
\end{minipage} & \begin{minipage}[t]{0.24\columnwidth}\raggedright
2011\strut
\end{minipage} & \begin{minipage}[t]{0.28\columnwidth}\raggedright
Three-stage lifecycle validation\strut
\end{minipage}\tabularnewline
\begin{minipage}[t]{0.40\columnwidth}\raggedright
PAT Framework\strut
\end{minipage} & \begin{minipage}[t]{0.24\columnwidth}\raggedright
2004\strut
\end{minipage} & \begin{minipage}[t]{0.28\columnwidth}\raggedright
Real-time monitoring and control\strut
\end{minipage}\tabularnewline
\begin{minipage}[t]{0.40\columnwidth}\raggedright
Pharmaceutical cGMPs for the 21st Century\strut
\end{minipage} & \begin{minipage}[t]{0.24\columnwidth}\raggedright
2004\strut
\end{minipage} & \begin{minipage}[t]{0.28\columnwidth}\raggedright
Risk-based quality framework\strut
\end{minipage}\tabularnewline
\begin{minipage}[t]{0.40\columnwidth}\raggedright
SUPAC-IR\strut
\end{minipage} & \begin{minipage}[t]{0.24\columnwidth}\raggedright
1995\strut
\end{minipage} & \begin{minipage}[t]{0.28\columnwidth}\raggedright
Post-approval changes for IR oral solids\strut
\end{minipage}\tabularnewline
\begin{minipage}[t]{0.40\columnwidth}\raggedright
SUPAC-MR\strut
\end{minipage} & \begin{minipage}[t]{0.24\columnwidth}\raggedright
1997\strut
\end{minipage} & \begin{minipage}[t]{0.28\columnwidth}\raggedright
Post-approval changes for MR oral solids\strut
\end{minipage}\tabularnewline
\begin{minipage}[t]{0.40\columnwidth}\raggedright
SUPAC-SS\strut
\end{minipage} & \begin{minipage}[t]{0.24\columnwidth}\raggedright
1997\strut
\end{minipage} & \begin{minipage}[t]{0.28\columnwidth}\raggedright
Post-approval changes for semisolids\strut
\end{minipage}\tabularnewline
\begin{minipage}[t]{0.40\columnwidth}\raggedright
ANDAs: Stability Testing\strut
\end{minipage} & \begin{minipage}[t]{0.24\columnwidth}\raggedright
Various\strut
\end{minipage} & \begin{minipage}[t]{0.28\columnwidth}\raggedright
Stability testing guidance\strut
\end{minipage}\tabularnewline
\begin{minipage}[t]{0.40\columnwidth}\raggedright
Container Closure Systems\strut
\end{minipage} & \begin{minipage}[t]{0.24\columnwidth}\raggedright
Various\strut
\end{minipage} & \begin{minipage}[t]{0.28\columnwidth}\raggedright
Packaging requirements\strut
\end{minipage}\tabularnewline
\begin{minipage}[t]{0.40\columnwidth}\raggedright
Analytical Procedures and Methods Validation\strut
\end{minipage} & \begin{minipage}[t]{0.24\columnwidth}\raggedright
2015\strut
\end{minipage} & \begin{minipage}[t]{0.28\columnwidth}\raggedright
Complementing ICH Q2\strut
\end{minipage}\tabularnewline
\begin{minipage}[t]{0.40\columnwidth}\raggedright
CMC Information for Human Gene Therapy INDs\strut
\end{minipage} & \begin{minipage}[t]{0.24\columnwidth}\raggedright
2020\strut
\end{minipage} & \begin{minipage}[t]{0.28\columnwidth}\raggedright
Gene therapy CMC requirements\strut
\end{minipage}\tabularnewline
\begin{minipage}[t]{0.40\columnwidth}\raggedright
Control of Nitrosamine Impurities in Human Drugs\strut
\end{minipage} & \begin{minipage}[t]{0.24\columnwidth}\raggedright
2021/2023\strut
\end{minipage} & \begin{minipage}[t]{0.28\columnwidth}\raggedright
Nitrosamine risk assessment and control\strut
\end{minipage}\tabularnewline
\begin{minipage}[t]{0.40\columnwidth}\raggedright
Development of Therapeutic Protein Biosimilars\strut
\end{minipage} & \begin{minipage}[t]{0.24\columnwidth}\raggedright
2019\strut
\end{minipage} & \begin{minipage}[t]{0.28\columnwidth}\raggedright
Comparative analytical assessment for biosimilars\strut
\end{minipage}\tabularnewline
\begin{minipage}[t]{0.40\columnwidth}\raggedright
ANDA Submissions --- Content and Format\strut
\end{minipage} & \begin{minipage}[t]{0.24\columnwidth}\raggedright
Various\strut
\end{minipage} & \begin{minipage}[t]{0.28\columnwidth}\raggedright
Generic drug CMC requirements\strut
\end{minipage}\tabularnewline
\begin{minipage}[t]{0.40\columnwidth}\raggedright
Data Integrity and Compliance with Drug CGMP\strut
\end{minipage} & \begin{minipage}[t]{0.24\columnwidth}\raggedright
2018\strut
\end{minipage} & \begin{minipage}[t]{0.28\columnwidth}\raggedright
ALCOA+ principles, electronic records\strut
\end{minipage}\tabularnewline
\begin{minipage}[t]{0.40\columnwidth}\raggedright
Combination Products --- cGMP Requirements\strut
\end{minipage} & \begin{minipage}[t]{0.24\columnwidth}\raggedright
Various\strut
\end{minipage} & \begin{minipage}[t]{0.28\columnwidth}\raggedright
21 CFR Part 4, device constituent GMP\strut
\end{minipage}\tabularnewline
\begin{minipage}[t]{0.40\columnwidth}\raggedright
Metered Dose Inhaler and DPI Drug Products\strut
\end{minipage} & \begin{minipage}[t]{0.24\columnwidth}\raggedright
1998/2018\strut
\end{minipage} & \begin{minipage}[t]{0.28\columnwidth}\raggedright
OINDP CMC and bioequivalence\strut
\end{minipage}\tabularnewline
\bottomrule
\end{longtable}

\begin{center}\rule{0.5\linewidth}{0.5pt}\end{center}

\hypertarget{glossary-of-key-cmc-terms}{%
\section{34. Glossary of Key CMC
Terms}\label{glossary-of-key-cmc-terms}}

\begin{longtable}[]{@{}ll@{}}
\toprule
\begin{minipage}[b]{0.33\columnwidth}\raggedright
Term\strut
\end{minipage} & \begin{minipage}[b]{0.61\columnwidth}\raggedright
Definition\strut
\end{minipage}\tabularnewline
\midrule
\endhead
\begin{minipage}[t]{0.33\columnwidth}\raggedright
\textbf{API}\strut
\end{minipage} & \begin{minipage}[t]{0.61\columnwidth}\raggedright
Active Pharmaceutical Ingredient --- the drug substance\strut
\end{minipage}\tabularnewline
\begin{minipage}[t]{0.33\columnwidth}\raggedright
\textbf{APQR}\strut
\end{minipage} & \begin{minipage}[t]{0.61\columnwidth}\raggedright
Annual Product Quality Review\strut
\end{minipage}\tabularnewline
\begin{minipage}[t]{0.33\columnwidth}\raggedright
\textbf{ASD}\strut
\end{minipage} & \begin{minipage}[t]{0.61\columnwidth}\raggedright
Amorphous Solid Dispersion\strut
\end{minipage}\tabularnewline
\begin{minipage}[t]{0.33\columnwidth}\raggedright
\textbf{ASMF}\strut
\end{minipage} & \begin{minipage}[t]{0.61\columnwidth}\raggedright
Active Substance Master File (EU equivalent of DMF)\strut
\end{minipage}\tabularnewline
\begin{minipage}[t]{0.33\columnwidth}\raggedright
\textbf{ATP}\strut
\end{minipage} & \begin{minipage}[t]{0.61\columnwidth}\raggedright
Analytical Target Profile\strut
\end{minipage}\tabularnewline
\begin{minipage}[t]{0.33\columnwidth}\raggedright
\textbf{BCS}\strut
\end{minipage} & \begin{minipage}[t]{0.61\columnwidth}\raggedright
Biopharmaceutics Classification System\strut
\end{minipage}\tabularnewline
\begin{minipage}[t]{0.33\columnwidth}\raggedright
\textbf{BLA}\strut
\end{minipage} & \begin{minipage}[t]{0.61\columnwidth}\raggedright
Biologics License Application\strut
\end{minipage}\tabularnewline
\begin{minipage}[t]{0.33\columnwidth}\raggedright
\textbf{CAPA}\strut
\end{minipage} & \begin{minipage}[t]{0.61\columnwidth}\raggedright
Corrective Action and Preventive Action\strut
\end{minipage}\tabularnewline
\begin{minipage}[t]{0.33\columnwidth}\raggedright
\textbf{CBE}\strut
\end{minipage} & \begin{minipage}[t]{0.61\columnwidth}\raggedright
Changes Being Effected (supplement type)\strut
\end{minipage}\tabularnewline
\begin{minipage}[t]{0.33\columnwidth}\raggedright
\textbf{CCI}\strut
\end{minipage} & \begin{minipage}[t]{0.61\columnwidth}\raggedright
Container Closure Integrity\strut
\end{minipage}\tabularnewline
\begin{minipage}[t]{0.33\columnwidth}\raggedright
\textbf{CEP}\strut
\end{minipage} & \begin{minipage}[t]{0.61\columnwidth}\raggedright
Certificate of Suitability (EDQM)\strut
\end{minipage}\tabularnewline
\begin{minipage}[t]{0.33\columnwidth}\raggedright
\textbf{CPP}\strut
\end{minipage} & \begin{minipage}[t]{0.61\columnwidth}\raggedright
Critical Process Parameter\strut
\end{minipage}\tabularnewline
\begin{minipage}[t]{0.33\columnwidth}\raggedright
\textbf{CPV}\strut
\end{minipage} & \begin{minipage}[t]{0.61\columnwidth}\raggedright
Continued Process Verification\strut
\end{minipage}\tabularnewline
\begin{minipage}[t]{0.33\columnwidth}\raggedright
\textbf{CQA}\strut
\end{minipage} & \begin{minipage}[t]{0.61\columnwidth}\raggedright
Critical Quality Attribute\strut
\end{minipage}\tabularnewline
\begin{minipage}[t]{0.33\columnwidth}\raggedright
\textbf{CRL}\strut
\end{minipage} & \begin{minipage}[t]{0.61\columnwidth}\raggedright
Complete Response Letter\strut
\end{minipage}\tabularnewline
\begin{minipage}[t]{0.33\columnwidth}\raggedright
\textbf{CTD}\strut
\end{minipage} & \begin{minipage}[t]{0.61\columnwidth}\raggedright
Common Technical Document\strut
\end{minipage}\tabularnewline
\begin{minipage}[t]{0.33\columnwidth}\raggedright
\textbf{DMF}\strut
\end{minipage} & \begin{minipage}[t]{0.61\columnwidth}\raggedright
Drug Master File\strut
\end{minipage}\tabularnewline
\begin{minipage}[t]{0.33\columnwidth}\raggedright
\textbf{DoE}\strut
\end{minipage} & \begin{minipage}[t]{0.61\columnwidth}\raggedright
Design of Experiments\strut
\end{minipage}\tabularnewline
\begin{minipage}[t]{0.33\columnwidth}\raggedright
\textbf{DP}\strut
\end{minipage} & \begin{minipage}[t]{0.61\columnwidth}\raggedright
Drug Product\strut
\end{minipage}\tabularnewline
\begin{minipage}[t]{0.33\columnwidth}\raggedright
\textbf{DS}\strut
\end{minipage} & \begin{minipage}[t]{0.61\columnwidth}\raggedright
Drug Substance\strut
\end{minipage}\tabularnewline
\begin{minipage}[t]{0.33\columnwidth}\raggedright
\textbf{DSC}\strut
\end{minipage} & \begin{minipage}[t]{0.61\columnwidth}\raggedright
Differential Scanning Calorimetry\strut
\end{minipage}\tabularnewline
\begin{minipage}[t]{0.33\columnwidth}\raggedright
\textbf{E\&L}\strut
\end{minipage} & \begin{minipage}[t]{0.61\columnwidth}\raggedright
Extractables and Leachables\strut
\end{minipage}\tabularnewline
\begin{minipage}[t]{0.33\columnwidth}\raggedright
\textbf{EC}\strut
\end{minipage} & \begin{minipage}[t]{0.61\columnwidth}\raggedright
Established Conditions (ICH Q12)\strut
\end{minipage}\tabularnewline
\begin{minipage}[t]{0.33\columnwidth}\raggedright
\textbf{EMA}\strut
\end{minipage} & \begin{minipage}[t]{0.61\columnwidth}\raggedright
European Medicines Agency\strut
\end{minipage}\tabularnewline
\begin{minipage}[t]{0.33\columnwidth}\raggedright
\textbf{FMEA}\strut
\end{minipage} & \begin{minipage}[t]{0.61\columnwidth}\raggedright
Failure Mode and Effects Analysis\strut
\end{minipage}\tabularnewline
\begin{minipage}[t]{0.33\columnwidth}\raggedright
\textbf{GMP/cGMP}\strut
\end{minipage} & \begin{minipage}[t]{0.61\columnwidth}\raggedright
(current) Good Manufacturing Practice\strut
\end{minipage}\tabularnewline
\begin{minipage}[t]{0.33\columnwidth}\raggedright
\textbf{HBEL}\strut
\end{minipage} & \begin{minipage}[t]{0.61\columnwidth}\raggedright
Health-Based Exposure Limit\strut
\end{minipage}\tabularnewline
\begin{minipage}[t]{0.33\columnwidth}\raggedright
\textbf{HCP}\strut
\end{minipage} & \begin{minipage}[t]{0.61\columnwidth}\raggedright
Host Cell Protein\strut
\end{minipage}\tabularnewline
\begin{minipage}[t]{0.33\columnwidth}\raggedright
\textbf{HME}\strut
\end{minipage} & \begin{minipage}[t]{0.61\columnwidth}\raggedright
Hot Melt Extrusion\strut
\end{minipage}\tabularnewline
\begin{minipage}[t]{0.33\columnwidth}\raggedright
\textbf{HVLD}\strut
\end{minipage} & \begin{minipage}[t]{0.61\columnwidth}\raggedright
High-Voltage Leak Detection\strut
\end{minipage}\tabularnewline
\begin{minipage}[t]{0.33\columnwidth}\raggedright
\textbf{ICH}\strut
\end{minipage} & \begin{minipage}[t]{0.61\columnwidth}\raggedright
International Council for Harmonisation\strut
\end{minipage}\tabularnewline
\begin{minipage}[t]{0.33\columnwidth}\raggedright
\textbf{IND}\strut
\end{minipage} & \begin{minipage}[t]{0.61\columnwidth}\raggedright
Investigational New Drug Application\strut
\end{minipage}\tabularnewline
\begin{minipage}[t]{0.33\columnwidth}\raggedright
\textbf{IPC}\strut
\end{minipage} & \begin{minipage}[t]{0.61\columnwidth}\raggedright
In-Process Control\strut
\end{minipage}\tabularnewline
\begin{minipage}[t]{0.33\columnwidth}\raggedright
\textbf{IQ/OQ/PQ}\strut
\end{minipage} & \begin{minipage}[t]{0.61\columnwidth}\raggedright
Installation/Operational/Performance Qualification\strut
\end{minipage}\tabularnewline
\begin{minipage}[t]{0.33\columnwidth}\raggedright
\textbf{LNP}\strut
\end{minipage} & \begin{minipage}[t]{0.61\columnwidth}\raggedright
Lipid Nanoparticle\strut
\end{minipage}\tabularnewline
\begin{minipage}[t]{0.33\columnwidth}\raggedright
\textbf{MCB/WCB}\strut
\end{minipage} & \begin{minipage}[t]{0.61\columnwidth}\raggedright
Master Cell Bank / Working Cell Bank\strut
\end{minipage}\tabularnewline
\begin{minipage}[t]{0.33\columnwidth}\raggedright
\textbf{MVTR}\strut
\end{minipage} & \begin{minipage}[t]{0.61\columnwidth}\raggedright
Moisture Vapor Transmission Rate\strut
\end{minipage}\tabularnewline
\begin{minipage}[t]{0.33\columnwidth}\raggedright
\textbf{NDA}\strut
\end{minipage} & \begin{minipage}[t]{0.61\columnwidth}\raggedright
New Drug Application\strut
\end{minipage}\tabularnewline
\begin{minipage}[t]{0.33\columnwidth}\raggedright
\textbf{NOR}\strut
\end{minipage} & \begin{minipage}[t]{0.61\columnwidth}\raggedright
Normal Operating Range\strut
\end{minipage}\tabularnewline
\begin{minipage}[t]{0.33\columnwidth}\raggedright
\textbf{OOS/OOT}\strut
\end{minipage} & \begin{minipage}[t]{0.61\columnwidth}\raggedright
Out of Specification / Out of Trend\strut
\end{minipage}\tabularnewline
\begin{minipage}[t]{0.33\columnwidth}\raggedright
\textbf{PACMP}\strut
\end{minipage} & \begin{minipage}[t]{0.61\columnwidth}\raggedright
Post-Approval Change Management Protocol\strut
\end{minipage}\tabularnewline
\begin{minipage}[t]{0.33\columnwidth}\raggedright
\textbf{PAI}\strut
\end{minipage} & \begin{minipage}[t]{0.61\columnwidth}\raggedright
Pre-Approval Inspection\strut
\end{minipage}\tabularnewline
\begin{minipage}[t]{0.33\columnwidth}\raggedright
\textbf{PAR}\strut
\end{minipage} & \begin{minipage}[t]{0.61\columnwidth}\raggedright
Proven Acceptable Range\strut
\end{minipage}\tabularnewline
\begin{minipage}[t]{0.33\columnwidth}\raggedright
\textbf{PAS}\strut
\end{minipage} & \begin{minipage}[t]{0.61\columnwidth}\raggedright
Prior Approval Supplement\strut
\end{minipage}\tabularnewline
\begin{minipage}[t]{0.33\columnwidth}\raggedright
\textbf{PAT}\strut
\end{minipage} & \begin{minipage}[t]{0.61\columnwidth}\raggedright
Process Analytical Technology\strut
\end{minipage}\tabularnewline
\begin{minipage}[t]{0.33\columnwidth}\raggedright
\textbf{PDE}\strut
\end{minipage} & \begin{minipage}[t]{0.61\columnwidth}\raggedright
Permitted Daily Exposure\strut
\end{minipage}\tabularnewline
\begin{minipage}[t]{0.33\columnwidth}\raggedright
\textbf{PPQ}\strut
\end{minipage} & \begin{minipage}[t]{0.61\columnwidth}\raggedright
Process Performance Qualification\strut
\end{minipage}\tabularnewline
\begin{minipage}[t]{0.33\columnwidth}\raggedright
\textbf{PQS}\strut
\end{minipage} & \begin{minipage}[t]{0.61\columnwidth}\raggedright
Pharmaceutical Quality System\strut
\end{minipage}\tabularnewline
\begin{minipage}[t]{0.33\columnwidth}\raggedright
\textbf{QbD}\strut
\end{minipage} & \begin{minipage}[t]{0.61\columnwidth}\raggedright
Quality by Design\strut
\end{minipage}\tabularnewline
\begin{minipage}[t]{0.33\columnwidth}\raggedright
\textbf{QP}\strut
\end{minipage} & \begin{minipage}[t]{0.61\columnwidth}\raggedright
Qualified Person (EU)\strut
\end{minipage}\tabularnewline
\begin{minipage}[t]{0.33\columnwidth}\raggedright
\textbf{QTPP}\strut
\end{minipage} & \begin{minipage}[t]{0.61\columnwidth}\raggedright
Quality Target Product Profile\strut
\end{minipage}\tabularnewline
\begin{minipage}[t]{0.33\columnwidth}\raggedright
\textbf{RTRT}\strut
\end{minipage} & \begin{minipage}[t]{0.61\columnwidth}\raggedright
Real-Time Release Testing\strut
\end{minipage}\tabularnewline
\begin{minipage}[t]{0.33\columnwidth}\raggedright
\textbf{SCT}\strut
\end{minipage} & \begin{minipage}[t]{0.61\columnwidth}\raggedright
Safety Concern Threshold\strut
\end{minipage}\tabularnewline
\begin{minipage}[t]{0.33\columnwidth}\raggedright
\textbf{SPC}\strut
\end{minipage} & \begin{minipage}[t]{0.61\columnwidth}\raggedright
Statistical Process Control\strut
\end{minipage}\tabularnewline
\begin{minipage}[t]{0.33\columnwidth}\raggedright
\textbf{SUPAC}\strut
\end{minipage} & \begin{minipage}[t]{0.61\columnwidth}\raggedright
Scale-Up and Post-Approval Changes\strut
\end{minipage}\tabularnewline
\begin{minipage}[t]{0.33\columnwidth}\raggedright
\textbf{TTC}\strut
\end{minipage} & \begin{minipage}[t]{0.61\columnwidth}\raggedright
Threshold of Toxicological Concern\strut
\end{minipage}\tabularnewline
\begin{minipage}[t]{0.33\columnwidth}\raggedright
\textbf{XRPD}\strut
\end{minipage} & \begin{minipage}[t]{0.61\columnwidth}\raggedright
X-Ray Powder Diffraction\strut
\end{minipage}\tabularnewline
\begin{minipage}[t]{0.33\columnwidth}\raggedright
\textbf{ADC}\strut
\end{minipage} & \begin{minipage}[t]{0.61\columnwidth}\raggedright
Antibody-Drug Conjugate\strut
\end{minipage}\tabularnewline
\begin{minipage}[t]{0.33\columnwidth}\raggedright
\textbf{ALCOA+}\strut
\end{minipage} & \begin{minipage}[t]{0.61\columnwidth}\raggedright
Attributable, Legible, Contemporaneous, Original, Accurate + Complete,
Consistent, Enduring, Available\strut
\end{minipage}\tabularnewline
\begin{minipage}[t]{0.33\columnwidth}\raggedright
\textbf{ANDA}\strut
\end{minipage} & \begin{minipage}[t]{0.61\columnwidth}\raggedright
Abbreviated New Drug Application\strut
\end{minipage}\tabularnewline
\begin{minipage}[t]{0.33\columnwidth}\raggedright
\textbf{APSD}\strut
\end{minipage} & \begin{minipage}[t]{0.61\columnwidth}\raggedright
Aerodynamic Particle Size Distribution\strut
\end{minipage}\tabularnewline
\begin{minipage}[t]{0.33\columnwidth}\raggedright
\textbf{ASO}\strut
\end{minipage} & \begin{minipage}[t]{0.61\columnwidth}\raggedright
Antisense Oligonucleotide\strut
\end{minipage}\tabularnewline
\begin{minipage}[t]{0.33\columnwidth}\raggedright
\textbf{CDMO}\strut
\end{minipage} & \begin{minipage}[t]{0.61\columnwidth}\raggedright
Contract Development and Manufacturing Organization\strut
\end{minipage}\tabularnewline
\begin{minipage}[t]{0.33\columnwidth}\raggedright
\textbf{CCS}\strut
\end{minipage} & \begin{minipage}[t]{0.61\columnwidth}\raggedright
Contamination Control Strategy (EU GMP Annex 1)\strut
\end{minipage}\tabularnewline
\begin{minipage}[t]{0.33\columnwidth}\raggedright
\textbf{CMO}\strut
\end{minipage} & \begin{minipage}[t]{0.61\columnwidth}\raggedright
Contract Manufacturing Organization\strut
\end{minipage}\tabularnewline
\begin{minipage}[t]{0.33\columnwidth}\raggedright
\textbf{CSV}\strut
\end{minipage} & \begin{minipage}[t]{0.61\columnwidth}\raggedright
Computerized System Validation\strut
\end{minipage}\tabularnewline
\begin{minipage}[t]{0.33\columnwidth}\raggedright
\textbf{DAR}\strut
\end{minipage} & \begin{minipage}[t]{0.61\columnwidth}\raggedright
Drug-Antibody Ratio\strut
\end{minipage}\tabularnewline
\begin{minipage}[t]{0.33\columnwidth}\raggedright
\textbf{DDU}\strut
\end{minipage} & \begin{minipage}[t]{0.61\columnwidth}\raggedright
Delivered Dose Uniformity\strut
\end{minipage}\tabularnewline
\begin{minipage}[t]{0.33\columnwidth}\raggedright
\textbf{FDC}\strut
\end{minipage} & \begin{minipage}[t]{0.61\columnwidth}\raggedright
Fixed-Dose Combination\strut
\end{minipage}\tabularnewline
\begin{minipage}[t]{0.33\columnwidth}\raggedright
\textbf{GDP}\strut
\end{minipage} & \begin{minipage}[t]{0.61\columnwidth}\raggedright
Good Distribution Practice\strut
\end{minipage}\tabularnewline
\begin{minipage}[t]{0.33\columnwidth}\raggedright
\textbf{NDMA}\strut
\end{minipage} & \begin{minipage}[t]{0.61\columnwidth}\raggedright
N-Nitrosodimethylamine\strut
\end{minipage}\tabularnewline
\begin{minipage}[t]{0.33\columnwidth}\raggedright
\textbf{NDSRI}\strut
\end{minipage} & \begin{minipage}[t]{0.61\columnwidth}\raggedright
Nitrosamine Drug Substance Related Impurity\strut
\end{minipage}\tabularnewline
\begin{minipage}[t]{0.33\columnwidth}\raggedright
\textbf{RABS}\strut
\end{minipage} & \begin{minipage}[t]{0.61\columnwidth}\raggedright
Restricted Access Barrier System\strut
\end{minipage}\tabularnewline
\begin{minipage}[t]{0.33\columnwidth}\raggedright
\textbf{RLD}\strut
\end{minipage} & \begin{minipage}[t]{0.61\columnwidth}\raggedright
Reference Listed Drug\strut
\end{minipage}\tabularnewline
\begin{minipage}[t]{0.33\columnwidth}\raggedright
\textbf{siRNA}\strut
\end{minipage} & \begin{minipage}[t]{0.61\columnwidth}\raggedright
Small Interfering RNA\strut
\end{minipage}\tabularnewline
\begin{minipage}[t]{0.33\columnwidth}\raggedright
\textbf{WFI}\strut
\end{minipage} & \begin{minipage}[t]{0.61\columnwidth}\raggedright
Water for Injection\strut
\end{minipage}\tabularnewline
\bottomrule
\end{longtable}

\begin{center}\rule{0.5\linewidth}{0.5pt}\end{center}

\hypertarget{document-information}{%
\section{Document Information}\label{document-information}}

\begin{itemize}
\tightlist
\item
  \textbf{Compiled:} March 2026
\item
  \textbf{Sources:} ICH Guidelines (Q1-Q14, M4Q, M7), FDA 21 CFR
  Regulations, FDA Guidance Documents, EMA Guidelines, USP/Ph. Eur.
  Standards, Pharmaceutical Industry Best Practices
\item
  \textbf{Scope:} Small molecules, biologics, cell therapies, gene
  therapies, mRNA therapeutics
\item
  \textbf{Intended audience:} CMC professionals, regulatory affairs
  strategists, pharmaceutical development teams, consulting engagements
\end{itemize}

\begin{center}\rule{0.5\linewidth}{0.5pt}\end{center}

\emph{This document provides a comprehensive reference for CMC in drug
development. For the most current versions of regulatory guidance
documents, consult \href{https://www.fda.gov}{FDA.gov},
\href{https://www.ich.org}{ICH.org}, and
\href{https://www.ema.europa.eu}{EMA.europa.eu}.}

\end{document}
